% LaTeX source generated from paper3_material_ledgers_v32.md by wave13/build_latex_v13.py (batch 7 (audits of agent arena 1 paper rewrites)/wave13).
% LaTeX source generated by revision/v7/revisions_v33_tex.py from wave13's v32 source plus
% the 30 logged operations of revision/v6/revisions_v33.py (markdown mirror:
% revision/v6/paper3_material_ledgers_v33.md).  Content identical to that mirror; the v32
% source in the clone is untouched.  Compiles error-free with tectonic 0.17.0.
% Pandoc body identical to the wave-9/11 output (asserted at build time), with the
% declaration subsections relocated to a trailing Declarations section.
% Front matter: Amin Abaee, Independent Researcher, clickable ORCID
%(https://orcid.org/0000-0002-0019-1842) and email (amin_abaee@ut.ac.ir)
% beneath the name; date September 6, 2026.
% Back matter: the paper's own declarations, each a titled subsection, plus
% the AI declaration (GLM (Z.ai), Qwen (Alibaba Cloud) and DeepSeek AI
% assisted with drafting and iterative review).
% References carry the owner's Zenodo DOIs (see wave13/apply_md_doi.py).
% Edit the markdown source and re-run the build script to regenerate.
% Compiles error-free with tectonic (also compatible with pdflatex/xelatex).
\documentclass[11pt]{article}
\usepackage[a4paper,margin=1in]{geometry}
\usepackage{amsmath,amssymb}
\usepackage{graphicx}
\usepackage{booktabs,longtable,array,calc}
\usepackage{xcolor}
\usepackage[colorlinks=true,allcolors=blue!45!black]{hyperref}
\usepackage[font=small,labelfont=bf]{caption}
\usepackage{etoolbox}
\AtBeginEnvironment{longtable}{\small}  % dense data tables
\setlength{\tabcolsep}{4pt}
\providecommand{\tightlist}{\setlength{\itemsep}{0pt}\setlength{\parskip}{0pt}}
\setcounter{secnumdepth}{-1}
\emergencystretch=3em
\graphicspath{{../}}
\begin{document}
\title{Typed Flux Ledgers and Depletion Arithmetic: Conservation, Componentwise Diagnostics, and the Semantics of Depletion Horizons}
\author{Amin Abaee\\[0.35em]
{\small Independent Researcher}\\[0.55em]
{\small\href{https://orcid.org/0000-0002-0019-1842}{ORCID: 0000-0002-0019-1842}}\\[0.3em]
{\small\href{mailto:amin\_abaee@ut.ac.ir}{amin\_abaee@ut.ac.ir}}}
\date{September 6, 2026}
\maketitle

\begin{abstract}


Depletion numbers circulate under one label while answering different
questions. Reserve-life ratios, trend-persistence indices, and
removals-only pressure scales are read as ``time to depletion'' despite
measuring different things; compensatory aggregation hides a deficit
behind a positive scalar.

We separate them with a typed stock--flow accounting layer --- a
per-moiety ledger that keeps conservation laws typed, so biomass, money,
and biodiversity are not summed into one scalar. Conservation follows
from the incidence structure of the compartment--flux network;
positivity from donor limitation (each primitive outflow vanishes when
its donor is empty); services are readouts, not conserved mass.

Three certification layers carry a flux-reconstruction identity (an
algebraic relation reconstructing unobserved internal fluxes from
observed stock changes), a conservation-law reduction, and a
flux-bounding envelope theorem. The closed finite-donor ledger satisfies
the natural-block mass identity, orthant invariance (state stays
nonnegative), no interior rest at positive effort, the
vanishing-extraction rest set, and extraction integrability. Depletion
time separates into three non-interchangeable quantities --- gross
turnover intensity, a frozen-rate ratio, and a scenario-conditioned
hitting time --- with uniform-drift bounds. Three public-data
applications are classified: G3P anomaly index is a statistical index,
not a stock ratio; phosphate reserve-life ratio is arithmetic, not a
forecast; fisheries removals-only time is a pressure scale, not a
depletion diagnostic. No nonnegative weighting certifies componentwise
nonnegativity, and five double-counting rules block phantom mass.
First-passage semantics on declared stochastic surrogates --- stochastic
passage times against a relative barrier --- carry explicit non-claims.
Weak and strong sustainability are two regimes of one system,
distinguished by whether the material cycle closes at the rate of use.

Each depletion claim carries the predicate it actually establishes, with
an interface contract fixing the shared object with delay-based
institutional dynamics.

\end{abstract}

\textbf{Keywords:} material flow accounting; stock--flow ledger;
depletion indicators; first-passage time; conservation laws; composite
indicators; reserve life

\begin{center}\rule{0.5\linewidth}{0.5pt}\end{center}

\subsection{1. Introduction}\label{introduction}

\subsubsection{1.1 Failure modes and the two
confusions}\label{failure-modes-and-the-two-confusions}

Sustainability accounting fails in two characteristic ways. The first is
\emph{compensatory aggregation}: heterogeneous physical stocks and
service flows are summarized by scalar indices whose cross-component
trades are never declared as mathematics, so a severe deficit in one
component can coexist with a positive aggregate. The composite-indicator
and weak-versus-strong-sustainability literatures document the failure
and its noncompensatory remedies (Munda and Nardo, 2009; Ekins et al.,
2003; Neumayer, 2013). The second is \emph{classification drift}:
quantities with the units of time --- reserve-life ratios,
trend-persistence indices, removals-only pressure scales --- circulate
as if they were one thing, ``time to depletion,'' when they answer
different questions under different assumptions.

The drift is concrete. A reserve-life ratio divides a reserve figure by
a production figure and calls the quotient a horizon. A groundwater
anomaly index fits a trend to a satellite product and reports the fitted
distance to the series' own minimum divided by the fitted rate. The
result is a number in units of years that is not a time to any physical
event. A fisheries pressure indicator divides a log biomass margin by a
fishing mortality and presents the result as a time scale. Each of these
quantities is informative about something. None is what it is typically
taken to be.

The first failure mode has a public flagship object. Earth Overshoot Day
aggregates component flows into a single calendar date, so a severe
deficit in one component can coexist with a date that still falls late
in the year (Lin et al., 2018; Wackernagel and Beyers, 2019; Blomqvist
et al., 2013). The systems-dynamics overshoot models raise the same
aggregation question in dynamic form (Meadows et al., 1972). The
replacements built in this article are component-resolved by
construction: the per-asset, per-pool depletion horizons of Section
6.5.2 are each reported beside the pool it draws on and never summed
into a scalar date.

A further failure mode is what we call the \emph{productivity illusion}:
the appearance that a system is delivering adequately while the base
that sustains the delivery is being reduced. It has two senses, and they
are distinct. The first is arithmetic and is the
compensatory-aggregation failure above --- a deficit in one component
offset by a surplus in another, so that a positive aggregate reads as
adequate. This article formalises this sense as the aggregation
obstruction of Section 10.1.

The second sense is dynamical and is yield inflation. A measured yield
can exceed the true sustainable yield when it is maintained by drawing
down the support pool --- groundwater, soil carbon, bioavailable
nutrients --- rather than by that pool's regeneration. The resource
appears productive while what sustains it is liquidated silently. The
support pool need not be a superficially non-renewable reserve for this
diagnosis to run, and it can be read in more than one way at once --- as
natural capital, as a stock, or as a slowly regenerating flow of
services. These are overlapping readings of the same pool, not mutually
exclusive ones. Every such pool is regenerative on some timescale --- a
crop within a season, an aquifer within years to decades, a mineral
deposit over geological time --- and the failure is the same in each
case. What decides is the renewal rate of the pool relative to the rate
of use. A use that falls on the pool's regeneration is sustained. One
that falls on the pool itself is liquidation. A drawdown is recoverable
only insofar as the pool regenerates faster than it is taken; if a pool
regenerates too slowly for the rate at which it is taken, drawdown is
still liquidation. The size of the pool is a property of the resource.
Whether a drawdown is recoverable is a property of the rate.

The two-pool architecture of this article exists to detect exactly this
yield-inflation sense. The depletion horizons of Section 6.5.2 are
component-resolved: each is reported on the object it is defined on and
at its own classification status, never summed into a scalar, so no
component deficit is masked by a positive aggregate. Where the
yield-inflation reading calls for an explicit two-pool model --- a
resource beside a slow-regenerating support pool --- the article does
not overreach: no such model is claimed as established. The groundwater
two-pool model is a registered open gap whose admitted applied object is
the one-pool affine approximation (Section 8.2), and spawning biomass is
explicitly not an abiotic support pool (Section 6.5.4). The temporal
asymmetry that makes the yield-inflation sense dangerous has a
mechanical picture. An elevator rated for ten people will hold fourteen
for a long time: the cable does not part on the fourteenth passenger but
only after the accumulated wear it never shows. During the overshoot
nothing announces the damage --- the wear is diffuse, gradual, and
invisible to everyone riding the car --- while the snap, when it comes,
is sudden and total. The yield-inflation sense is the wear phase:
measured yield stays up while nothing in the yield series records the
drawdown. The two-pool architecture exists for exactly this reason: to
make the wear measurable while the cable still holds, not to predict the
snap (Section 7 states what the passage-time results do and do not
condition on). What Section 6.5.2 actually reports is per-resource,
per-pool component-resolved horizons at each object's own status --- not
a supporting-pool column that the applied data do not carry.

The classification-drift failure is well documented on the data side.
Reserve-life ratios are arithmetic: remaining reserves divided by
current production. Because reserves are an economic classification that
changes with prices, technology, exploration, and regulation --- not a
fixed physical stock --- the ratio is not an exhaustion forecast. The
critique has been made forcefully for phosphate. Illakwahhi, Vegi, and
Srivastava (2024) show that the influential ``depletion within a
century'' estimates rest on single-source U.S. Geological Survey (USGS)
data whose credibility is questionable, ignore recovery and recycling,
and that timeframes estimated from inadequate reserve data ``are
somewhat misleading.'' The same point is standard in mineral economics,
where the distinction between reserves and resources has been central at
least since Tilton (2003). The United States' phosphate reserves have
remained near \(1{,}000{,}000\) kt for decades while cumulative
production since 1996 is of order \(600{,}000\) kt. The classification
replenishes itself, and a ratio built on it inherits that behaviour. For
copper, Tilton and Lagos (2007) document reserves growing through a
century of rising production --- the same replenishment under the same
classification.

The same conflation pervades accounting itself. Material flow analysis
supplies the bookkeeping of society's material throughput (Brunner and
Rechberger, 2004; Eurostat, 2001; Fischer-Kowalski et al., 2011). Its
incidence structure is shared with reaction-network theory, where the
sign pattern of the stoichiometric matrix is a conservation object
(Feinberg, 2019). But bookkeeping balance, stoichiometric conservation,
thermodynamic admissibility, and sustainability safety are four
different predicates, and the literature routinely slides between them.
A mass-balanced ledger can be chemically impossible. A chemically
consistent ledger can violate every declared barrier. A ledger
satisfying all declared barriers can fail conservation. Separating these
layers and proving their relationships, so that each claim about a
material system carries the predicate it actually establishes, is this
article's first task (Section 3).

The second task is noncompensation. The weak-comparability thesis of
ecological economics holds that values relevant to environmental
decisions may not be commensurable in a single metric (Martinez-Alier,
Munda, and O'Neill, 1998). At the level of material accounting this
thesis has a precise algebraic form, which Section 10.1 states: no
nonnegative weighting of component balances certifies that every
component satisfies its floor. A positive weighted sum never certifies
componentwise nonnegativity. Scalar summaries may rank and communicate;
certification requires the vector.

The third task is to locate substitution within the ledger, rather than
alongside it. Weak and strong sustainability are not competing
hypotheses but two regimes of one dynamic system, distinguished not by
whether substitution alone keeps pace with depletion but by whether the
material cycle can be closed at the rate of use --- a reading of the two
regimes developed for the ledger, not the received distinction of the
literature, which turns on the substitutability of natural capital
(Neumayer, 2013; Ekins et al., 2003). Weak sustainability is the
idealized regime: humans consume and populate slowly enough that
technological substitution and natural regeneration together
redistribute matter so that it is used as it arises. On human-relevant
timescales substitution is the dominant term; natural regeneration acts
far more slowly --- often on deep-time scales --- and is included for
physical completeness rather than as a co-equal mechanism (Daly, 1990).
The deep-time clause is scoped to the slow compartments --- geological
donors and mineral stocks, whose renewal runs on deep time. The
regenerative compartments this article's applied record tabulates renew
on human timescales --- a crop within a season, an aquifer within years
to decades, a fish stock within years --- so for them regeneration is a
co-equal or dominant term, and the regime question is decided by the
rate of use relative to that renewal, not by any deep-time regeneration.
In this idealized closure the byproducts of use --- carbon drawn from
the atmosphere, chemical substances released to air, water, and soil ---
are returned to use in time and are therefore not waste. Waste is a
relational status, not an intrinsic property of any material. It is the
matter that, under the present relation and circumstances, accumulates
because it cannot be put to immediate use for lack of knowledge,
technology, or timely redistribution. So a chemical substance that is a
product in one context is waste-in-waiting in another, and no substance
is waste by its nature.

Strong sustainability is the regime in which that closure fails ---
either because no identified physical pathway re-routes the extracted
matter, or because depletion outruns the rate at which such a pathway
could be deployed (Neumayer, 2013; Ekins et al., 2003). Substitution is
admissible only where an identified physical pathway exists and does not
draw down a different critical stock. No substitute is admitted merely
because an aggregate production function or a weighted index permits it.
In ledger terms, a substitute is either a recycled or recovered flux
returned to the regenerating pool, or a non-renewable drawdown on a
second compartment. The two are different ledger entries and carry
different statuses (Section 2.3, Section 6).

The two regimes correspond to two readings of the same object. The
scalar reading \(B = b\cdot M\) --- the aggregate regeneration flow,
with \(b\) the per-unit-mass regeneration rate and \(M\) the
natural-block mass; letters local to the introduction, since from
Section 2.2 on \(B\) names the gross turnover \(R + T\) while \(M\)
keeps the natural-block-mass reading --- against consumption is the
weak-sustainability flow check: does the cycle close in aggregate? The
multi-compartment ledger is the vector reading: does it close only by
drawing down a critical compartment or by over-filling the waste
compartment --- that is, are local depletions or waste accumulation
merely being pushed off-book? Both readings are required, and neither
demotes stocks. The ledger keeps the stock compartments first-class and
supplies the thresholds, hitting times, and sink constraints from which
\(B\) is derived. The scalar reading is the operational overshoot test
on top of it. Both readings are instantaneous; a single-period check has
no view of what must also hold over time, namely that the checks do not
drift into overshoot. The balance must be maintained as neither
consumption nor population grows faster than the productivity that
supports them, and the drawdown must not be recoverable only on a
timescale longer than it is taken. The scenario-conditioned hitting time
of Section 6.5.2 exists to give the drift a horizon. The reserve
classification itself encodes this. Reserves change with prices,
technology, exploration, and regulation, so an exhaustion-horizon
estimate built on a reserve figure carries the substitution and
technology premises of the classification, not a physical forecast.

This article builds the accounting layer that resists these failures.
Its objects are typed ledgers in which conservation follows from the
incidence structure, positivity follows from donor limitation, services
are readouts rather than conserved mass, and each depletion quantity
carries its own classification, stated at its actual strength.

\subsubsection{1.2 Contributions}\label{contributions}

\begin{enumerate}
\def\labelenumi{\arabic{enumi}.}
\item
  \textbf{A typed primitive-flux ledger.} Compartments carry a material
  identity, a boundary, and a unit. Non-negative primitive fluxes
  connect them through a signed incidence matrix (the compartment--flux
  bookkeeping of material flow analysis: Brunner and Rechberger, 2004;
  Fischer-Kowalski et al., 2011; the incidence formalism of
  reaction-network theory: Feinberg, 2019). Conversions between types
  appear only as explicit stoichiometric coefficients. Every one-way
  transfer is donor-limited.
\item
  \textbf{Three certification layers, separated and proved.} Accounting
  consistency (the balance law holds), conservation consistency
  (declared moieties are invariant up to boundary flows), and barrier
  safety (declared lower and upper barriers hold) are distinct
  predicates with distinct proof obligations. The flux-reconstruction
  identity, the conservation-law reduction, and a flux-bounding envelope
  theorem --- with a barrier certificate corollary --- are proved in
  full (Section 3).
\item
  \textbf{The closed finite-donor theorem set.} For the closed ledger in
  which the geological donor is a state and no derived target appears:
  the natural-block mass identity, orthant invariance, absence of any
  interior rest point at positive effort, the vanishing-extraction rest
  set (extinction, carrying capacity, and the frozen-biomass face), and
  integrability of extraction on the finite donor budget (Section 4).
  Positivity is proved face by face via the tangent cone (Aubin, 1991).
  The classical compartmental-systems nonnegativity theory (Jacquez and
  Simon, 1993) is the lineage.
\item
  \textbf{Depletion arithmetic.} Depletion time is separated into three
  non-interchangeable quantities --- gross turnover intensity, a
  frozen-rate ratio, and a scenario-conditioned hitting time ---
  extended to upper barriers, infimum exit times, and infinite-horizon
  maintainability, with uniform-drift bounds proved and the failure of
  ``positive extraction implies finite exhaustion'' exhibited by
  counterexample (Sections 6.1--6.5).
\item
  \textbf{Application classifications at exact status.} The G3P
  anomaly-persistence index, the phosphate reserve-life ratio, and the
  fisheries removals-only pressure time are classified as what they are
  --- statistical index, arithmetic ratio, pressure scale --- against
  the published critiques they corroborate (Section 6.5).
\item
  \textbf{First-passage semantics on declared surrogates.} The
  inverse-Gaussian groundwater passage and the geometric-Brownian
  fisheries passage are stated with complete derivations, the
  record-relative-barrier discipline, and seven explicit non-claims
  separating them from the ledger (Section 7).
\item
  \textbf{The interface with institutional dynamics.} The exact shared
  object with delay-based institutional models is the single-resource
  deficit identity \(qEN - R = -\dot N\). The boundary is the
  non-reduction theorem: no exact dynamic reduction from the closed
  ledger to the open working system, for five stated mathematical
  reasons (Section 9).
\end{enumerate}

\textbf{What is not claimed.} The article is equally explicit about what
it does not establish. No stochastic completion of the ledger is
claimed: the surrogate processes of Section 7 do not conserve the
ledger's mass and are not perturbations of its dynamics. No
thermodynamic admissibility is claimed. No identification of the
two-pool groundwater hypothesis is claimed (its identification
requirements are registered in Section 8, not discharged). And no
empirical finding about any basin, aquifer, or fishery is claimed beyond
the descriptive status of the tabulated indicators.

\textbf{Antecedents.} The bookkeeping idiom and its reconciliation practice are those of material flow
analysis (Section 1.1; Brunner and Rechberger, 2004), and the incidence apparatus is that of
compartmental systems and reaction-network theory (Jacquez and Simon, 1993; Feinberg, 2019). The
regime distinction the ledger formalises is the contested distinction of the weak-versus-strong
sustainability and critical-natural-capital literatures (Ekins et al., 2003; Neumayer, 2013), and
the aggregate-versus-component question is the non-compensatory-aggregation question posed for
composite indicators (Munda and Nardo, 2009), where the footprint accounts that animate part of that
debate are themselves under assessment (Wackernagel and Beyers, 2019; Lin et al., 2018; Blomqvist
et al., 2013). The terminal-condition tradition is viability theory (Aubin, 1991); the horizon
arithmetic sits on a bioeconomic and extractive-price literature (Clark, 1990; Tilton, 2003; Tilton
and Lagos, 2007), as do the growth-limit arguments with which depletion horizons are confused
(Meadows et al., 1972; Daly, 1990). The statistical standards are closer than that genealogy
suggests, and they make this article's point about themselves: the 2025 System of National Accounts
records the depletion of non-produced natural resources as a cost of production, so that net domestic
product is gross product less depreciation and less depletion, where the 2008 revision had recorded
the same depletion as an other change in the volume of assets (United Nations, 2025), a treatment
carried into the framework from the SEEA Central Framework, which standardised the definition and
recording of depletion (United Nations, 2014). No extraction rate changed between the two revisions:
a net aggregate moved because a classification moved, which is the reserve-classification behaviour
this article treats as reclassification, exhibited on the accounts themselves. What the standards
settle is the recording and the valuation of a depletion flow, and their own physical definition of it
--- extraction in excess of the resource's growth --- is a comparison of two rates; the step from a
rate to a horizon is taken nowhere in them, and it is not taken here either. The direction of the
comparison is one-way and is meant to be read that way: the accounts supply an instance of the phenomenon this
article analyses, not a premise for it, and nothing in Sections 2 to 10 depends on any claim made about the
standards here --- not the arithmetic, not the certificates, and not the interpretation boundaries. Their
revision is also silent on the boundary of an asset in physical terms: for the renewable-resource categories
the 2025 revision admits, what counts as an asset is drawn by viability under prevailing technology and
prices, which is an eligibility statement of exactly the kind Remark 36 treats as reclassification rather than
as a property of the stock. None of these sources is used as a premise here.
The objects proved about are the ledger's own.


\subsubsection{1.3 Organization}\label{organization}

Section 2 defines the typed ledger. Section 3 develops the certification
layers and the accounting theorems. Section 4 proves the closed-ledger
theorem set. Section 5 adds the service layer and the componentwise
deficit. Section 6 develops the depletion taxonomy, the uniform-drift
bounds, and the application classifications. Section 7 supplies
first-passage semantics. Section 8 records the domain templates at
registered status. Section 9 fixes the interface with delay dynamics.
Section 10 states what the ledger does not support, and Section 11
concludes.

\begin{center}\rule{0.5\linewidth}{0.5pt}\end{center}

\subsection{2. The Typed Primitive
Ledger}\label{the-typed-primitive-ledger}

\textbf{Notation.} One letter carries one sort wherever a computation is
displayed; section-local aliases are declared where they occur, and the
incidence operator is never written \(N\) (which is reserved for the
living stock). The table sits at the head of Section 2 so that every
alias is declared before use:

\begin{longtable}[]{@{}
  >{\raggedright\arraybackslash}p{(\columnwidth - 4\tabcolsep) * \real{0.3333}}
  >{\raggedright\arraybackslash}p{(\columnwidth - 4\tabcolsep) * \real{0.3333}}
  >{\raggedright\arraybackslash}p{(\columnwidth - 4\tabcolsep) * \real{0.3333}}@{}}
\toprule\noalign{}
\begin{minipage}[b]{\linewidth}\raggedright
Symbol
\end{minipage} & \begin{minipage}[b]{\linewidth}\raggedright
Meaning
\end{minipage} & \begin{minipage}[b]{\linewidth}\raggedright
Where
\end{minipage} \\
\midrule\noalign{}
\endhead
\bottomrule\noalign{}
\endlastfoot
\(N\) & living stock; nutrient stock (local to \S{}2.4) & \S{}2.2; \S{}2.4 \\
\(S_{\mathcal{T}}\) & typed stoichiometric (incidence) operator & \S{}2.1,
Lemma 3, Proposition 4, Theorem 5, Theorem 8 \\
\(v\) & non-negative primitive flux vector & \S{}2.1 \\
\(S\) & moiety readout \(S = Cx\) & Lemma 3, \S{}6 \\
\(s\) & support factor \(A^{\mathrm{act}}/(A^{\mathrm{act}}+A_0)\) &
\S{}2.2 \\
\(\sigma\) & donor fraction
\(A^{\mathrm{geo}}/(A^{\mathrm{geo}}+A_{g0})\) & \S{}2.2 \\
\(\varsigma\) & noise scale of the stochastic surrogates & \S{}7 \\
\(M\) & natural-block mass
\(N + A^{\mathrm{act}} + A^{\mathrm{geo}} + U\) & Theorems 7, 14 \\
\(\widehat{M}\) & demand-coverage matrix of the physical deficit &
\S{}5.4 \\
\(K\) & carrying capacity; sink stock (local to \S{}2.4); maintainability
kernel \(K_{\mathrm{maint}}\) (subscripted) & \S{}2.2; \S{}2.4; \S{}6.3 \\
\(T\) & gross uptake \(\kappa_A N s\); finite horizon (local to each
statement) & \S{}2.2; Theorems 1, 5 \\
\(C\) & moiety-composition matrix; operative extraction-law readout;
mining intensity \(C^A\) & Lemma 3; \S{}5.4; \S{}2.2 \\
\(B\) & \(R + T\); barriers (local to \S{}3.1); aggregate regeneration flow
\(b \cdot M\) (local to \S{}1.1); fisheries biomass \(B_t\) and barrier
\(B_{\min}\) (local to \S{}7.5); reference \(B_{\lim}\) (local to \S{}6.5.2,
\S{}6.5.4) & \S{}2.2; \S{}3.1; \S{}1.1; \S{}6.5.2 \\
\(b\) & boundary-transfer term; service balance \(b_i\) (subscripted);
specific regeneration rate (local to \S{}1.1) & Lemma 3, Proposition 4,
Theorem 5; \S{}5.1; \S{}1.1 \\
\(R\) & net regeneration; log-margin \(R_B\) (local to \S{}6.5.4) & \S{}2.2;
\S{}6.5.4 \\
\(r\) & intrinsic growth rate; net boundary inflow \(r_m\) (subscripted,
\S{}3.3) & \S{}2.2; \S{}2.6; \S{}3.3 \\
\(G\) & geological pool (scaffold); reserves (local to \S{}6.5.3, \S{}7.6);
donor stock \(G_0\) (subscripted) & \S{}2.3; \S{}6.5.3; \S{}9 \\
\(I\) & inert sink; boundary input \(I_N\) (subscripted) & \S{}2.2; \S{}2.4 \\
\(z\) & six-compartment state (local to \S{}2.3) & \S{}2.3 \\
\(h\) & harvest primitive; net boundary outflow \(h_m\) (subscripted,
\S{}3.3); geometric-Brownian drift (local to \S{}7.5) & \S{}2.3; \S{}3.3; \S{}7.5 \\
\(x\) & ledger state & \S{}2.1, Lemma 3, Proposition 4, Theorem 5 \\
\(y\) & declared boundary states feeding the primitive fluxes & \S{}2.1,
eq. (1) \\
\(\theta\) & declared constitutive parameter vector; sink-generation
fraction \(\theta_K\) (subscripted); parameter set \(\Theta\) and
fisheries pressure time \(\Theta_F\) (distinct objects) & \S{}2.1, eq. (1);
\S{}2.4; \S{}6.4, \S{}6.5.4 \\
\(E\) & extraction effort, \(E \ge 0\) along classical solutions & \S{}2.2;
\S{}4.5; Theorem 14 \\
\(\alpha\) & harvest routing fraction; directional support fraction
\(\alpha_{\mathrm{reg}}\) (subscripted) & \S{}2.2--\S{}2.3, \S{}4.2; \S{}5.3 \\
\(\rho_P\) & product-retirement fraction routing \(r_P\) to \(U\) versus
\(W\) & \S{}2.3; \S{}4.3; \S{}4.4; \S{}8.1 \\
\(\mu, \nu, \rho\) & product, waste, and price parameters of the
unreduced ledger (zero in the single-resource specialization); \(\mu\)
also the growth parameter of Theorem 1 and the drift of the \S{}7
surrogates; \(\nu\) also the inverse-Gaussian mean (local to \S{}7.3) &
\S{}2.2; \S{}5.4; \S{}9 \\
\(\varepsilon\) & drift bracket (Proposition 17); probabilistic level
(\S{}6.4); resource-threshold fraction (\S{}6.5.2, \S{}7.6); donor-draw
diagnostic \(\varepsilon_G\) (subscripted); slack (\S{}10.1) & \S{}6.2; \S{}6.4;
\S{}6.5.2; \S{}9; \S{}10.1 \\
\(\tau\) & hitting and exit times (\(\tau_B\), \(\tau_m^{\pm}\),
\(\tau_{\mathrm{exit}}\)) & \S{}3.6; \S{}6.3 \\
\(d\) & disturbance (\S{}2.1); demand vector (\S{}5.2); drift distance (local
to \S{}7.3) & \S{}2.1; \S{}5.2; \S{}7.3 \\
\(A_0\) & half-saturation constant (\S{}2.2); latest observed anomaly
(local to \S{}7.2--\S{}7.4) & \S{}2.2; \S{}7.2 \\
\(\lambda\) & inverse-Gaussian shape parameter & \S{}7.3 \\
\(\chi, \eta\) & hybrid state and primitive-flux vector of Conditional
Theorem 15 & \S{}4.8 \\
\(P\) & product compartment; production rate (local to \S{}6.5.3, \S{}7.6) &
\S{}2.2--\S{}2.3; \S{}6.5.3 \\
\end{longtable}

\subsubsection{2.1 Typed stocks, primitive fluxes, and the incidence
discipline}\label{typed-stocks-primitive-fluxes-and-the-incidence-discipline}

Material flow accounting begins with a small set of distinctions.
Heterogeneous substances must be tracked separately. Spatial and
functional locations must be distinguished. Conversions between chemical
forms must be made explicit. This section sets out the incidence
discipline that makes conservation a property of the structure rather
than an assumption.

The ledger separates four concepts that accounting practice often
merges. A \emph{moiety} is a conserved substance class (an element, or a
declared conserved combination) --- the only object to which a
conservation law attaches. A \emph{species} is a chemical or biological
form of a moiety. A \emph{compartment} is a spatial or functional
location holding a stock. A \emph{stock} is a compartment's current
amount of a species, with a physical unit. ``Carbon in the atmosphere''
is a location-specific stock, not a moiety: carbon is the moiety, and
the atmosphere is a compartment. Conservation laws are stated per
moiety; nothing is conserved merely by being a compartment.

A ledger state \(x \in \mathbb{R}^m_+\) collects compartments, each
entry carrying a material identity, spatial support, and physical unit.
Internal dynamics use non-negative primitive fluxes:
\[\dot x = S_{\mathcal{T}} v(x, y, \theta) + B_{\mathcal{T}} u_{\partial}(t) + d_x(t), \qquad v \ge 0, \tag{1}\]
where \(S_{\mathcal{T}}\) is the typed stoichiometric (incidence)
operator, \(y\) collects the declared boundary states --- environmental
or companion variables external to the ledger --- \(\theta\) the
declared constitutive parameter vector, \(B_{\mathcal{T}} u_{\partial}\)
collects declared boundary transfers, and \(d_x\) belongs to a stated
disturbance class. Entries are added within a row only when their types and units agree; a conversion between types is
represented by an explicit stoichiometric coefficient, never by an implicit sum. That sentence presumes a
declaration, and the declaration is Definition 47: types, units and conversion coefficients are declared with
the ledger, and \(S_{\mathcal{T}}\) carries the subscript for that reason. If
\(L^\top S_{\mathcal{T}} = 0\), then
\[\frac{d}{dt}(L^\top x) = L^\top B_{\mathcal{T}} u_{\partial} + L^\top d_x,\]
one conservation law per conserved moiety and boundary; the identity does not create a scalar
sustainability mass across incommensurable systems, and Proposition 42 states exactly in what sense it cannot. Three clarifications are part of the statement. First, \(d_x\)
must itself be typed: a physical disturbance on represented material is
a different object from a structural discrepancy term. Second,
\(S_{\mathcal{T}}\) may contain signed entries even though \(v \ge 0\):
the sign pattern of the incidence matrix and the nonnegativity of the
primitives are separate declarations. Third, forward invariance is a
separate requirement: every primitive outflow must vanish or be limited
when its donor compartment is empty, and a target-relaxation flux from a
finite donor is admissible only after donor limitation is made explicit.

\textbf{Definition 21 (Closure cone).} The regimes above can be stated on the ledger's own objects rather
than as attitudes toward them. Let \(S_{\mathcal{T}}\) be the incidence operator declared here, let \(v\ge0\)
be the primitive fluxes under declared capacities \(\bar v\), let \(B_Tu_\partial\) collect the
declared boundary transfers, and let \(P\) project a flux vector onto its use columns. Write
\(\mathcal K=\{v\ge0:\ S_Tv+B_Tu_\partial=0,\ v\le\bar v\}\) for the admissible stationary flux
patterns. A demanded use vector \(D\) closes at the rate of use if and only if
\(D\in P(\mathcal K)\). The closure capacity
\(\Lambda^{*}=\max\{\mu:\ \mu D\in P(\mathcal K)\}\) is a linear programme on the declared
capacities, and its feasibility is the cut condition on return capacity (Gale, 1957; Ahuja et al.,
1993). Where \(\Lambda^{*}\ge1\), the cycle closes at the demanded rate. Where \(\Lambda^{*}<1\),
every trajectory meeting \(D\) has a non-stationary stock, and by conservation the shortfall appears
as support drawdown together with sink accumulation; the deficit share is not a free parameter. Yield
inflation is the same statement read outside the projection: demand met with no stationary pattern
behind it.

Left-null vectors express conservation; right-null vectors express circulation. The regimes of
Section 1.1 are statements about the second set, and the depletion readouts of Section 6 are
statements about the first; neither set's membership follows from the other's.

\textbf{Definition 22 (Closure deficit at the use timescale).} For each moiety \(m\), let
\(\kappa_m(\tau;\Theta)\in[0,1]\) be the fraction of mobilised flux returned to a usable compartment
within lag \(\tau\) under the declared technology and process graph \(\Theta\), and let
\(\delta_m=1-\kappa_m(\tau_{\mathrm{use}};\Theta)\) be the closure deficit at the declared use
timescale. The deficit is a property of a declared graph at a declared timescale, not a rate of the
material itself. Where a mobilisation \(g_m\) is sustained with \(\delta_m\ge\underline\delta>0\),
the net drawdown of the support pool is \(\delta_mg_m\), so
\[ \int_0^T(\text{liquidation})\,dt\ \ge\ \underline\delta\int_0^Tg_m\,dt, \qquad T\ \le\ \frac{m_0}{\underline\delta\,g_m}, \]
the second inequality being the frozen-rate horizon read on the liquidation flux rather than on gross
use. Both bounds are one-sided by construction, and neither reverses in the limit
\(\delta_m\to0^+\). A positive deficit is not a barrier-reachability claim: a deficit that decays
with the stock need never reach a barrier, and where the binding constraint is an upper barrier the
relevant object is the complementary headroom.

Two readings follow from the definition. An indicator that reports a return fraction with no lag
argument reports \(\kappa_m(\infty)\) where \(\kappa_m(\tau_{\mathrm{use}})\) is being read; the
recycled-input and circularity-rate families are of that form, since a ratio of fluxes carries no
timescale. And "a material is waste" is, on this reading, the statement
\(\kappa_m(\tau_{\mathrm{use}};\Theta)<1\) for the declared graph, so unlocking a stream is an
increase of \(\kappa\) at fixed \(\tau\): measurable, and not in need of a new substance category.


\subsubsection{2.2 The closed finite-donor
ledger}\label{the-closed-finite-donor-ledger}

The closed ledger of this article is the finite-donor primitive system.
Let
\[x_L = (N, A^{\mathrm{act}}, A^{\mathrm{geo}}, U), \qquad s = \frac{A^{\mathrm{act}}}{A^{\mathrm{act}} + A_0}, \qquad \sigma = \frac{A^{\mathrm{geo}}}{A^{\mathrm{geo}} + A_{g0}},\]
with \(N\) the living stock, \(A^{\mathrm{act}}\) the active abiotic
pool, \(A^{\mathrm{geo}}\) the geological donor, and \(U\) the detritus
compartment. Net regeneration and gross uptake are the constitutive laws
\[R(N, A^{\mathrm{act}}) = rN\left(1 - \frac{N}{K}\right)s, \qquad T = \kappa_A N s, \qquad B = R + T,\]
and the four geo-interface and closure primitives are
\[e_{GA} = \omega_A A^{\mathrm{eq,intrinsic}} \sigma, \qquad e_{AG} = \omega_A A^{\mathrm{act}}, \qquad C^{A,\mathrm{lim}} = C^A \sigma, \qquad \gamma_U U \ \text{(detritus return)}.\]

These donor primitives instantiate one of three recharge laws that
appear across this article and the companion delay-dynamics analysis (Abaee, 2026, doi:10.5281/zenodo.22554217); the three are distinct objects,
tabulated once so that none is silently substituted for another:

\begin{longtable}[]{@{}
  >{\raggedright\arraybackslash}p{(\columnwidth - 4\tabcolsep) * \real{0.3333}}
  >{\raggedright\arraybackslash}p{(\columnwidth - 4\tabcolsep) * \real{0.3333}}
  >{\raggedright\arraybackslash}p{(\columnwidth - 4\tabcolsep) * \real{0.3333}}@{}}
\toprule\noalign{}
\begin{minipage}[b]{\linewidth}\raggedright
Recharge law
\end{minipage} & \begin{minipage}[b]{\linewidth}\raggedright
Form
\end{minipage} & \begin{minipage}[b]{\linewidth}\raggedright
Status
\end{minipage} \\
\midrule\noalign{}
\endhead
\bottomrule\noalign{}
\endlastfoot
Primitive donor-limited exchange &
\(e_{GA} = \omega_A A^{\mathrm{eq,intrinsic}} \sigma\) & The closed
block's law (Section 2.2): the forward rate depends on the donor alone,
not on how empty the receiving pool is --- linear donor-limited
exchange, whose rest is
\(A^{\mathrm{act}} = A^{\mathrm{eq,intrinsic}}\sigma\) (Theorem 13). \\
Target-relaxation & \(\omega_A(A^{\mathrm{eq}} - A)\) & Banned unless
donor-limited (Section 4.4): runs backward at an empty donor; admissible
only with the source declared an effectively infinite external
reservoir, making the system open. \\
Working derived target &
\(A^{\mathrm{eq,W}} = A^{\mathrm{eq,intrinsic}} + \kappa_A K/\omega_A\)
& The delay-dynamics analysis's working completion (that analysis; this
article's Section 9): not a closed-block law, and the reason the two
systems do not reduce. \\
\end{longtable}

In the closed block no derived target appears. Recharge is donor-limited
and cannot run backward (\(e_{GA} = 0\) at \(A^{\mathrm{geo}} = 0\); the
positive-part convention \([\cdot]_+\), read as a one-way valve at a
nonpositive target, never binds here because the registered intrinsic
target is positive), and mining \(C^{A,\mathrm{lim}}\) is donor-limited
the same way extraction is. With \(A_{g0} > 0\) the donor fraction
\(\sigma\) is smooth and strictly increasing in the donor level.

Under the institutional-failure specialization (\(\mu = \nu = \rho = 0\)
and \(C^A = 0\) --- the product, waste, and price parameters
\(\mu, \nu, \rho\) of the unreduced ledger, its macroeconomic-feedback,
recycling, and price-response channels, set to zero together with the
mining intensity \(C^A\); the parameters are glossed at their Section
5.4 site) the closed natural block is \[\begin{aligned}
\dot N &= R - qEN, \\
\dot A^{\mathrm{act}} &= -B + e_{GA} - e_{AG} + \gamma_U U, \\
\dot A^{\mathrm{geo}} &= -e_{GA} + e_{AG}, \\
\dot U &= T - \gamma_U U,
\end{aligned} \tag{2}\] together with a memory--effort pair \((Z, E)\)
driven by \(qEN - R\) (never by mining; \(q\) the per-effort extraction
coefficient of the harvest law). The pair is the registered object of
the companion delay-dynamics analysis (Abaee, 2026, doi:10.5281/zenodo.22554217; its eq. (1), where the memory--effort pair is defined on the
gated three-state core with the stock-decline rate as memory input, and its Section 2.4, which relates that
core to the four-state working model in which the input is \(qEN - R(N, A)\)) and is not analysed in this
article. Net regeneration is the difference of two non-negative
primitives --- gross regeneration \(rNs\) (support \(\to\) stock) and
density-dependent return \(rN^2 s/K\) (stock \(\to\) support) --- so (2)
stays within the primitive-flux discipline of Section 2.1 despite the
signed entry. The block's harvest routing is the \(\alpha = 0\) corner
of Section 2.3: harvest \(qEN\) exits the natural block entirely as
product, and a positive detritus-routed fraction \(\alpha > 0\) would
add \(\alpha qEN\) to \(\dot U\) and reduce the block export to
\((1-\alpha)qEN\); the mass identity of Theorem 7 is stated for the
declared routing. The registered parameterization is \(r = 0.02\),
\(K = 100\), \(q = 0.001\), \(\kappa_A = 0.05\), \(\omega_A = 10^{-3}\),
\(A_0 = 1\), \(A^{\mathrm{eq,intrinsic}} = 50\), \(\gamma_U = 0.2\); the
geological half-saturation \(A_{g0}\) is declared positive (smoothness
of the donor fraction \(\sigma\)) under the separation-of-scale
condition \(A^{\mathrm{geo}} \gg A_{g0}\), in which regime
\(\sigma \approx 1\); the scale separation is registered rather than a
numerical value, and the \(A_{g0} = 0\) corner is the
discontinuous-perturbation limit, not the registered regime. With
\(A_0 > 0\) and \(A_{g0} > 0\) the right-hand side of (2) is locally
Lipschitz on the closed orthant, and the comparison
\(\dot N \le rN(1 - N/K)\) (with \(\dot N \le 0\) once \(N \ge K\))
bounds the stock by \(\max\{N(0), K\}\). Classical solutions therefore
exist globally and stay in the orthant by Theorem 10 --- the existence
clause behind every `classical solution' statement below.

The incidence discipline of Section 2.1, written out for the closed
block, is the four-row block incidence --- rows
\((N, A^{\mathrm{act}}, A^{\mathrm{geo}}, U)\), columns gross
regeneration \(rNs\), density-dependent return \(rN^2 s/K\), harvest
\(qEN\), uptake \(T = \kappa_A N s\), detritus return \(\gamma_U U\),
geological recharge \(e_{GA}\), geological return \(e_{AG}\), mining
\(C^{A,\mathrm{lim}}\): \[S_{\mathrm{block}} =
\begin{pmatrix}
1 & -1 & -1 & 0 & 0 & 0 & 0 & 0 \\
-1 & 1 & 0 & -1 & 1 & 1 & -1 & 0 \\
0 & 0 & 0 & 0 & 0 & -1 & 1 & -1 \\
0 & 0 & 0 & 1 & -1 & 0 & 0 & 0
\end{pmatrix}.\] Every column is a two-compartment transfer or a
block-boundary export: the six internal columns sum to zero, and the two
exports --- harvest and mining --- carry the column sums \(-1\) each, so
summing the four rows reads off \(\dot M = -qEN - C^{A,\mathrm{lim}}\),
the mass identity proved as Theorem 7, directly from the display (under
the institutional-failure specialization the mining column is inactive,
\(C^A = 0\)). The harvest column is written at the declared
\(\alpha = 0\) routing, under which harvest exits the block; a
detritus-routed fraction \(\alpha > 0\) moves \(\alpha\) into the \(U\)
entry of that column and reduces the block export to \((1-\alpha)qEN\)
--- the full-ledger routing displayed in Section 4.2.

When product, waste, and the inert sink are restored with the same
donor-limited routing, the full ledger
\(N + A^{\mathrm{act}} + A^{\mathrm{geo}} + U + P + W + I\) is closed
(Section 4.2). The geological donor is an internal state throughout ---
no infinite reservoir is declared --- and the block boundary is crossed
only by harvest (to product) and, when restored, mining (to product or
waste). The full seven-compartment incidence is the nested completion of
the four-block routing. Uptake \(T = \kappa_A N s\) never enters
\(\dot N\): it is an \(A^{\mathrm{act}} \to U\) throughput with the
living stock entering as a catalytic factor, not stored biomass ---
legitimate in a monomaterial projection; the masses below are defined on
this routing.

\subsubsection{2.3 The six-compartment
illustration}\label{the-six-compartment-illustration}

For one conserved limiting material, the accounting scaffold is
instantiated by six compartments --- living biomass \(X\), detritus or
recoverable residual \(U\), active abiotic pool \(A\), geological or
slowly available pool \(G\), product or in-use stock \(P\), and
absorbing or currently unavailable stock \(W\) --- with eight
non-negative primitive fluxes: assimilation \(g(X,A)\), mortality
\(m(X)\), harvest \(h(X,E)\), decomposition \(d_U(U)\),
geological-to-active transfer \(e_{GA}(G,A)\), active-to-geological
transfer \(e_{AG}(A,G)\), direct mining \(c_G(G,E_G)\), and product
retirement \(r_P(P)\). With harvest fraction \(\alpha \in [0,1]\) routed
to \(U\) and retirement fraction \(\rho_P \in [0,1]\) returning to \(U\)
rather than \(W\),
\[\dot z = S(\alpha, \rho_P)\, v(z, u), \qquad z = (X, U, A, G, P, W)^\top, \qquad v = (g, m, h, d_U, e_{GA}, e_{AG}, c_G, r_P)^\top,\]
where \[S(\alpha, \rho_P) =
\begin{pmatrix}
1 & -1 & -1 & 0 & 0 & 0 & 0 & 0 \\
0 & 1 & \alpha & -1 & 0 & 0 & 0 & \rho_P \\
-1 & 0 & 0 & 1 & 1 & -1 & 0 & 0 \\
0 & 0 & 0 & 0 & -1 & 1 & -1 & 0 \\
0 & 0 & 1-\alpha & 0 & 0 & 0 & 1 & -1 \\
0 & 0 & 0 & 0 & 0 & 0 & 0 & 1-\rho_P
\end{pmatrix}, \qquad \mathbf{1}^\top S = 0.\] The zero column sums are
the incidence statement of mass conservation (proved for the full system
in Section 4.2). The matrix makes the routing choices visible, and the
constitutive choices are features of this example, not properties of
every typed ledger: the constant splits \(\alpha\) and \(\rho_P\), the
compartment set, and the absorbing-sink convention are declared choices.
The construction is a monomaterial projection. Coupled multi-element
accounts require additional typed rows and a conservation matrix. If an
application makes recovery claims, \(U\) and \(P\) must be split by
material, location, and quality grade, with declared yields, residual
routes, and exergy or capacity inputs --- the conservation argument then
applies to the expanded typed incidence system, not automatically to an
undifferentiated quality-neutral loop. The four-block system (2) is not
a specialization of this scaffold: in the scaffold assimilation \(g\) is
a slow \(A \to X\) flux and mortality \(m\) a slow \(X \to U\) flux,
while in (2) the uptake \(T\) transfers \(A^{\mathrm{act}} \to U\) with
the living stock catalytic and no separate mortality primitive --- two
different timescale lumpings of the same physical story, and no
incidence specialization maps one onto the other.

\subsubsection{2.4 The four-stock resource--sink--nutrient--product
system}\label{the-four-stock-resourcesinknutrientproduct-system}

A second exact specialization closes a resource--sink system with a
nutrient stock and product stock. The state is
\((S, K, N, P) \in \mathbb{R}^4_+\) (in this block's local notation,
\(S\) is the resource stock, \(K\) the sink stock, \(N\) the nutrient
stock, and \(P\) the input flux; the carrying capacity and living stock
of Section 2.2 do not enter), with
\[\dot S = g(S,N) - H, \qquad \dot K = \theta_K H - \theta_\delta K, \qquad \dot N = -g(S,N) + \theta_\delta K + I_N, \qquad \dot P = (1 - \theta_K)H - Q_P,\]
where \(\theta_K\) is the sink-generation fraction, \(\theta_\delta\)
the assimilation rate, \(I_N\) external nutrient input, and \(Q_P\)
product disposal. Adding the four equations gives the mass balance
\[\frac{d}{dt}(S + K + N + P) = I_N - Q_P,\] so total mass is conserved
exactly when both boundary transfers vanish. This mass balance is an
exact specialization of the incidence discipline of Section 2.1: every
internal transfer cancels in the column sum, and the boundary terms
survive as the ledger's declared inputs and outputs.

\textbf{Sink obstructions independent of the stock.} The mass balance
has a sink-side physical reading with two empty-kernel mechanisms that
operate whatever the resource stock does. With sink loading \(w(H)\),
assimilation \(\delta(K)\), and a harvest floor \(H \ge H_{\min} > 0\)
(in the four-stock specialization, \(w(H) = \theta_K H\) and
\(\delta(K) = \theta_\delta K\)): under \emph{no assimilation}
(\(\delta \equiv 0\)), \(\dot K \ge w(H_{\min}) > 0\), and the sink
exceeds any finite ceiling \(K_{\max}\) in finite time; under \emph{weak
assimilation} (\(\delta(K_{\max}) < w(H_{\min})\)), the sink load at the
ceiling is still positive ---
\(\dot K = w(H_{\min}) - \delta(K_{\max}) > 0\) at \(K = K_{\max}\) ---
so \(K\) exits above \(K_{\max}\) in finite time, the explicit negation
of the ceiling condition \(\delta(K^\dagger) = w(H_{\min})\) with
\(K^\dagger \le K_{\max}\). In both cases the viability kernel (Aubin,
1991) of the constraint set
\(\{S \ge S_{\min},\ 0 \le K \le K_{\max}\}\) --- the states from which
some admissible harvest keeps both constraints for all time --- is
empty. The obstruction needs the sink-generation fraction
\(\theta_K > 0\): if \(\theta_K = 0\), harvest never loads the sink and
the loading argument does not apply --- emptiness would then have to
come from the resource constraint or an undeclared ceiling on the
product stock.

The closed-ledger corollary is the same mechanism in ledger language. In
a closed ledger without recycling, where \(w(H)\) enters the sink
irreversibly and \(\delta = 0\), the sink rises monotonically against
the finite total mass, and any positive output floor forces an empty
viability kernel. Which constraint fails first depends on the total mass
\(M\), the ceiling \(K_{\max}\), and the remaining stock: a ceiling
below the sink's reachable mass share is crossed in finite time, while a
ceiling at or above the total mass is unreachable and the violated
constraint is the resource floor (the finite-budget bound of Theorem
14). For material-flows accounting this is the precise sense in which a
``balanced'' mass ledger can still be physically inadmissible: the
bookkeeping is exact, but no harvest schedule respects both the resource
floor and the sink ceiling simultaneously.

\subsubsection{2.5 Mechanism typing: routing is never determined by
diagnostic
labels}\label{mechanism-typing-routing-is-never-determined-by-diagnostic-labels}

Extraction has at least three distinct physical meanings ---
standing-stock culling (present extraction removes reproductive stock
directly), recruitment suppression (present use prevents future recruits
without immediate adult removal), and weak viability coupling (use has
limited or indirect effect on reproduction). In the ledger,
standing-stock culling enters as an outflow from the standing-stock
compartment. The typing is the physical module's, not the diagnostic's:
a diagnostic label such as ``unsustainable portion'' never determines
physical destination. Material routing is determined by the typed
physical module alone, and the diagnostic threshold that flags a flow
has no standing in the incidence matrix.

\textbf{The split-assignment evidence requirement.} Where an application
splits extraction between standing-stock removal and recruitment
suppression --- \(C_{\mathrm{stock}} = \psi qEN\) and
\(C_{\mathrm{recruit}} = (1-\psi)qEN\) for \(\psi \in [0,1]\) --- the
assignment requires evidence per channel. The dominant physical
mechanism sets it, never a diagnostic label. Illustrative assignments
through a logistic two-channel proxy (stated as illustrative, calibrated
examples, not constitutive claims for the named domains): soil zinc
under crop export, an existing-unit removal, at \(\psi = 0.85\), against
impaired mineralisation, a replenishment degradation, at
\(\psi = 0.25\); pollinators under adult mortality at \(\psi = 0.70\)
against brood failure at \(\psi = 0.20\) --- and across such mechanism
pairs the trough depth varies by a factor of about \(1.5\) from
mechanism alone. The mass-routing discipline is the typing made
explicit. Literally harvesting pre-recruit stages is a harvest of
existing units and routes to the product and waste fractions.
Habitat-induced failed recruitment is a prevented inflow --- routing it
into product or waste would create mass that was never in the stock.
Damage to the capital stock itself (aquifer compaction, severe soil
loss) is not a split-assignment channel at all, but a slow drift in
capacity or a transfer to the inert sink.

\subsubsection{2.6 Support saturation and the logistic
limit}\label{support-saturation-and-the-logistic-limit}

Two results control what the ledger's stock equation becomes when its
support pool saturates. Both are singular reductions with explicit
scope, and neither is a full-system reduction.

\textbf{Theorem 1 (Support-saturated logistic stock limit).} \emph{Fix
\(T < \infty\) and non-negative parameters \(\mu, \delta, c, q\). For
\(\kappa > 0\), assume:}

\emph{(H1) \(A_\kappa\) is measurable with \(A_\kappa(t) \ge a_0 > 0\)
and \(0 \le X_\kappa(t) \le X_{\max}\).}

\emph{(H2) Common effort \(E \in L^\infty([0,T])\).}

\emph{Let \(X_\kappa\) solve
\(\dot X_\kappa = \mu X_\kappa A_\kappa/(\kappa + A_\kappa) - \delta X_\kappa - cX_\kappa^2 - qE(t)X_\kappa\)
and \(X_0\) solve the limiting equation with the same initial value.
Then \(\sup_{t \le T} |X_\kappa(t) - X_0(t)| = O(\kappa)\); if
\(\mu > \delta\) and \(c > 0\) the limit is
\(\dot X_0 = rX_0(1 - X_0/K_{\log}) - qE(t)X_0\) with
\(r = \mu - \delta\) and \(K_{\log} = (\mu - \delta)/c\).}

\emph{Proof.} Write \(e(t) = |X_\kappa(t) - X_0(t)|\). The saturation
defect satisfies
\[\left| \frac{A_\kappa}{\kappa + A_\kappa} - 1 \right| = \frac{\kappa}{\kappa + A_\kappa} \le \frac{\kappa}{a_0},\]
so the vector-field defect obeys
\[|\dot X_\kappa - \dot X_0| \le L_1 \kappa + L_2 e(t), \qquad L_1 = \frac{\mu X_{\max}}{a_0}, \qquad L_2 = \mu + \delta + 2cX_{\max} + q\|E\|_\infty,\]
using \(|\mu - \delta| \le \mu + \delta\) and
\(c(X_\kappa + X_0) \le 2cX_{\max}\). Gronwall's inequality with the
Lipschitz constant \(L_2\) and the particular-term scale \(L_1\) gives
\(e(t) \le (L_1/L_2)\kappa\,(e^{L_2 t} - 1) \le C_T \kappa\) on
\([0,T]\). The bound \(0 \le X_\kappa \le X_{\max}\) is satisfiable in
the registered family: under \(\mu > \delta\) the comparison
\(\dot X_\kappa \le X_\kappa(\mu - \delta - cX_\kappa)\) keeps
\(X_\kappa \le \max\{X(0), (\mu-\delta)/c\}\), so
\(X_{\max} = \max\{X(0), K_{\log}\}\) suffices. The limiting equation is
\(\dot X_0 = \mu X_0 - \delta X_0 - cX_0^2 - qE X_0\),
i.e.~\(\dot X_0 = (\mu - \delta)X_0 (1 - X_0/K_{\log}) - qE X_0\) with
\(K_{\log} = (\mu - \delta)/c\). \ensuremath{\square}

\textbf{Remark 2 (Registered-family support-saturated identity).}
\emph{In the primitive-flux core with \(g(X,A) = \mu XA/(K_A + A)\),
\(m(X) = dX + cX^2\), \(h(X,E) = qEX\), the support-saturated stock
equation is, for each fixed interior \(A > 0\) in the limit
\(K_A \to 0\),}
\[\dot X = (\mu - d)X - cX^2 - qEX = rX\left(1 - \frac{X}{K}\right) - qEX, \qquad r = \mu - d, \quad K = \frac{\mu - d}{c},\]
\emph{requiring \(\mu > d\) and \(c > 0\). The identity is pointwise on
the interior support region and not uniform through the depleted-pool
boundary: for every \(K_A > 0\), \(A/(K_A + A) = 0\) at \(A = 0\).}

\emph{Proof.} At fixed \(A > 0\), \(A/(K_A + A) \to 1\) as
\(K_A \to 0\), so \(g \to \mu X\) pointwise and
\(g - m \to (\mu - d)X - cX^2\); the algebraic reduction to the logistic
form with \(r = \mu - d\), \(K = (\mu - d)/c\) is immediate. The
non-uniformity statement is the identity \(A/(K_A + A) = 0\) at
\(A = 0\) for every \(K_A > 0\), which the pointwise limit does not
touch. The scope is restricted: the limit does not eliminate the
detritus compartment \(U\), does not make \(A\) constant near its
boundary, and does not transform the memory or effort laws --- an
ecological stock-equation identity, not a full-system reduction and not
a transfer principle for bifurcation thresholds. In the correspondence,
the logistic law of (2) is this saturated, mortality-folded readout of
the Theorem 1 and Remark 2 family --- not the vector field of the closed
four-tuple of Section 2.2 --- and bifurcation numbers of the two
families do not transfer between them. \ensuremath{\square}

The implication for material-flows accounting is that a logistic stock
equation is a saturated support-pool readout, not a primitive physical
law: substituting it for the underlying stock-support dynamics is
legitimate only where the support pool is at interior saturation and
only on the timescales over which the saturated approximation holds.

\begin{center}\rule{0.5\linewidth}{0.5pt}\end{center}

\subsection{3. Certification Layers and the Accounting
Theorems}\label{certification-layers-and-the-accounting-theorems}

\subsubsection{3.1 Three predicates,
separated}\label{three-predicates-separated}

The four predicates of Section 1.1 --- bookkeeping balance,
stoichiometric conservation, thermodynamic admissibility, and
sustainability safety --- are this section's working layer: it separates
the first, second, and fourth and proves their relationships (the third,
thermodynamic admissibility, is out of scope here, per Proposition 2's
layering), for the reasons Section 1.1 states --- a point at which the
Daly/Ayres tradition of throughput accounting meets the formalism of
reaction-network theory (Feinberg, 2019; Brunner and Rechberger, 2004).

The ledger supports three distinct predicates that are related but not
identical.

\textbf{Layer 1: Accounting consistency.} The balance law (1) holds
almost everywhere on \([0,T]\).

\textbf{Layer 2: Conservation consistency.}
\(\ell^\top S_{\mathcal{T}} = 0\) for every declared conserved quantity
\(\ell\).

\textbf{Layer 3: Barrier safety.} For declared lower and upper barriers
\(\underline{B}_m(t) \le S_m(t) \le \overline{B}_m(t)\) for every
component \(m\) and every \(t \in [0,T]\), where \(S = Cx\) is the
moiety-composition readout.

Layer 2 is a structural predicate on the incidence operator alone.
Layers 1 and 3 are properties of a trajectory triple \((x, v, b)\). The
logical relations are the content of the next two propositions.
(Numbering note: the two layering propositions of this section carry
their own counter, Propositions 1--2; every other numbered statement of
this article runs on the single 1--47 sequence counter, so ``Proposition
4'', ``Proposition 6'', ``Proposition 17'', ``Proposition 18'', and
``Proposition 20'' are main-counter statements, not further members of
the layering counter ; the statements added at this revision carry the consecutive labels Definitions 21--23 and 34--35 and 38 and 40--47, Lemma 4, Theorem 24, and Propositions 25--32 and 36--37 and 39--43, with Remarks 33--37, so no label is repeated --- so every label is unique and resolves directly;
the supplementary's statement inventory and this status-word offset are
reconciled in its S6.)

\textbf{Proposition 1.} \emph{Conservation consistency implies
accounting consistency for the conserved quantities: if
\(\ell^\top S_{\mathcal{T}} = 0\), then
\(\frac{d}{dt}(\ell^\top x) = \ell^\top b\), and in a closed system
(\(b = 0\)), \(\ell^\top x\) is invariant.}

\emph{Proof.}
\(\frac{d}{dt}(\ell^\top x) = \ell^\top \dot x = \ell^\top(S_{\mathcal{T}} v + b) = (\ell^\top S_{\mathcal{T}}) v + \ell^\top b = \ell^\top b\).
\ensuremath{\square}

\textbf{Proposition 2.} \emph{Barrier safety does not follow from
accounting consistency: a trajectory can be perfectly mass-balanced
while violating a declared barrier; and a trajectory can satisfy
declared barriers while violating a conservation law or a stoichiometric
constraint. Conversely, thermodynamic admissibility (energy
conservation, entropy-production nonnegativity, reaction feasibility)
implies accounting consistency, but the converse does not hold; this
article establishes Layers 1--3 only.}

\emph{Proof.} First clause: on the closed ledger (2), constant
extraction at a rate exceeding regeneration --- a comparison flux, not a
donor-limited primitive of the ledger's own discipline --- is exactly
mass-balanced (Theorem 7 states the identity) and drives the living
stock through any positive lower barrier in finite time; likewise the
trajectory with \(N \equiv 0\) (extinction rest of Section 4.6) is
mass-balanced yet violates any positive lower barrier on \(N\).
Conversely, a trajectory satisfying the barriers may be generated by
fluxes that fail a stoichiometric constraint --- mass balance does not
certify that the flux decomposition is physically realizable --- or by
bookkeeping that silently drops a moiety through the yield-routing
violation of rule (iii) in Section 10.2, failing conservation without
touching the barriers. Second clause: thermodynamic admissibility
presupposes a mass balance, but a mass-balanced flux decomposition need
not satisfy energy or entropy constraints; establishing those requires
structure outside the present scope. \ensuremath{\square}

The implication for industrial-ecology measurement is that mass-balance
closure, stoichiometric consistency, and barrier compliance must each be
audited separately. A single integrated assessment that reports only one
of them declares less than it appears to.

\textbf{Certification state.} The obligations developed in this section are reported as a vector
\[ \mathrm{Cert}=(\mathsf{Typed},\ \mathsf{Balanced},\ \mathsf{Conserved},\ \mathsf{Positive},\ \mathsf{Admissible},\ \mathsf{Safe},\ \mathsf{Adequate\ service},\ \mathsf{Closed}), \]
each entry taking one of three values: established, not established, and not applicable to this
object. The third value is not a failure and the second is not a refutation. The same three values are the
vocabulary of the per-parameter field of the declaration protocol, applied to identification rather than to
proof: a declared parameter is \emph{established} when the observation map pins it, \emph{not established} when the
record leaves it free, and \emph{not applicable} when no identification question is live for it, as for a
normalising convention. Status is a property of the pair (model, record) and never of a parameter alone, so a
value that is unidentifiable for one readout may be identifiable for another, and the field is reported per
readout --- which is the entry-level content of Proposition 39. A classification record
that establishes an accounting identity and nothing else reports \(\mathsf{Safe}\) as not applicable,
not as refuted; a record that establishes every entry except closure is not thereby unsafe. The
predicates are separately certified and generally non-equivalent rather than pairwise independent:
the implications that hold among them are the ones stated in Propositions 1 and 2 and in the
thermodynamic clause of Section 3.2, and an unlisted implication is not available by inference. The
proof obligation attached to each entry, and the certificate vectors of the indicators classified in
Section 6.5, are tabulated in the supplementary material.


\subsubsection{3.2 The typed safety set}\label{the-typed-safety-set}

The typed safety set at time \(t\) is
\[\mathcal{K}(t) = \{ x \ge 0 : \underline{B}(t) \le Cx \le \overline{B}(t) \},\]
and the non-compensatory assessment reads
\[\mathcal{V}_{\mathrm{typed}} = \{ x(\cdot) : x(t) \in \mathcal{K}(t)\ \forall t \in [0,T] \}.\]
This is a conjunctive criterion: all moiety barriers must be satisfied
simultaneously, and no weighted aggregate \(\sum_m w_m S_m\) is used as
the decision criterion. The force of the conjunction is the subject of
Section 10.1.

\subsubsection{3.3 The flux-reconstruction
identity}\label{the-flux-reconstruction-identity}

\textbf{Lemma 3 (Flux reconstruction under a typed balance law).}
\emph{Let \(x : [0,T] \to \mathbb{R}_+^n\) be absolutely continuous with
\(\dot x(t) = S_{\mathcal{T}} v(t) + b(t)\) almost everywhere,
\(v \in L^1([0,T]; \mathbb{R}_+^J)\),
\(b \in L^1([0,T]; \mathbb{R}^n)\), and \(S = Cx\) with
\(C \in \mathbb{R}^{M \times n}\) the moiety-composition matrix. Then}
\[S(t) = S(0) + \int_0^t \bigl( C S_{\mathcal{T}} v(\tau) + Cb(\tau) \bigr) d\tau,\]
\emph{and for continuous barriers \(\underline{B}, \overline{B}\) the
trajectory satisfies
\(\underline{B}_m(t) \le S_m(t) \le \overline{B}_m(t)\) for all
\(t \in [0,T]\) if and only if the corresponding integrated inequalities
hold for all \(t \in [0,T]\).}

\emph{Proof.} Integrate
\(\dot S = C\dot x = C(S_{\mathcal{T}} v + b) = C S_{\mathcal{T}} v + Cb\)
from \(0\) to \(t\); this is valid because \(x\) is absolutely
continuous and \(C\) linear. The barrier equivalence follows because
\(S_m\) is continuous (as \(x\) is absolutely continuous) and the
barriers are continuous, so a pointwise inequality violation is an
integrated-inequality violation at the same time. \ensuremath{\square}

Two cases are distinguished. \textbf{Prescribed or observed fluxes:} if
\(v(t)\) and \(b(t)\) are known and integrable, stock balances are
reconstructed by integration without solving the internal constitutive
dynamics --- the computationally simple auditing case.
\textbf{Endogenous fluxes:} if \(v(t) = v(x(t), u(t), d(t))\), flux-only
auditing is not generally possible without solving, estimating, or
bounding the coupled system; the identity still holds, but the integral
cannot be evaluated without determining \(v(t)\). In both cases the
barriers are declared, not computed: the flux data do not compute
ecological threshold values, aquifer collapse thresholds, or
concentration boundaries --- the theorem establishes trajectory
compliance with declared barriers, not derivation of the barriers
themselves.

With \(h_m(t) \ge 0\) the net outflow of moiety \(m\) across the system
boundary and \(r_m(t) \ge 0\) the net inflow, the sign convention reads
\[S_m(t) = S_m(0) + \int_0^t \bigl( r_m(\tau) - h_m(\tau) \bigr) d\tau,\]
and the barrier-safety (non-depletion) condition is
\[S_m(0) + \int_0^t \bigl( r_m(\tau) - h_m(\tau) \bigr) d\tau \ge \underline{B}_m(t).\]

\subsubsection{3.4 The conservation-law
reduction}\label{the-conservation-law-reduction}

\textbf{Proposition 4 (Conservation-law reduction).} \emph{If
\(\ell^\top S_{\mathcal{T}} = 0\) for some vector
\(\ell \in \mathbb{R}^n\), then}
\[\ell^\top x(t) = \ell^\top x(0) + \int_0^t \ell^\top b(\tau) d\tau;\]
\emph{in a closed system (\(b = 0\)), \(\ell^\top x\) is invariant.}

\emph{Proof.} This is Proposition 1 integrated:
\(\frac{d}{dt}(\ell^\top x) = \ell^\top b\); integrate. \ensuremath{\square}

A conserved moiety is a row \(c_m^\top\) of \(C\) with
\(c_m^\top S_{\mathcal{T}} = 0\); then \(S_m = c_m^\top x\) satisfies
\(\dot S_m = c_m^\top b\). In a closed system \(S_m\) is constant, and
in an open system \(S_m\) changes only through boundary flows.
Conservation does not imply barrier safety: if a closed conserved moiety
has fixed total stock \(S_m(t) = S_m(0)\), a time-varying barrier can
become infeasible solely because the barrier moves. The three objects
are distinct: the conservation invariant (\(S_m\) constant), the barrier
tube (\(\underline{B}_m \le S_m \le \overline{B}_m\)), and the
intersection of the invariant manifold with the barrier tube. A
conservation law alone does not imply that the trajectory satisfies the
barrier. The precise sense is linear: the set
\(\{x \ge 0 : \ell^\top x = \ell^\top x(0),\ \underline{B} \le Cx \le \overline{B}\}\)
may be empty even though each of the three objects is separately
well-defined. Barrier--conservation compatibility is a linear
feasibility programme, not a slogan.

\textbf{Definition 23 (Bounded residual).} Conservation consistency is stated with a declared residual
budget. For every left-null vector \(\ell\) of \(S_{\mathcal{T}}\), with \(C\) the moiety-composition matrix of Section 2.1 (Lemma 3),
\[ \Bigl|\int_0^T\ell^{\top}C\,d_x(t)\,dt\Bigr|\ \le\ \epsilon_\ell(T), \qquad \epsilon_\ell\ \text{declared with the model}. \]
Proposition 1 then reads \(\ell^{\top}Cx(T)=\ell^{\top}Cx(0)+\int_0^T\ell^{\top}C\,b(t)\,dt\) with
the disturbance term carried rather than absorbed, and \(\epsilon_\ell\) is the quantity
material-flow reconciliation books as statistical discrepancy. An envelope computed without
\(\epsilon_\ell\) supports no barrier certificate once the residual-inflated envelope leaves the
declared barriers, and where \(\epsilon_\ell\) is undeclared the conservation predicate is reported
as not established rather than assumed at zero. Where no budget is declared for \(d_x\), the
residual is not established at zero and is not bounded: the predicate carries the gap rather than
closing it.


\subsubsection{3.5 The flux-bounding envelope
theorem}\label{the-flux-bounding-envelope-theorem}

\textbf{Theorem 5 (Flux-bounding envelopes).} \emph{Assume:}

\emph{(H1) The primitive fluxes and boundary transfers satisfy
componentwise bounds \(v(t) \in [\underline{v}(t), \overline{v}(t)]\)
and \(b(t) \in [\underline{b}(t), \overline{b}(t)]\) for all
\(t \in [0,T]\).}

\emph{For any matrix \(\mathsf{A}\) write
\(\mathsf{A}^{+} = \max\{\mathsf{A}, 0\}\) and
\(\mathsf{A}^{-} = \max\{-\mathsf{A}, 0\}\) entrywise, and define for
each moiety \(m\) the envelope integrands}
\[\varphi_m(\tau) = (C S_{\mathcal{T}})_{m}^{+} \underline{v}(\tau) - (C S_{\mathcal{T}})_{m}^{-} \overline{v}(\tau) + C_{m}^{+} \underline{b}(\tau) - C_{m}^{-} \overline{b}(\tau),\]
\[\psi_m(\tau) = (C S_{\mathcal{T}})_{m}^{+} \overline{v}(\tau) - (C S_{\mathcal{T}})_{m}^{-} \underline{v}(\tau) + C_{m}^{+} \overline{b}(\tau) - C_{m}^{-} \underline{b}(\tau),\]
\emph{and the envelopes}
\[\underline{S}_m(t) = S_m(0) + \int_0^t \varphi_m(\tau)\, d\tau, \qquad \overline{S}_m(t) = S_m(0) + \int_0^t \psi_m(\tau)\, d\tau.\]
\emph{Then \(S_m(t) \in [\underline{S}_m(t), \overline{S}_m(t)]\) for
all \(t \in [0,T]\) and all \(m\).}

\emph{Proof.} By Lemma 3, \(\dot S_m = (C S_{\mathcal{T}} v + Cb)_m\).
For each row \(m\) and time \(\tau\), the bilinear form
\((C S_{\mathcal{T}})_m v\) is linear in \(v\) with coefficient vector
\((C S_{\mathcal{T}})_m\), whose positive and negative parts give, over
the box \([\underline{v}(\tau), \overline{v}(\tau)]\), the pointwise
bounds
\[(C S_{\mathcal{T}})_{m}^{+} \underline{v}(\tau) - (C S_{\mathcal{T}})_{m}^{-} \overline{v}(\tau) \;\le\; (C S_{\mathcal{T}})_m v(\tau) \;\le\; (C S_{\mathcal{T}})_{m}^{+} \overline{v}(\tau) - (C S_{\mathcal{T}})_{m}^{-} \underline{v}(\tau),\]
attained at the extreme points of the box; the same argument applies to
\(C_m b(\tau)\), and adding gives
\(\varphi_m(\tau) \le \dot S_m(\tau) \le \psi_m(\tau)\). Integrating
over \([0,t]\) yields the stated envelope. \ensuremath{\square}

\textbf{Corollary (Flux-derived barrier certificate).} \emph{If
\(\underline{S}_m(t) \ge \underline{B}_m(t)\) and
\(\overline{S}_m(t) \le \overline{B}_m(t)\) for all \(t \in [0,T]\) and
all \(m\), then every trajectory compatible with the flux bounds is
barrier-safe on \([0,T]\).}

\emph{Proof.} By Theorem 5 every such trajectory satisfies
\(\underline{S}_m(t) \le S_m(t) \le \overline{S}_m(t)\); the certificate
conditions sandwich \(S_m\) between the barriers. \ensuremath{\square}

Two qualifications are part of the theorem. The bounds are conservative:
they hold for \emph{all} flux selections in the declared boxes,
including selections that are not jointly realizable by the coupled
dynamics, so the certificate may fail when a trajectory with jointly
realizable fluxes would pass; attainability requires solving or bounding
the coupled system. When the fluxes are state-dependent, \(v = v(x)\),
the declared box must additionally be forward-invariant under the
coupled dynamics for the envelopes to bound the system's reachable set;
without that condition the corollary is a certificate for
flux-admissible paths only, not for the ODE's trajectories. And the
envelope is an interval computation on the flux data, not a forecast: it
says nothing about what the fluxes will be, only about what every
admissible flux path implies for the stock. Stoichiometric and
donor-limit constraints make the jointly admissible flux selections a
polytope rather than a box; the tight certificate is the linear
programme over that polytope, and the box envelope above is its auditing
relaxation --- the box is what is audited, the polytope what is
realizable. The envelope is the interval-arithmetic counterpart, at the
level of declared flux bounds, of the data-reconciliation practice of
material flow analysis (Brunner and Rechberger, 2004).

\textbf{Worked envelope on the closed block.} On (2) with declared boxes
\(N \in [0, K]\) and \(E \in [0, E_{\max}]\), the mass row of the
incidence gives
\(\dot M = -qEN - C^{A,\mathrm{lim}} \in [-(qE_{\max}K + C^A),\, 0]\),
hence \(M(t) \in [M(0) - (qE_{\max}K + C^A)t,\ M(0)]\). The conservatism
is visible in the extremes: the maximal extraction \(qE_{\max}K\) is
realizable only at \(N = K\), where regeneration vanishes, and maximal
recharge coincides with minimal extraction --- box extremes the coupled
dynamics cannot realize jointly.

\subsubsection{3.6 Finite exhaustion under uniform drift, and its
failure
mode}\label{finite-exhaustion-under-uniform-drift-and-its-failure-mode}

\textbf{Proposition 6 (Finite exhaustion under uniform negative drift).}
\emph{Assume:}

\emph{(H1) \(S\) is absolutely continuous with
\(\dot S(t) \le -\varepsilon < 0\) whenever \(S(t) > B\), where \(B\) is
a constant lower barrier and \(\varepsilon > 0\) a uniform drift bound.}

\emph{(H2) \(S(0) > B\).}

\emph{Then the first hitting time satisfies}
\[\tau_B = \inf\{ t \ge 0 : S(t) \le B \} \le \frac{S(0) - B}{\varepsilon}.\]

\emph{Proof.} While \(S(t) > B\), the drift bound integrates to
\(S(t) \le S(0) - \varepsilon t\); the right-hand side reaches \(B\) at
\(t = (S(0) - B)/\varepsilon\), so by continuity of \(S\) the crossing
occurs no later. \ensuremath{\square}

The uniform-margin assumption is essential, and its absence is the
classical failure mode.

\textbf{Counterexample (proportional extraction).} \emph{For
donor-controlled proportional extraction \(\dot S = -kS\) with \(k > 0\)
and \(S(0) > 0\), the stock satisfies \(S(t) = S(0)e^{-kt} > 0\) for
every finite \(t\): the stock approaches zero asymptotically and is
never exhausted in finite time, \(\tau_0 = \infty\). The time to a
positive barrier \(B > 0\) is \(\tau_B = k^{-1}\log(S(0)/B)\), finite
for \(B > 0\) but diverging as \(B \to 0\).}

In particular, the claim that positive extraction implies finite
exhaustion is false, and every exhaustion statement must name its
referent. Internal transfers do not exhaust a total moiety: in a closed
system, internal conversion redistributes a conserved moiety but does
not exhaust it. Exhaustion of a compartment requires a boundary outflow,
an irreversible conversion into an uncounted or unavailable form,
destruction of the relevant function, or a barrier defined on a
particular compartment rather than on the total moiety. Whether
``exhaustion'' refers to the total conserved moiety, a compartment
stock, an accessible or functional stock, a stock above a lower barrier,
or an economically recoverable reserve --- these are different
quantities, and the theorem must specify which (Section 6.4). The
counterexample and this taxonomy are the content of the following named
statement, which the rest of the article uses as its boundary
discipline.

\textbf{Proposition (Depletion is compartmental).} \emph{On a closed
typed ledger, the total mass \(\mathbf{1}^\top x\) of each conserved
moiety is invariant along trajectories; consequently every finite
hitting time of a zero readout is the hitting time of a compartment or
of a barrier on a readout, never of total mass, and ``exhaustion of the
natural block'' (Theorems 7 and 14) is transfer across the block
boundary into the product, waste, and inert compartments.}

The boundary discipline this proposition fixes is the
ecological-economics one of Daly (1990): depletion is not loss of
matter, it is loss of access to matter in the form and place that
supports the service in question. The remainder of the article keeps
that distinction in force.

\begin{center}\rule{0.5\linewidth}{0.5pt}\end{center}

\textbf{Theorem 24 (Critical-margin budget).} Let the declared barrier margins be affine,
\(m(x)=Gx+a\ge0\), and let the delivered service rate be \(y=c^{\top}v\) with \(c\ge0\). Suppose some
\(\lambda\ge0\) satisfies \(\lambda^{\top}GS_T+c^{\top}\le0\) componentwise, with \(\lambda_j\)
carrying the units that convert margin \(j\) into cumulative service. Then \(V=\lambda^{\top}m(x)\)
obeys \(\dot V\le-y+\lambda^{\top}Gb\) along every admissible trajectory, and for every trajectory
that remains inside the declared barriers
\[ \int_0^Ty(t)\,dt\ \le\ V(x(0))+\int_0^T\lambda^{\top}Gb(t)\,dt . \]
If additionally \(y\ge y_{\mathrm{req}}\) and \(\lambda^{\top}Gb\le\beta<y_{\mathrm{req}}\), then
\(T\le V(x(0))/(y_{\mathrm{req}}-\beta)\).

\emph{Proof.} \(\dot V=\lambda^{\top}G(S_Tv+b)=(\lambda^{\top}GS_T)v+\lambda^{\top}Gb\le -c^{\top}v+\lambda^{\top}Gb\) since \(v\ge0\); integrate and use \(V(T)\ge0\), which holds because
\(m\ge0\) and \(\lambda\ge0\). \ensuremath{\square}

The search for \(\lambda\) is a linear programme, and its infeasibility is no evidence of safety: it
certifies only that this multiplier family does not close. The theorem bounds how long adequate
service can coexist with componentwise safety; it is not a scalar certificate, and \(V>0\)
establishes nothing about the sign of any individual margin. The uniform-drift bound above is its
single-margin instance, and the two-sided demand version replaces no statement of Section 5.



\subsubsection{3.7 Composition of ledgers and the calculus of certificates}\label{composition-of-ledgers-and-the-calculus-of-certificates}

Everything above is stated for one ledger. Applications assemble ledgers --- a resource module, a
process module, a territory --- and an assembled object inherits its obligations only if the
assembly is itself an object of the theory. This section supplies that missing definition, and with it
the two results the rest of the article needs: closure is not compositional, and the margin an
interface costs is computable.

\textbf{Definition 34 (Composition of typed ledgers).} Let \(L_1 = (x^1, v^1, S_1, B_1, C_1)\) and
\(L_2 = (x^2, v^2, S_2, B_2, C_2)\) be ledgers over disjoint compartment sets, and let an \emph{interface}
\(\Theta = (J, \Phi, \bar v)\) consist of a finite set \(J\) of declared identifications pairing a
compartment of \(L_1\) with a compartment of \(L_2\), a finite set \(\Phi\) of declared exchange fluxes
with boxes \(0 \le f_\varphi \le \bar v_\varphi\), and, for each \(\varphi\), the two endpoints it drains
and fills. The composition \(L_1 \oplus_\Theta L_2\) is the ledger on the quotient compartment set
(disjoint union modulo \(J\)) whose incidence is the block-diagonal \(\mathrm{diag}(S_1, S_2)\) extended
by the identification rows and by the endpoint columns of \(\Phi\), with readouts carried over unchanged
on each side. The composition is \emph{admissible} when three conditions hold on the declared data:

1. like with like: every identification pairs compartments of the same material type and the same
   unit, and the \emph{interface defect}
   \(d_\Theta = \sum_{(a,b) \in J} \bigl[\,\mathbf 1\{\mathrm{type}(a) \ne \mathrm{type}(b)\} + \mathbf 1\{\mathrm{unit}(a) \ne \mathrm{unit}(b)\}\bigr]\) vanishes;
2. joint boxes declared: no capacity constrains fluxes of both ledgers unless that sharing is one of
   the declared objects, in which case the shared box is a constraint of the composition and not of
   either part;
3. antisymmetry: each exchange enters one ledger as an outflow and the other as an inflow of the same
   magnitude, \(\textstyle\sum_{i} B_i\, e_\varphi = 0\) on the identified compartments.

The admissible compositions are exactly those for which the flux polytope of the whole is the fibre
product of the parts' polytopes over the interface coordinates. The defect and the antisymmetry
residual are computed from the declaration, not assumed away; a composition with \(d_\Theta > 0\) is not
a ledger, and no statement below applies to it.

\textbf{Definition 35 (Interface price).} Suppose each part is certified as in Theorem 24, with multiplier
\(\lambda_i \ge 0\), margins \(m_i(x^i) = G_i x^i + a_i\), and readout \(y_i = c_i^{\top} v^i\). For a
declared exchange \(\varphi\) draining compartment \(c\) on the paying side, its \emph{price} is
\[
\pi_\varphi = \bigl(\lambda_1^{\top} G_1\bigr)_c - \bigl(\lambda_2^{\top} G_2\bigr)_c ,
\]
evaluated at the identified compartments. A price is margin consumed (positive) or margin released (negative) per unit of interface flux per
unit of time, \emph{in the declared orientation}: what the composition must budget for is the magnitude
\(|\pi_\varphi|\,\bar v_\varphi T\), and the sign only decides which side pays. Reversing an exchange's
declared direction leaves that magnitude unchanged and moves the charge to the partner ledger, so no claim
of a free interface can rest on the sign of \(\pi_\varphi\) alone. The additivity of these charges is
asserted for a two-part composition; a composition of three or more parts sharing a compartment needs a
declared global interface, because a flux the whole admits need not decompose into part-wise admissible
exchanges, and two parts each delivering into a third attain service their summed budgets do not cover.

The two ways of reading the display above are different constructions, and the distinction is
load-bearing, so both are given. \emph{Rows.} Take the disjoint union, keep every declared coordinate, and add the
equalities \(x_a^1 = x_b^2\) for each \((a,b) \in J\) as constraint rows; capacities stay attached to the flux
columns. \emph{Quotient.} Let \(R\) be the class-incidence matrix of the equivalence relation generated by \(J\), with
\(R_{Ca} = 1\) exactly when \(a \in C\), and set \(x_C = R x_0\), \(S_C = R\,[\,\mathrm{diag}(S_1, S_2)\ \ E\,]\),
\(B_C = R B_0\), where an exchange \(\varphi\) draining \(a\) and filling \(b\) enters as the column
\(E_\varphi = -e_a + e_b\), and the classes are read off by connected components. The row form is the one this
article uses, because the certificates are stated on the declared coordinates and because the row form stays
meaningful when the interface defect of condition 1 fails; the quotient form is meaningful only when types and
units agree \emph{throughout} each class rather than pair by pair, and then the two coincide on the feasible
flux set. Condition 1's \(\mathrm{type}(\cdot)\) is \(\mathrm{ty}\) of Definition 47, which is what makes the
agreement check decidable from the declaration. Neither construction carries certificates across on its own:
the descent of a part's conservation law to the composition is exactly the agreement condition of
Proposition 36, and it can fail.

\textbf{Proposition 36 (Conservation composes; closure does not).} Let \(\Theta\) be admissible. The conserved quantities of the composition are exactly the pairs of part-wise left-null vectors that
agree on every identified compartment: such a pair lifts to a conservation of the whole, and every
conservation of the whole arises that way. The lift is injective and need not be onto, so a conserved total
of a part can be destroyed by an identification the typing permits. Closure of
the cycle at a demanded rate does not compose: the closure capacity of the composition can be strictly
below the minimum of the parts' closure capacities. Instance: two cycles, each with demand \(1\) unit
per year and return capacity exactly \(1\) unit per year, each therefore closing tightly at \(\Lambda^{*} = 1\); when the two return fluxes draw on a single declared capacity of \(1.5\) units per year, the composition has \(\Lambda^{*} = 0.75\) and closure deficit \(0.25\). Every trajectory of the
composition meeting the demanded rate then has non-stationary stock, so by Definition 22 the shortfall
appears as support drawdown together with sink accumulation, at a rate the linear programme fixes.

\emph{Proof.} The identification rows are the only further constraints, so a covector is left-null on the
quotient precisely when its pullback is left-null on each part and constant on each identification class;
the pullback is therefore injective, with image the compatible subspace, and it is not onto. Witness: two
compartments in series in each part carry two moieties, and after the admissible identification of the two
middle compartments the quotient carries one, so the moiety of the first part is gone. Which subspace
survives is a rank computation on the declared data --- with \(D_J\) the matrix of merged-compartment rows,
the surviving conserved covectors are \((U_1 \oplus U_2) \cap \ker D_J^{\top}\) --- the same computation
that decides closure. For the instance, the composition's programme is \(\max \mu\) subject to \(q = \mu\), \(s = \mu\), \(q + s \le 1.5\)
with \(q, s \ge 0\), whose optimum is \(\mu = 0.75\); the parts' programmes are \(q \le 1\) and \(s \le 1\), both with optimum \(1\). \ensuremath{\square}

\textbf{Proposition 37 (Margin needed, and the admissible exchange rate).} Let \(\Theta\) be admissible and
let each part satisfy Theorem 24 with budget \(B_i = V_i(x^i(0)) + \int_0^T \lambda_i^{\top} G_i b_i\,dt\). Then the composition satisfies the same bound for the summed service \(y_1 + y_2\) with budget
\[
B_1 + B_2 + \sum_{\varphi \in \Phi} |\pi_\varphi|\, \bar v_\varphi\, T .
\]
The sum is the \emph{margin needed} to absorb the interface; the parts' budgets add without remainder
exactly when every declared interface either releases margin on the side that carries the budget or is
slack at its declared bound. Under the outflow-negative convention of Definition 35 the sharper charge
\((-\pi_\varphi)^{+}\bar v_\varphi T\) applies and the interface is free when \(\pi_\varphi \ge 0\).
With a single exchange of box \([0, L]\), non-zero price \(\pi\) and horizon \(T\), this multiplier pair
certifies the composition whenever
\[
L \;\le\; L_{\max} := \frac{B_1 + B_2}{|\pi|\, T},
\]
and at \(L = L_{\max}\) the interface charge equals the combined budget by construction. The converse is not
claimed: another multiplier pair can certify a rate this one does not, and the exact threshold for the
declared horizons is the value of the joint programme over the fibre product of the two declared polytopes
--- not the scalars \(\lambda_i\), \(G_i\), \(\pi\) quoted here, which do not determine it on their own.
Instance: \(G_1 = G_2 = 1\), \(b_1 = b_2 = 0\), \(\lambda_1 = 1\), \(\lambda_2 = 2\), so \(\pi = -1\) and the
exchange drains the ledger holding the cheaper margin; with \(B_1 + B_2 = 9\) units of margin and \(T = 30\)
yr, \(L_{\max} = 0.30\) unit per year --- an exchange at \(0.2\) stays inside the combined budget (\(6 \le 9\))
and one at \(0.5\) does not (\(15 > 9\)). Interchanging the two multipliers leaves the magnitude unchanged and
moves the charge to the other ledger.

\emph{Proof.} Take \(V = V_1 + V_2\). Along the composition, \(\dot V \le -(y_1 + y_2) + \lambda_1^{\top} G_1 b_1 + \lambda_2^{\top} G_2 b_2 - \sum_\varphi \pi_\varphi f_\varphi\): the exchange terms are the only ones that do not
decouple, because a flux that leaves margin in one part enters it in the other at the partner's
multiplier. Each interface term is then bounded over its box by its worst case,
\(\max_{0 \le f \le \bar v_\varphi}(-\pi_\varphi f) = (-\pi_\varphi)^{+}\bar v_\varphi \le |\pi_\varphi|\bar v_\varphi\), and integrating with \(V(T) \ge 0\) gives the stated budget. The
one-dimensional case is the equality \(L|\pi|T = B_1 + B_2\) solved for \(L\). \ensuremath{\square}

\textbf{Definition 40 (Compensation as a decidable predicate).} Fix a partition of the declared
compartments into donor and recipient sides, a demanded service rate, and the joint polytope of the
composition. \emph{Compensation holds at fraction \(\kappa\)} when one linear programme on that polytope is
feasible: maximise the recipient's sink-weighted margin release net of the donor's, subject to the joint
polytope, to the demanded service, and to the donor's budgets being drawn at no more than \(\kappa\) times
what its own service costs; the predicate is true when the optimum covers the demanded drawdown. It is a
feasibility test on the fibre product rather than a price comparison, and it is the exact form of the
question the aggregation cannot answer. Prices do not substitute for it: the shared return capacity that
leaves each of the two cycles of Proposition 36 closing at \(\Lambda^{*} = 1\) individually puts the
composition at \(\Lambda^{*} = 0.75\), and cutting one side's own return capacity to \(10^{-4}\) units per
year leaves that shared-capacity programme well posed and finite while the cut side cannot close at all.
Substitutability is decided by the joint programme, not by whether a capacity price is finite.
The verdict has three branches and they should not be collapsed. Feasibility establishes steady-rate
compensation \emph{under the declared inputs}. Infeasibility establishes failure of that programme, and it comes
with its own witness: a nonnegative combination of the constraint rows of the displayed system reducing to
\(0 \le -1\), the linear-programming alternative read in the direction that certifies failure rather than the
one that certifies safety. Where the declaration is incomplete --- a substitution column missing, a shared
constraint unlisted, a target unnamed --- neither branch is available, and the predicate is reported as not
established, which is a statement about the declaration and not about the ledger. The programme is a
steady-rate object throughout: it says nothing about the trajectory between endpoints, and it establishes no
dynamic safety, no exit-time bound and no corridor invariance; those are the questions of Sections 3.5 to 3.6
and they need their own hypotheses. A positive interface charge does not prohibit compensation either: it
prices it, and with adequate stock, capacity and budget the composition may close --- which is why the charge
enters Definition 34's budgets and never its admissibility conditions.

\textbf{Remark 34 (What the multiplier search can and cannot prove).} The certificates above search over the
conservation multipliers \(\lambda\) alone, because those are what the declared conservation structure
supplies; the capacity and box rows of the declared polytope carry no multiplier, and the gap is visible
in a single-compartment example: with \(\dot m = -y + \sigma\), \(y \in [0, 1]\), \(m(0) = \sigma(0) = 0\)
and \(m(T) \le 1\), a certificate built from \(\lambda\) only returns \(100.0\) units of cumulative service
where the exact value over the horizon is \(1.0\), and pricing the capacity row as well --- the full dual
pair \(\lambda = 0\), \(\rho = 1\) --- returns \(1.0\) exactly. The same refinement applies to the interface
bound of Proposition 37. In the other direction the search is complete for the terminal affine relaxation
of a linear differential inclusion, and on the joint polytope exactness follows from Proposition 37's
joint programme; it is not complete for quadratic certificates, since a valid quadratic margin need not be a sum
of squares --- Motzkin's polynomial is the standard non-negative-not-sum-of-squares obstruction --- so a
search restricted to diagonal multipliers and sums of squares can miss a certificate that exists, and that
failure is no evidence of unsafety.

Two dependencies belong with this reading. The interface price of Definition 35 is a \emph{certificate-relative}
number: with unit service coefficients on both sides, \(\lambda = (2,1)\) and \(\lambda = (1,2)\) are both
feasible part certificates for the same pair of drain ledgers, and they return \(\pi_\varphi = +1\) and
\(\pi_\varphi = -1\) for the same declared exchange. Physical admissibility cannot turn on which certificate was
quoted, and in this article it does not: Definition 34's three conditions involve no multiplier, and
Definition 40's programme decides substitution without reference to either sign. What is certificate-relative
is the charge's size, so it is reported with the certificate that produced it rather than as a property of
the interface. A capacity multiplier is a selection in the same sense: the value of the shared-capacity
programme is a concave function of the capacity, at a kink its subdifferential is an interval --- in the
two-cycle instance of Definition 40, where the capacity equals the sum of the two cycles' own returns, that
interval is \([0, 1]\) in units of one cycle's marginal return --- and its one-sided derivatives at the kink
differ. A single quoted price there is a convention, and it is properly named as a supergradient of the value
function rather than as the shadow value of the capacity.

\textbf{Definition 41 (Control margin).} Let \((\mathcal K_\rho)_{\rho \ge 0}\) be the declared nested family of
corridor sets. The \emph{control margin} of a state is
\[ \mu(x_0) \;=\; \sup\{\, \rho \;:\; x_0 \in \mathcal K_\rho \,\}. \]
Two readings are licensed and no third is: with a normalised margin \(\mu\) is dimensionless, and it is a
min-type quantity in the sense that \(\mu(x_0) \ge \rho\) if and only if \(x_0 \in \mathcal K_\rho\), so what it
answers is whether a declared corridor contains the state. It prices no recovery, it bounds no exit time on
its own, and it inherits every conservatism of the envelope that produced the family.

\textbf{Definition 42 (Affinity, the fourth admissibility predicate).} Conservation, capacity and typing
(Definition 47) are the three predicates a declared flux must pass. Where the declaration also carries a specific Gibbs energy
\(g_j\) for each compartment, a fourth is available: a conversion with stoichiometric row \(\nu\) (products
positive) has \emph{affinity} \(\mathsf{A} = -\nu^{\top} g\), and the conversion is admissible only when
\(\mathsf{A} \cdot J \ge 0\) along its declared rate \(J\) --- equivalently, only when it does not raise the
total Gibbs energy in the direction it is booked. A flux with \(\mathsf{A} \cdot J < 0\) is not a noisy
observation of an allowed process but an impossible one, and the reparation is to re-declare the ledger,
never to re-estimate the flux. Where energies are not declared the predicate is reported as not
established, exactly as the conservation residual of Definition 23 is.

\textbf{Proposition 40 (An empty corridor comes with a named conflict).} For declared bounds \(A v \le b\) and
declarative identities \(C v = d\), the corridor \(\{v : A v \le b,\ C v = d\}\) is empty if and only if there
are \(y \ge 0\) and free \(z\) with
\[
y^{\top} A + z^{\top} C = 0 \qquad \text{and} \qquad y^{\top} b + z^{\top} d \;<\; 0 .
\]
Any such pair is a \emph{corridor-infeasibility witness}, and its support --- the rows with \(y_i > 0\) --- names
the declared floors and ceilings that cannot hold at once. The empty-corridor findings of Section 6.5 are of
this kind: they say not that the system is unstable but that the declaration is jointly unsatisfiable, and
the remedy is to withdraw a bound, not to re-tune a controller.

\emph{Proof.} Append the identities as \(C v \le d\) and \(-C v \le -d\) and apply Gale's theorem of the
alternative (Gale, 1957; Ahuja et al., 1993): its second system is the displayed pair, with \(z\) the
difference of the two multiplier blocks and the strict inequality normalised from \(-1\). \ensuremath{\square}

\textbf{Definition 43 (Circulation time).} On a stationary pattern in which a cycle of \(n\) pools carries the
common rate \(q > 0\) past masses \(m_1, \dots, m_n\), the \emph{circulation time} is
\[ \tau^{\mathrm{circ}} \;=\; \frac{\sum_i m_i}{q} \;=\; \sum_i \frac{m_i}{q} , \]
the mean time a unit of the moiety takes per revolution of that cycle. It is a property of the declared
pattern, not of a trajectory, and it is not any pool's residence time: stationarity forces one common rate
around a simple cycle, so a pool that several cycles share has a residence time per source rather than one
overall, and where several cycles coexist the quantity is a vector indexed by cycle --- a scalar quotation
must name its cycle. Yield inflation is invisible here by construction: a demand met with no stationary pattern behind it
has no circulation time to report, which is the statement of Section 6.6 in reciprocal units.
\textbf{Definition 44 (Governability ratio).} Let \(h\) be the declared review interval of Section 9, \(\tau\) the
declared response delay of the assessment channel, and \(T_{\mathcal{K},m}\) the margin-return time of moiety
\(m\): the time its margin needs to re-enter the declared corridor \(\mathcal K\) under the command in force.
The \emph{governability ratio} is
\[ G_{\mathcal{K},m} \;=\; \frac{T_{\mathcal{K},m}}{\tau + h} . \]
The reading is a comparison of clocks, not a stability theorem. Where \(G_{\mathcal{K},m} > 1\) the corridor is
left and re-entered faster than the periodic assessment can act on it, so review cannot be the instrument
holding that moiety and only the automatic triggers named with the declaration act within the corridor's own
time scale; where \(G_{\mathcal{K},m} \le 1\) the channel acts at least once per excursion. The ratio is
dimensionless and per moiety, so a scalar quotation must name its moiety and its corridor. In the registered
delay family of the companion analysis the same comparison sits between the two thresholds reported there: the
regime change its owner reports near \(\tau \approx 3.7\) yr, and the review interval above which the same
analysis reports restabilisation, about \(6.5\) yr (Section 9). Both are witnesses inside a declared model, not
constants for the class.

\textbf{Lemma 4 (A moving baseline contributes a covariance term).} For an index reported as the difference
\(Y - B\) of an estimate and its own baseline,
\[ \operatorname{Var}(Y - B) \;=\; \operatorname{Var} Y + \operatorname{Var} B - 2\operatorname{Cov}(Y, B) , \]
so a declining reported anomaly is compatible with a stationary process whose baseline alone moved: the part
of the variance manufactured by the baseline is \(\operatorname{Var} B - 2\operatorname{Cov}(Y, B)\), and it is
invisible unless both are declared. This is the covariance form of a warning the index already carries, and
it is why an anomaly is reported together with the construction of its baseline rather than against an
implicit one.
\textbf{Definition 45 (Latest safe intervention time).} Let a critical support pool have margin
\(m_0 = A_0 - A_{\min} > 0\) over its declared barrier and net drawdown \(d_0 > 0\), so \(\dot A = -d\).
Let \(\tau\) be the declared unavoidable delay before drawdown can begin to fall, and let \(\rho > 0\) be the
declared maximum rate of that fall, with the reaching of \(d = 0\) declared achievable. The \emph{latest safe
intervention time} is the largest further postponement that still admits a service-preserving transition:
\[ t_{\mathrm{last}} \;=\; \sup\Bigl\{\, L \;:\; m_0 \;\ge\; \inf_{d(\cdot)} \int_0^{\,\tau+L+d_0/\rho} \!\! d(t)\,dt \;\text{over}\; d \equiv d_0 \;\text{on}\; [0,\tau],\ \dot d \ge -\rho,\ d \ge 0 \Bigr\} . \]
It is a feasibility deadline, not a forecast: the quantifier order is "there exists one response, and every
declared disturbance is then respected", not "every disturbance has some response of its own", and the inner
infimum prices the best admissible response rather than an expected one.

\textbf{Proposition 41 (The deadline has a closed form, and it is not the horizon).} Under the declaration of
Definition 45 the least attainable additional loss is
\[ m_{\mathrm{needed}} \;=\; d_0\tau + \frac{d_0^{2}}{2\rho} , \]
a transition respecting \(A \ge A_{\min}\) exists exactly when \(m_0 \ge m_{\mathrm{needed}}\), and
\[ t_{\mathrm{last}} \;=\; \frac{m_0}{d_0} - \tau - \frac{d_0}{2\rho} \;=\; H^{\mathrm{loc}} - \tau - \frac{d_0}{2\rho} , \]
with \(H^{\mathrm{loc}} = m_0/d_0\) the frozen-rate ratio of Proposition 26, read on the support pool in
the manner of Definition 22. The two objects differ by the
physical cost of delay and of finite deployment speed, so \(t_{\mathrm{last}} < H^{\mathrm{loc}}\) whenever
\(\tau\) and \(\rho\) are finite, and the gap is not a discount applied to the horizon but a different
quantifier over the same declaration. Where \(t_{\mathrm{last}} < 0\) the pool can lie above its barrier,
every trajectory inside the declared corridor, while the service-preserving transition is already
infeasible --- which is the case a horizon cannot express.
Illustration, declared as illustrative and not empirical: \(m_0 = 100\) units, \(d_0 = 10\) units per year,
\(\tau = 2\) yr, \(\rho = 1\) unit per year per year gives \(H^{\mathrm{loc}} = 10\) yr and \(m_{\mathrm{needed}} = 70\),
so ten years of coverage leaves three years to begin: \(t_{\mathrm{last}} = 10 - 2 - 5 = 3\) yr. With
\(m_0 = 25\), \(d_0 = 4\), \(\tau = 3\) yr and \(\rho = 0.5\) the horizon reads \(6.25\) yr while
\(t_{\mathrm{last}} = -0.75\) yr. Replacing \(H^{\mathrm{loc}}\) by the true horizon \(T\) is a step Proposition 26
licenses in direction only: where the depletion law is nondecreasing in the stock, \(T \ge H^{\mathrm{loc}}\)
and the deadline moves later, never earlier, while the two ramp terms are untouched; the closed form above is
then the conservative side of the substituted expression, and no equality is claimed for it.

\emph{Proof.} The integrand \(d\) is non-negative and the constraint set is lower-bounded in slope, so the loss is
minimised by letting \(d\) fall as early and as fast as it is allowed to: any postponement of the ramp adds
area and nothing else. On \([0, \tau]\) the drawdown is held at \(d_0\), contributing \(d_0 \tau\); the ramp from
\(d_0\) to \(0\) at slope \(-\rho\) lasts \(d_0/\rho\) and contributes the triangle \(\tfrac{1}{2} d_0^{2}/\rho\).
Adding gives \(m_{\mathrm{needed}}\); the equivalence is the definition of the supremum; and solving
\(d_0(\tau + L) + d_0^2/(2\rho) = m_0\) for \(L\) gives the stated \(t_{\mathrm{last}}\). \ensuremath{\square}

\textbf{Definition 46 (Supportable-output envelope).} For a declared ledger, barrier family and horizon,
\[ \mathcal Y(T) \;=\; \bigl\{\, y \ge 0 \;:\; \exists\ \text{an admissible trajectory on } [0,T] \ \text{with service at least } y \ \text{throughout and } x(t) \in \mathcal K(t) \ \forall t \bigr\} \]
is the set of service rates supportable for the whole horizon. Three readings and no fourth. The set is
decreasing in \(T\) by restriction, so an envelope number is always a pair (rate, horizon). A single quoted
rate is a selection from the set, and Proposition 29 says which selections can be certified aggregators;
the envelope itself is the report. And at the frozen-rate corner the first empty horizon is the
\(T = m_0/(\underline\delta g_m)\) of Definition 22, so a published "years of supply" measures the declared
drawdown path together with the barrier, never the stock alone: it is an envelope statement, and a
sustained deficit below the declared \(\underline\delta\) moves it without moving the stock.

\textbf{Remark 35 (What a non-displacement gate can say on this apparatus).} An additionality or
non-displacement claim is a comparison of two declarations, not a forecast of a counterfactual world: on this
apparatus the intervention is non-displacing at the declared timescale when the closure deficit of
Definition 22, evaluated on the composition of the project block with the host ledger, is not below the
deficit of the host ledger alone. The gate is then checkable from the declarations, and it fails in the
ordinary way --- a project that re-labels an existing mobilisation as its own leaves the composed deficit
where it was, and a predicate on deficits reports that, where a predicate on attributions would not. What
the gate cannot say is carried by its baseline: the host ledger is declared, not estimated, and where the
two declarations share a compartment the deficit is read on the quotient, so it is not the sum of the two
deficits --- Proposition 36's warning about conserved quantities applies to deficits as well.
These three statements are the calculus the article's own citations had deferred to a companion: an
object of composition, the predicate that survives it, and the number that prices what does not. No
further deferral is made here, and nothing in Sections 4 to 10 depends on a claim about composition
that is not proved above.

\subsubsection{3.8 The type structure the operator's subscript abbreviates}\label{3-8-the-type-structure-the-operator-s-subscript-abbreviates}


\textbf{Definition 47 (Type structure).} Equation (1) writes the incidence operator as \(S_{\mathcal{T}}\), and the
subscript has not been defined. It is not decoration. The sentence of Section 2.1 that entries may be added
within a row "only when their types and units agree", the requirement that a disturbance \(d_x\) be itself typed,
Definition 42's listing of typing as one of the three predicates a declared flux must pass, and Section 6.5's
charge that a score mixes objects which are incommensurable under the typing it invokes --- all of these
presuppose an object that is never given. A \emph{type structure} \(\mathcal{T}\) on a ledger is a single declaration consisting of three
items: a set \(\mathsf{Ty}\) of types; for each compartment \(i\), a type \(\mathrm{ty}_i \in \mathsf{Ty}\) and a
unit \(\mathrm{un}_i\), so that an entry of the state vector carries the pair
\((\mathrm{ty}_i, \mathrm{un}_i)\) rather than a number alone; and a declared set \(\mathsf{Cv}\) of \emph{conversion
coefficients}, each an ordered triple \((\alpha, \beta, c)\) with \(\alpha\) and \(\beta\) distinct types and
\(c > 0\), attached to a named conversion process and read as "one unit of \(\beta\) is obtained from \(c\) units of
\(\alpha\) by that process". Nothing in this list is derived from the stoichiometry; all of it is declared
alongside it.

Admissibility then has content. A sum \(\sum_{i \in I} x_i\) is admissible as one ledger quantity only when
\(\mathrm{ty}_i\) and \(\mathrm{un}_i\) agree for every \(i \in I\); across types there is no sum, only a conversion,
and a conversion appears as a signed pair of entries in the column of the primitive flux that performs it. A
primitive flux is \emph{typed-admissible} when its column is of exactly one of two kinds: a \emph{transfer}, all of whose
nonzero entries are \(\pm 1\) among compartments of one type and one unit; or a \emph{conversion}, whose entries on
the two compartments it joins are \(-c\) and \(+1\) for a declared \((\alpha, \beta, c) \in \mathsf{Cv}\); and no
column is both. Where the declaration draws no type distinction beyond units, every column is a transfer and
the predicate is reported as not applicable rather than as passed. Changing \(\mathsf{Ty}\), \(\mathrm{ty}\) or
\(\mathsf{Cv}\) is a change of ledger and not a change of notation: the admissible flux set \(\mathcal{K}\) of
Definition 21 is defined through \(S_{\mathcal{T}}\), so it moves with the declaration, and every certificate
built on \(\mathcal{K}\) moves with it.

\textbf{Proposition 42 (No conservation law crosses a type class).} \emph{Form the graph on \(\mathsf{Ty}\) in which two
types are joined when \(\mathsf{Cv}\) declares a conversion between them, and call the pullback of a connected
component a type class. Then no admissible column of \(S_{\mathcal{T}}\) has nonzero entries in two distinct
classes, and for a permutation of rows and columns}
\[ S_{\mathcal{T}} \;=\; \bigoplus_{\gamma} S_{\gamma}, \qquad \ker S_{\mathcal{T}}^{\top} \;=\; \bigoplus_{\gamma} \ker S_{\gamma}^{\top} . \]
\emph{Every conservation law of the ledger is therefore carried by a single class; no conserved quantity of the
ledger prices one class against another; and an aggregate formed by adding across classes is not a consequence
of equation (1) but an additional declaration, admitted or refused on Definition 23's terms and propagating
through the certificates exactly as Proposition 36 describes.}

\emph{Proof.} A transfer column meets one type, hence one class. A conversion column meets two types joined by a
declared coefficient, hence one class. So every admissible column has support inside a single class, which is
block-diagonality of \(S_{\mathcal{T}}\) once rows are grouped by class; a vector \(L\) satisfies
\(L^{\top} S_{\mathcal{T}} = 0\) if and only if its restriction to each block does, which is the splitting in the
display. The final clause is read off that splitting: a left-null vector is a tuple of left-null vectors, so
nothing in the kernel relates one block to another. \ensuremath{\square}

\textbf{Remark 36 (What the type structure buys, and what it does not).} Three consequences are worth separating,
because the literature runs them together. \emph{Reporting boundaries.} Since the classes are built from the
declared conversions, drawing a reporting boundary differently --- merging two type names into one report
label, or splitting one --- cannot change \(\ker S_{\mathcal{T}}^{\top}\) unless it declares or withdraws a
conversion. That is the invariance the composition calculus is sometimes asked to supply, with the hypothesis
that makes it true and with the boundary of its scope: redrawing \emph{inside} a class is a different act, decided
by Proposition 36's compatibility condition and not by this one. \emph{Valuation.} The predicate is a
well-posedness condition on declarations, not a thesis about worth. It says that a ledger which has added a
gigajoule to an hour of labour has not declared what it did; it says nothing about whether the two are
comparable in worth. A material-and-energy-value position stated as a substantive claim about worth is
therefore not a corollary of this apparatus, and this article does not make it: what is available here is the
position's checkable content --- Definition 47's predicate, and Proposition 29's aggregator, which does not
compensate a deficit in one compartment with a surplus in another. \emph{Reclassification.} A reserve class
entering or leaving an asset boundary is not a change of \(\mathsf{Cv}\) but a boundary transfer, and it is
priced by Definition 22's closure deficit with its declared elasticity, not by the typing predicate. The
statistical standards recalled in Section 1.2 draw their own asset boundary in economic terms of exactly that
kind, so the exposure is common to the frameworks rather than peculiar to this one, and it is a reason the
horizon readouts of Section 6 are stated with their deficit attached.



\subsection{4. Conservation and Positivity of the Closed
Ledger}\label{conservation-and-positivity-of-the-closed-ledger}

\subsubsection{4.1 The natural-block mass
identity}\label{the-natural-block-mass-identity}

\textbf{Theorem 7 (Natural-block mass identity).} \emph{Let
\(M = N + A^{\mathrm{act}} + A^{\mathrm{geo}} + U\). Along every
trajectory of the closed natural block (2) with optional mining
restored,} \[\dot M = -qEN - C^{A,\mathrm{lim}},\] \emph{i.e.~mass
leaves the natural block exactly at the extraction rate, plus the
donor-limited mining rate; under the institutional-failure
specialization (\(C^A = 0\)), \(\dot M = -qEN\). The identity is stated
for the declared harvest routing \(\alpha = 0\) of Section 2.2; with a
detritus-routed harvest fraction \(\alpha > 0\) the block export is
\((1-\alpha)qEN\) and the identity reads
\(\dot M = -(1-\alpha)qEN - C^{A,\mathrm{lim}}\).}

\emph{Proof.} Sum the four equations of (2), with the mining term
subtracted from \(\dot A^{\mathrm{geo}}\):
\[\dot M = (R - qEN) + (-B + e_{GA} - e_{AG} + \gamma_U U) + (-e_{GA} + e_{AG} - C^{A,\mathrm{lim}}) + (T - \gamma_U U) = R - B + T - qEN - C^{A,\mathrm{lim}},\]
and \(R - B + T = R - (R + T) + T = 0\). The mined fraction routes out
of the four-coordinate natural block; the full-ledger theorems of
Sections 4.2--4.3 record the mining column as an internal transfer
between compartments outside the block --- consistent because the block
boundary, not the ledger boundary, is crossed. \ensuremath{\square}

\subsubsection{4.2 Stoichiometric conservation of the full
ledger}\label{stoichiometric-conservation-of-the-full-ledger}

\textbf{Theorem 8 (Stoichiometric conservation).} \emph{Let
\(X = (N, P, W, I, U, A^{\mathrm{act}}, A^{\mathrm{geo}})\) be the mass
compartments of one resource system and \(S_{\mathcal{T}}\) the
incidence matrix of its flux ledger. One-way transfers are non-negative
and donor-limited; net regeneration is the difference of two such
primitives and is signed when \(N > K\). Under the unit-sum routing
constraints with \(0 \le \alpha \le 1\),}
\[\dot X = S_{\mathcal{T}} F(X), \qquad \frac{d}{dt}\mathbf{1}^\top X = 0.\]

\emph{Proof.} Every primitive is a transfer between two compartments, or
a pair of opposite primitives implementing a two-way exchange; the
corresponding column of \(S_{\mathcal{T}}\) has entries \(+1\) and
\(-1\) in the receiving and donating rows and zeros elsewhere. Routing
tensors are column-stochastic in the destination-indexed convention by
construction: each unit of a split flux sums to one across destinations.
Hence \(\mathbf{1}^\top S_{\mathcal{T}} = 0\) and
\(\mathbf{1}^\top \dot X = \mathbf{1}^\top S_{\mathcal{T}} F = 0\). The
theorem is an exact conservation identity under the routing constraints.
\ensuremath{\square}

The seven-compartment incidence claimed by Theorem 8 --- in the
compartment order of its statement, with the closed block's primitive
fluxes, the pattern of the displayed six-compartment
\(S(\alpha, \rho_P)\) of Theorem 9 with the inert column (no outflow
from the inert compartment) appended and the harvest column split by
\((\alpha, 1-\alpha)\) --- is displayed here. Rows
\((N, P, W, I, U, A^{\mathrm{act}}, A^{\mathrm{geo}})\); columns gross
regeneration, density-dependent return, harvest, uptake, detritus
return, \(e_{GA}\), \(e_{AG}\), mining (to product), product retirement,
inert-bound transfer (waste \(\to\) inert): \[S_{\mathcal{T}} =
\begin{pmatrix}
1 & -1 & -1 & 0 & 0 & 0 & 0 & 0 & 0 & 0 \\
0 & 0 & 1-\alpha & 0 & 0 & 0 & 0 & 1 & -1 & 0 \\
0 & 0 & 0 & 0 & 0 & 0 & 0 & 0 & 1-\rho_P & -1 \\
0 & 0 & 0 & 0 & 0 & 0 & 0 & 0 & 0 & 1 \\
0 & 0 & \alpha & 1 & -1 & 0 & 0 & 0 & \rho_P & 0 \\
-1 & 1 & 0 & -1 & 1 & 1 & -1 & 0 & 0 & 0 \\
0 & 0 & 0 & 0 & 0 & -1 & 1 & -1 & 0 & 0
\end{pmatrix}.\] Every column is a two-compartment transfer under the
unit-sum routing constraints, so \(\mathbf{1}^\top S_{\mathcal{T}} = 0\)
column by column --- Theorem 8's conservation, at sight --- and the
natural-block rows \((N, U, A^{\mathrm{act}}, A^{\mathrm{geo}})\)
reproduce the four-row display of Section 2.2, with the harvest column
now carrying its full routing (\(\alpha\) to \(U\), \(1-\alpha\) to
\(P\)) instead of the block-export convention of the \(\alpha = 0\)
corner. The waste-routed mining variant appends a column with \(-1\) in
the \(A^{\mathrm{geo}}\) row and \(+1\) in the \(W\) row under the same
pattern. The retirement split \((\rho_P, 1-\rho_P)\) and the inert-bound
source (the absorbing stock \(W\)) are declared routing choices of this
displayed instance; the incidence pattern --- every primitive a
two-compartment transfer, column sums zero --- is the theorem's content,
and no classification depends on the declared choices.

\subsubsection{4.3 Conservation of the six-compartment
ledger}\label{conservation-of-the-six-compartment-ledger}

\textbf{Theorem 9 (Six-compartment conservation).} \emph{For the system
of Section 2.3, \(\frac{d}{dt}\mathbf{1}^\top z = 0\); the total mass
\(M_6 = X + U + A + G + P + W\) is constant along every trajectory on
which the classical solution is defined.}

\emph{Proof.} With \(M_6 = \mathbf{1}^\top z\),
\(\dot M_6 = \mathbf{1}^\top S(\alpha,\rho) v = 0\) because each column
of \(S\) sums to zero; term by term: assimilation gives \(g - g = 0\);
mortality gives \(-m + m = 0\); decomposition gives \(-d_U + d_U = 0\);
geological exchange gives \(e_{GA} - e_{GA} = 0\) and
\(-e_{AG} + e_{AG} = 0\); mining gives \(-c_G + c_G = 0\); and the
remaining harvest and retirement terms satisfy
\(-h + \alpha h + (1-\alpha)h = 0\) and
\(\rho_P r_P - r_P + (1-\rho_P)r_P = 0\). \ensuremath{\square}

Two scope notes are part of the theorem. The conservation argument
applies to the expanded typed incidence system when quality grades are
split --- not automatically to an undifferentiated quality-neutral loop.
And open systems are explicit: imports, exports, atmospheric losses, and
cross-boundary transport enter as typed boundary fluxes, giving
\(\dot M_6 = I_{\partial} - O_{\partial}\); writing these flows
explicitly is preferable to preserving a nominal invariant by allowing
an unobserved or finite donor compartment to become negative. Theorems
7--9 are three instances of the conservation lemma of Proposition 1 ---
each is the identity \(\frac{d}{dt}(\ell^\top x) = \ell^\top b\) for a
declared ledger and boundary with \(\ell = \mathbf{1}\) ---
differentiated only by which compartments the declaration includes.

\subsubsection{4.4 Orthant invariance}\label{orthant-invariance}

\textbf{Theorem 10 (Orthant invariance of the closed ledger).} \emph{The
nonnegative orthant in \((N, A^{\mathrm{act}}, A^{\mathrm{geo}}, U)\) is
forward invariant for the closed natural block (2).}

\emph{Proof.} The right-hand side is locally Lipschitz on a
neighbourhood of the closed orthant: each Michaelis--Menten factor
\(s\), \(\sigma\) is \(C^\infty\) on the nonnegative half-line because
the registered regime keeps \(A_0 > 0\) and \(A_{g0} > 0\). Face by
face. On \(A^{\mathrm{geo}} = 0\) one has \(\sigma = 0\), hence
\(e_{GA} = 0\) and
\(\dot A^{\mathrm{geo}} = e_{AG} = \omega_A A^{\mathrm{act}} \ge 0\). On
\(A^{\mathrm{act}} = 0\) one has \(s = 0\), so
\(R = B = T = e_{AG} = 0\) and
\(\dot A^{\mathrm{act}} = e_{GA} + \gamma_U U \ge 0\). On \(N = 0\),
extraction and uptake vanish and \(\dot N = 0\). On \(U = 0\),
\(\dot U = T \ge 0\). Nagumo's inward-pointing criterion (Aubin, 1991)
yields forward invariance of the orthant. \ensuremath{\square}

\textbf{Theorem 11 (Forward invariance of the six-compartment cone).}
\emph{Under the donor boundary assumptions of Section 2.3 (each
primitive flux vanishes when its donor is empty, fluxes continuous in
effort and locally Lipschitz in the state), \(\mathbb{R}^6_+\) is
forward invariant for the six-compartment system.}

\emph{Proof.} Face by face: at \(X = 0\), \(g = m = h = 0\) so
\(\dot X = 0\); at \(U = 0\),
\(\dot U = m + \alpha h + \rho_P r_P \ge 0\); at \(A = 0\),
\(\dot A = -g + d_U + e_{GA} - e_{AG} = d_U + e_{GA} \ge 0\), the two
negative terms vanishing by donor limitation (\(A\) is the donor of both
\(g\) and \(e_{AG}\)); at \(G = 0\), \(\dot G = e_{AG} \ge 0\); at
\(P = 0\), \(\dot P = (1-\alpha)h + c_G \ge 0\); at \(W = 0\),
\(\dot W = (1-\rho_P)r_P \ge 0\). The vector field belongs to the
tangent cone at every boundary point, and the tangent-cone invariance
theorem applies. Conservation and boundary admissibility are separate
obligations, and the finite-donor condition carries a discipline: a
target-relaxation law \(e_{GA} = \omega(A^{\mathrm{eq}} - A)\) does not
satisfy it unless also limited by \(G\); it may be used only with the
source declared an effectively infinite external reservoir, in which
case the system is open rather than closed. \ensuremath{\square}

The classical lineage of these statements is the compartmental-systems
nonnegativity theory (Jacquez and Simon, 1993); the donor-limitation
condition is the exact sufficiency requirement --- algebraic
cancellation alone does not establish invariance.

\subsubsection{4.5 No interior rest at positive
effort}\label{no-interior-rest-at-positive-effort}

\textbf{Theorem 12 (No interior rest at positive effort).}
\emph{Assume:}

\emph{(H1) \(E \equiv E_* > 0\) is constant.}

\emph{Then a rest point of the closed natural block satisfies
\(R + C^{A,\mathrm{lim}} = 0\) after restoring optional mining; with
\(C^A = 0\) this is \(R = 0\), hence \(N = 0\) or \(N = K\) or
\(A^{\mathrm{act}} = 0\). None of these is compatible with \(E_* > 0\)
and \(N_* > 0\): (i) \(N = K\) and \(E_* > 0\) give
\(\dot N = -qE_* K < 0\); (ii) \(A^{\mathrm{act}} = 0\) and
\(A^{\mathrm{geo}} > 0\) give
\(\dot A^{\mathrm{act}} = \omega_A A^{\mathrm{eq,intrinsic}} \sigma > 0\);
(iii) \(N = 0\) forces \(R = T = 0\) and reduces to the extinction
family \(\mathcal{R}_{\mathrm{ext}}\) of Theorem 13. In particular the
working point \((N^*, A^{\mathrm{act}*}) = (89.526, 397.87)\) is not a
rest point at \(E = E^* \approx 2.090\): \(\dot N = 0\) holds there by
construction (\(R^* = qE^*N^* \approx 0.187 > 0\)), and the rest
condition of the proof below fails on the abiotic pair.}

\emph{Proof.} At a rest point, \(\dot U = 0\) forces \(\gamma_U U = T\).
Adding \(\dot A^{\mathrm{act}} + \dot A^{\mathrm{geo}}\) gives
\(-B + \gamma_U U - C^{A,\mathrm{lim}} = 0\); with \(\gamma_U U = T\)
and \(B = R + T\) this is \(R + C^{A,\mathrm{lim}} = 0\). With mining
declared (\(C^A > 0\)) this forces \(R \le 0\); with \(C^A = 0\) it is
\(R = 0\), and from the constitutive law \(R = rN(1 - N/K)s = 0\)
implies \(N = 0\) or \(N = K\) or \(s = 0\) (that is,
\(A^{\mathrm{act}} = 0\)). Cases (i)--(iii) exclude each branch at
positive effort; in the mining case the contradiction is more direct ---
\(\dot N = 0\) with \(E_* > 0\) and \(N_* > 0\) gives
\(R = qE_*N_* > 0\), while the abiotic rest condition gives
\(R = -C^{A,\mathrm{lim}} \le 0\). At the working point \(\dot N = 0\)
by construction while rest would require \(R = 0\); with \(\dot U = 0\)
the abiotic pair would satisfy
\(\dot A^{\mathrm{act}} + \dot A^{\mathrm{geo}} = -R < 0\). \ensuremath{\square}

\subsubsection{4.6 The extinction--geochemical rest
set}\label{the-extinctiongeochemical-rest-set}

\textbf{Theorem 13 (Vanishing-extraction rest set).} \emph{With vanishing extraction (\(E \equiv 0\)), the rest points of the closed natural block (2) are exactly the three sets --- the extinction family \(\mathcal{R}_{\mathrm{ext}}\), the carrying-capacity family \(\mathcal{R}_K\), and the frozen-biomass face \(\mathcal{R}_{\mathrm{frozen}}\); the union symbol \(\mathcal{R}_0 = \mathcal{R}_{\mathrm{ext}} \cup \mathcal{R}_K\) is retained for the two geochemical families:}
\[ \mathcal{R}_{\mathrm{ext}} = \bigl\{ N = 0,\ U = 0,\ A^{\mathrm{act}} = A^{\mathrm{eq,intrinsic}}\,\sigma(A^{\mathrm{geo}}),\ A^{\mathrm{geo}} \ge 0 \bigr\}, \qquad \mathcal{R}_K = \bigl\{ N = K,\ U = \kappa_A K s/\gamma_U,\ A^{\mathrm{act}} = A^{\mathrm{eq,intrinsic}}\,\sigma,\ A^{\mathrm{geo}} \ge 0 \bigr\}, \]
\emph{where in the second family \(s = A^{\mathrm{act}}/(A^{\mathrm{act}} + A_0)\) is evaluated at the solution --- together with the frozen-biomass face \(\mathcal{R}_{\mathrm{frozen}} = \{(N, 0, 0, 0) : N \ge 0\}\), on which \(s = 0\) identically and the biomass is frozen at its initial value. With \(E > 0\) constant, no interior rest point (with \(N_* > 0\)) exists (Theorem 12); the extinction face \(\mathcal{R}_{\mathrm{ext}}\) of this set persists at positive effort, because extraction \(qEN\) vanishes identically at \(N = 0\), and it is a boundary rest rather than an interior one (Section 4.6). If \(A_{g0} = 0\) and \(\sigma \equiv 1\) is imposed for \(A^{\mathrm{geo}} > 0\), the shared active-pool ray of \(\mathcal{R}_{\mathrm{ext}}\) and \(\mathcal{R}_K\) is \(A^{\mathrm{act}} = A^{\mathrm{eq,intrinsic}}\), \(A^{\mathrm{geo}} > 0\) --- the endpoint \(A^{\mathrm{geo}} = 0\) is excluded, because there the donor-limited recharge vanishes and \(\dot A^{\mathrm{act}} = -\omega_A A^{\mathrm{eq,intrinsic}} < 0\). The constitutive laws carry no basal mortality independent of the support factor; adding one (\(\mu_{\mathrm{basal}} N\), stock \(\to\) detritus) collapses the frozen-biomass face and does not touch Theorems 7--12 or 14.}

\emph{Proof.} With \(E \equiv 0\), set the four derivatives to zero.
From \(\dot A^{\mathrm{geo}} = -e_{GA} + e_{AG} = 0\):
\(e_{GA} = e_{AG}\),
i.e.~\(\omega_A A^{\mathrm{eq,intrinsic}} \sigma = \omega_A A^{\mathrm{act}}\),
so \(A^{\mathrm{act}} = A^{\mathrm{eq,intrinsic}} \sigma\) --- the
geological-exchange balance, in which the active pool is pinned to the
donor-scaled intrinsic target. From \(\dot N = R = 0\):
\(rN(1 - N/K)s = 0\). If \(A^{\mathrm{geo}} > 0\) the geo-balance pins
\(A^{\mathrm{act}} > 0\), so \(s > 0\) and \(N = 0\) or \(N = K\). At
the boundary \(A^{\mathrm{geo}} = 0\) the geo-balance forces
\(A^{\mathrm{act}} = 0\) (since \(\sigma(0) = 0\)), hence \(s = 0\) and
\(\dot N = 0\) for every \(N \ge 0\); with \(U = 0\) the remaining
equations vanish identically, so the frozen-biomass face
\(\{(N, 0, 0, 0) : N \ge 0\}\) is a rest set. From
\(\dot U = T - \gamma_U U = 0\):
\(U = T/\gamma_U = \kappa_A N s/\gamma_U\), which vanishes in the
\(N = 0\) branch and is positive in the \(N = K\) branch. From
\(\dot A^{\mathrm{act}} = -B + e_{GA} - e_{AG} + \gamma_U U = -(R + T) + 0 + \gamma_U U\):
this vanishes in both branches, since \(\gamma_U U = T\) and \(R = 0\)
hold there. The two families together with the frozen-biomass face are
exactly the stated rest set. ``Geochemical'' names the mechanism of both
families' active-pool rest: the pool rests at the donor-scaled intrinsic
target; apart from the frozen-biomass face, no rest point exists away
from extinction or carrying capacity. The institutional memory yields
\(E \to E^*\) at \(N = 0\) with extraction vanishing identically ---
consistent with the rest set and not an interior rest. \ensuremath{\square}

\subsubsection{4.7 Extraction
integrability}\label{extraction-integrability}

\textbf{Theorem 14 (Integrable extraction).} \emph{Assume \(E(s) \ge 0\)
along the trajectory (effort is nonnegative; \(N \ge 0\) along classical
solutions is Theorem 10's). Let
\(M = N + A^{\mathrm{act}} + A^{\mathrm{geo}} + U\). Then
\(M(t) = M(0) - \int_0^t qE(s)N(s)\, ds \ge 0\), so}
\[\int_0^\infty qE(s)N(s)\, ds \le M(0) < \infty;\] \emph{in particular
\(qEN \in L^1(0,\infty)\), and no trajectory maintains extraction at the
working value \(qE^*N^* \approx 0.187\) for all time; with mining
restored,
\(\int_0^\infty \bigl( qE(s)N(s) + C^{A,\mathrm{lim}}(s) \bigr) ds \le M(0)\).}

\emph{Proof.} By Theorem 7, \(M(t) = M(0) - \int_0^t qE(s)N(s)\, ds\);
forward invariance (Theorem 10) gives \(M(t) \ge 0\), so the improper
integral is at most \(M(0)\). If \(qEN \equiv qE^*N^*\) for all
\(t \ge 0\), the integral would diverge. \ensuremath{\square}

This is the depletion-horizon semantics of the closed ledger in its
strongest form: the donor budget is finite and extraction is integrable
against it. A constant extraction flux \(c > 0\) --- a comparison flux
only, not a donor-limited primitive the ledger's own discipline admits
as a sustained law --- exhausts the budget in finite time (\(M\) reaches
its lower bound no later than \(M(0)/c\)), while proportional extraction
\(qEN\) need not drive \(M\) to zero in finite time --- the integral
bound of the theorem is the whole statement, and the hitting time of
\(M = 0\) may be infinite. This is the finite-budget fact that Section 9
turns into the non-reduction boundary with the open working system. The
theorem does not select among the vanishing-extraction rests of Theorem
13: integrable extraction is compatible with approach to either the
extinction family or the carrying-capacity--geochemical family, and the
\(L^1\) bound alone decides nothing between them.

\subsubsection{4.8 The conditional hybrid moiety
balance}\label{the-conditional-hybrid-moiety-balance}

\textbf{Conditional Theorem 15 (Hybrid moiety balance).} \emph{Let
\(\chi\) denote the hybrid state and \(\eta \ge 0\) its primitive-flux
vector --- letters local to this statement, chosen so that \(r\) stays
the growth rate of Section 2.2 and \(\nu\) a macro parameter of Section
5.4. Assume:}

\emph{(H1) \(\chi\) is absolutely continuous between locally finite
event times with left and right limits at events.}

\emph{(H2) \(\dot \chi = \mathsf{S}\eta + b\) with \(\eta \ge 0\),
separate reverse columns, and donor-limited negative boundary flows.}

\emph{(H3) \(\mathsf{L}^\top \mathsf{S} = 0\).}

\emph{Then}
\[\mathsf{L}^\top \chi(t) - \mathsf{L}^\top \chi(0) = \int_0^t \mathsf{L}^\top b\, ds + \sum_{t_k \le t} \mathsf{L}^\top \bigl[ \chi(t_k^+) - \chi(t_k^-) \bigr].\]

\emph{Proof.} Integrate the continuous balance between consecutive
events and telescope the left/right state differences. \ensuremath{\square}

The theorem is conditional, and its jump interpretation is part of the
content: an internal-transformation jump requires
\(\mathsf{L}^\top(\chi^+ - \chi^-) = 0\) or a jump incidence
factorization with left-kernel conservation; a boundary-crossing jump is
a boundary impulse and belongs in the boundary term. Two obligations
ride the theorem. The \emph{yield-routing obligation}: if a
transformation is represented with a yield below one for a declared
moiety, the omitted fraction must be routed to another represented
compartment or a declared boundary flow --- otherwise the claimed moiety
balance holds only after silently dropping that moiety from
\(\mathsf{L}\). And the \emph{separation obligation}: this is the hybrid
variant of Proposition 4, retained at its own conditional status; the
two statements are not merged.

\subsubsection{4.9 Cancellation is cheap}\label{cancellation-is-cheap}

Summing the six material equations of a ten-state admissibility template
gives the exact identity
\[\frac{d}{dt}\bigl( \bar X_A + X_J + P + U + A + G \bigr) = 0.\] This
is an algebraic cancellation only: it does not prove forward invariance
of the six material states or physical admissibility of every term. The
ghost-sink check is part of the discipline: the same birth-transfer rate
\(g_B\) enters \(\dot X_J\) and \(\dot A\) with opposite signs, so
material not transferred to juveniles remains in \(A\) --- there is no
unmatched sink in the six-state ledger. The identity is retained
precisely for its discipline: formal cancellation coexists with boundary
failure elsewhere in the same template (its geological exchange is not
donor-limited), and the cancellation by itself establishes nothing about
admissibility. Conservation (Theorems 7--9) and positivity (Theorems
10--11) are proved separately in every well-posed ledger of this
article, exactly because cancellation is cheap and admissibility is not.
The template's remaining negative witnesses --- a variance closure that
is not realizable by a non-negative spatial distribution, and an output
functional without a displayed state equation --- are recorded in the
supplementary material as audited admissibility failures.

\textbf{The closed-ledger portrait.} Theorems 7--14 assemble into a
complete qualitative portrait of the closed orthant: conservation
(Theorems 7--9), positivity (Theorems 10--11), no interior rest at
positive effort (Theorem 12), the two-family vanishing-extraction rest
set with the frozen-biomass face (Theorem 13), and the finite donor
budget (Theorem 14). The portrait is the source object handed to the
interface of Section 9: the closed system's candidate long-time set is
the rest set of Theorem 13, and the budget of Theorem 14 bounds how long
any positive-flux configuration can persist. For industrial-ecology
measurement the message is direct: a ``balanced'' closed ledger is a
finite-budget object, and any sustained extraction against it must
integrate to a quantity no greater than the initial budget.

\begin{center}\rule{0.5\linewidth}{0.5pt}\end{center}

\subsection{5. Service Readouts and the Componentwise
Deficit}\label{service-readouts-and-the-componentwise-deficit}

Services are observations or feasible outputs of the physical state, not
additional conserved mass. Internal physical transfers are not services
merely because they appear in a ledger: a typed readout identifies the
delivered flow, its boundary, and any unit conversion. This distinction
is the accounting counterpart of the ecological-economics point that a
service flow (Ayres' useful-work reading, Daly's throughput-of-services
reading) is not the same object as the mass that delivers it.

\subsubsection{5.1 The service readout and the contemporaneous
balance}\label{the-service-readout-and-the-contemporaneous-balance}

For services indexed by \(i = 1, \ldots, n\), write
\(s_i(t) = \mathcal{O}_i(x(t), u(t), \theta)\), where \(u\) denotes
admissible operating or extraction choices and \(s_i\) and the demand
\(d_i\) share service-specific units. Where delivered services are
selected or converted ledger fluxes, the readout is linear in the
primitives, \[s = \mathcal{O}(x, u, \theta) = Q(\theta)\, v(x, u),\]
with every row of \(Q\) declaring the delivery boundary and the
conversion into one service-specific unit; more general state-dependent
readouts are possible. The contemporaneous component balance is
\[b_i(t) = s_i(t) - d_i(t),\] and \(b_i(t) \ge 0\) means measured supply
meets measured demand for component \(i\) at that instant. It does not
by itself imply that the underlying trajectory is sustainable: a stock
can meet current demand while declining toward a threshold, and a stock
below a desired level can have a positive current balance while
recovering.

\subsubsection{5.2 The state-dependent feasible balance
domain}\label{the-state-dependent-feasible-balance-domain}

\textbf{Definition 1 (Feasible balance domain).} \emph{For an admissible
operating set \(\mathcal{U}(x,t)\) and a declared demand set
\(\mathcal{D}(t)\),}
\[\mathcal{B}(x,t) = \{ \mathcal{O}(x,u,\theta) - d : u \in \mathcal{U}(x,t),\ d \in \mathcal{D}(t) \}.\]

The geometry of the balance domain is state dependent and inherited
partly from the stock--flow model; no unrestricted argument can replace
an application-specific analysis of \(\mathcal{B}(x,t)\). This domain is
the object against which any scalar certificate claim must be checked
(Section 10.1): a weighted sum certifies componentwise nonnegativity on
\(\mathcal{B}(x,t)\) only through an implication proved from the
physical restrictions that define the domain.

\subsubsection{5.3 Support provenance and the directional support
gap}\label{support-provenance-and-the-directional-support-gap}

Current service adequacy and regenerative feasibility are different
claims. Let \(\Gamma_{\mathrm{all}}(x,t) \subseteq \mathbb{R}^n_+\)
contain the service vectors feasible through all pathways admitted by an
application, and
\(\Gamma_{\mathrm{reg}}(x,t) \subseteq \Gamma_{\mathrm{all}}(x,t)\) the
feasible set after imposing the declared regenerative-flow,
system-boundary, material-quality, and exergy or capacity restrictions.
These correspondences are application inputs obtained from a typed
pathway or technology model; the stock ledger alone does not construct
them.

\textbf{Definition 2 (Directional regenerative-support fraction and
gap).} \emph{Assume:}

\emph{(H1) \(0 \in \Gamma_{\mathrm{reg}}(x,t)\).}

\emph{(H2) A nonzero service direction \(\bar s \ge 0\) is chosen.}

\emph{Define}
\[\alpha_{\mathrm{reg}}(\bar s; x, t) = \sup\{ \alpha \in [0,1] : \alpha \bar s \in \Gamma_{\mathrm{reg}}(x,t) \}.\]
\emph{The vector \((1 - \alpha_{\mathrm{reg}})\bar s\) is the
directional support gap, measured in the same service units as
\(\bar s\). A realized service
\(s \in \Gamma_{\mathrm{all}} \setminus \Gamma_{\mathrm{reg}}\) is
support-dependent under that declaration even when \(s \ge d\).}

Attainment requires closedness: if \(\Gamma_{\mathrm{reg}}\) is not
closed the supremum may not be attained, and the gap is relative to a
supremal fraction, not necessarily an achievable boundary service. The
non-interpretation discipline is equally part of the definition: the
statement neither subtracts raw material from service nor proves that a
physical stock is declining; net depletion still requires a negative
stock balance or a trajectory argument. The provenance partition behind
\(\Gamma_{\mathrm{reg}}\) (renewable flow, recovered or recycled
material, imports, non-renewable drawdown) never adds unlike physical
units.

\subsubsection{5.4 The componentwise deficit and the specialization
identity}\label{the-componentwise-deficit-and-the-specialization-identity}

On the unreduced ledger the physical deficit is the diagnostic
\[\Delta^{\mathrm{phys}}(t) = C(t) - \widehat{M}^\top S(t),\] with \(C\)
the operative extraction-law readout and \(\widehat{M}\) the declared
demand-coverage matrix mapping the moiety readout \(S = Cx\) --- the
composition matrix \(C\) of Lemma 3, a different object from the
coverage vector \(C(t)\) despite the shared letter --- into the units of
that coverage vector --- rows indexed by covered services, columns by
moieties, entries the declared stoichiometric coefficients of the
coverage convention (the hat distinguishes the matrix from the scalar
natural-block mass \(M\) of Section 4.1). It does not drive the physical
equations, and it is not equal to \(-\dot N\) unless waste--product
feedback vanishes and the service is identified with regeneration. The
single-resource specialization (the omitted product, waste, and price
parameters \(\mu, \nu, \rho\) of the unreduced ledger set to zero,
together with \(C^A = 0\)) makes that identification, and on that class
--- and only on that class --- the deficit collapses to the
stock-decline rate.

\textbf{Remark 16 (Exact specialization deficit identity).} \emph{On
every trajectory of the specialized system, and of every reduced system
whose stock equation is \(\dot N = R(N,A) - qEN\),}
\[qEN - R(N,A) = -\dot N, \qquad \Lambda(t) := \bigl[ qEN - R \bigr]_+ = \bigl[ -\dot N \bigr]_+.\]

\emph{Proof.} Substitute the stock equation:
\(qEN - R = -(R - qEN) = -\dot N\). \ensuremath{\square}

The collapse is a property of the specialization, not a definition of
liquidation on the unreduced ledger; the general diagnostic remains
\(C - \widehat{M}^\top S\).

\textbf{Decline pressure.} In the registered delay family the
depletion-pressure classification is
\(\Lambda(t) = \max\{0, qE(t)N(t) - R(N(t), A^{\mathrm{act}}(t))\} = \max\{0, -\dot N(t)\}\):
the memory input of the institutional dynamics is a smoothed
stock-decline rate, exactly the positive part of the decline. It is not
a stock-level scarcity measure, not an unmet-consumption measure, and
not an independently observed service deficit. Since
\(qEN - R(N, A^{\mathrm{act}}) = O(N)\) as \(N \to 0\), the raw decline
input vanishes near extinction while the positive baseline source of the
effort law can still sustain commanded effort --- the incremental
decline amplification disappears, but the effort command need not; a
controller intended to respond to low stock irrespective of its current
rate of change requires a separately registered level-dependent channel.

\begin{center}\rule{0.5\linewidth}{0.5pt}\end{center}

\textbf{Proposition 25 (Accumulated liquidation certified by delivered service).} Let the support pool
obey \(\dot A=r-c\), and let the declared production relation require at least \(\alpha\ge0\) units
of support use per unit of delivered service, \(c(t)\ge\alpha Y(t)\). Then
\[ A(0)-A(T)\ \ge\ \alpha\int_0^TY(t)\,dt-\int_0^Tr(t)\,dt . \]
Where the right-hand side is positive, the delivered service certifies that much support liquidation,
whatever the interior allocation of fluxes.

\emph{Proof.} \(A(0)-A(T)=\int_0^T(c-r)\,dt\ge\alpha\int_0^TY\,dt-\int_0^Tr\,dt\). \ensuremath{\square}

The statement is falsifiable against the record rather than self-sealing: a stock change smaller than
the certified lower bound exposes a wrong coefficient, a missing flux, or a measurement inconsistency,
and each of the three is a registered obligation, not a licence to reconcile. The coefficient
\(\alpha\) is a declared relation, not an estimate; the proposition is conditional on it and is not
computed anywhere in this article.


\subsection{6. Depletion Arithmetic}\label{depletion-arithmetic}

The ledger supplies the net active-pool derivative needed to distinguish
gross throughput from net decline and from a model-conditioned threshold
time. The distinction matters because ``time to depletion'' is publicly
used as if all three were one quantity. They are not, and the worked
instances of this section make the differences explicit. Let
\(A_{\min}\) be a declared threshold for the active abiotic pool with
\(A > A_{\min}\).

\subsubsection{6.1 The three quantities}\label{the-three-quantities}

\textbf{Definition 3 (Gross turnover intensity and support coverage).}
\emph{With assimilation \(g(X,A) > 0\), the gross turnover intensity is
\(J_A^{\mathrm{gross}} = g(X,A)/A\) and the gross support-coverage ratio
is \(H_A^{\mathrm{gross}} = (A - A_{\min})/g(X,A)\).}

Neither is a time to depletion. The implication
\(g > 0 \Rightarrow \dot A < 0\) is false in general: at an interior
steady state, \(g\) can be positive while decomposition and geological
transfer balance it exactly, so that \(\dot A = 0\). Gross uptake
measures throughput or dependency; net depletion is a balance property.
This false-implication record is the first rung of the taxonomy and
governs every application below.

\textbf{Definition 4 (Local net-depletion ratio).}
\[H_A^{\mathrm{loc}}(t) = \frac{A(t) - A_{\min}}{\bigl[ -\dot A(t) \bigr]_+},\]
\emph{with the extended-real convention \(H_A^{\mathrm{loc}} = +\infty\)
when \(\dot A \ge 0\) --- correctly reporting no current net decline at
a stationary or replenishing state. The ratio is still not a trajectory
forecast: it freezes the current net rate. If the fluxes change with
\(A\), policy, climate, prices, or other states, the realized threshold
time can differ substantially.}

\textbf{Definition 5 (Scenario-conditioned hitting time).} \emph{For a
fully specified dynamical model, policy or scenario \(\pi\), disturbance
history \(d\), and initial state \(x_0\),}
\[T_A(x_0; \pi, d) = \inf\{ t \ge 0 : A^{\pi,d}(t; x_0) \le A_{\min} \},\]
\emph{with \(T_A = +\infty\) if the threshold is never reached. Under
parameter, observation, and scenario uncertainty the appropriate output
is a distribution or robust interval of \(T_A\), not a single universal
date.}

The three quantities answer different questions and must not share one
depletion-horizon label:

\begin{longtable}[]{@{}
  >{\raggedright\arraybackslash}p{(\columnwidth - 2\tabcolsep) * \real{0.5000}}
  >{\raggedright\arraybackslash}p{(\columnwidth - 2\tabcolsep) * \real{0.5000}}@{}}
\toprule\noalign{}
\begin{minipage}[b]{\linewidth}\raggedright
Quantity
\end{minipage} & \begin{minipage}[b]{\linewidth}\raggedright
Question answered
\end{minipage} \\
\midrule\noalign{}
\endhead
\bottomrule\noalign{}
\endlastfoot
\(J_A^{\mathrm{gross}}\), \(H_A^{\mathrm{gross}}\) & How strongly does
the system depend on, or turn over, the pool at the current gross
rate? \\
\(H_A^{\mathrm{loc}}\) & If the current net decline were frozen, what is
the local stock-to-rate ratio? \\
\(T_A\) & Under a stated model, policy, and disturbance scenario, when
is the threshold first reached? \\
\end{longtable}

\subsubsection{6.2 Uniform-drift bounds}\label{uniform-drift-bounds}

\textbf{Proposition 17 (Local threshold-horizon bracket).}
\emph{Assume:}

\emph{(H1) \(A : [0,T] \to \mathbb{R}\) is absolutely continuous with
\(A(0) > A_{\min}\).}

\emph{(H2) Constants \(v_0 > 0\) and \(0 < \varepsilon < 1\) are given;
set \(H_0 = (A(0) - A_{\min})/v_0\).}

\emph{(H3) \(T \ge H_0/(1-\varepsilon)\).}

\emph{(H4) \((1-\varepsilon)v_0 \le -\dot A(t) \le (1+\varepsilon)v_0\)
for almost every \(t\) while \(A\) stays above \(A_{\min}\).}

\emph{Then a first crossing time \(H\) exists no later than
\(H_0/(1-\varepsilon)\), and}
\[\frac{H_0}{1+\varepsilon} \le H \le \frac{H_0}{1-\varepsilon}, \qquad |H - H_0| \le \frac{\varepsilon}{1-\varepsilon}H_0.\]

\emph{Proof.} If no crossing occurs before
\(t_* = H_0/(1-\varepsilon)\), absolute continuity gives
\(A(t_*) \le A(0) - (1-\varepsilon)v_0 t_* = A_{\min}\), a
contradiction; hence \(H \le t_*\). Integrating both rate bounds over
\([0,H]\) and using \(A(0) - A(H) = v_0 H_0\) gives the two-sided
bracket: from the upper rate bound,
\(v_0 H_0 \le (1+\varepsilon) v_0 H\), and from the lower rate bound,
\(v_0 H_0 \ge (1-\varepsilon) v_0 H\). \ensuremath{\square}

This is a local diagnostic only: it fails when depletion reverses, the
rate approaches zero, or feedback moves the trajectory outside the
declared rate bounds. Its companion is the one-sided exhaustion
proposition of Section 3.6, whose counterexample --- proportional
extraction never exhausts in finite time --- shows that the uniform
margin \(\varepsilon > 0\) is load-bearing in both directions. The
bracket bounds the frozen-rate ratio's error under declared rate bounds,
and nothing more.

\textbf{When the clocks coincide.} Under the rate bracket the
frozen-rate ratio and the true hitting time agree within the bracket:
\(-\dot A \in [(1-\varepsilon)v_0, (1+\varepsilon)v_0]\) gives
\(H_A^{\mathrm{loc}} = (A(0) - A_{\min})/[-\dot A]_+ \in [H_0/(1+\varepsilon), H_0/(1-\varepsilon)]\),
hence \(|H - H_A^{\mathrm{loc}}| \le 2\varepsilon H_0/(1-\varepsilon)\)
--- the only regime in which the frozen-rate ratio is a horizon. At a
stationary state (\(\dot A = 0\)) with \(g > 0\),
\(H_A^{\mathrm{loc}} = +\infty\) while
\(H_A^{\mathrm{gross}} < \infty\): the three quantities of Section 6.1
coincide only under a declared rate bracket, and the false implication
\(g > 0 \Rightarrow \dot A < 0\) is the reason.

\textbf{Proposition 26 (Sign of the frozen-rate error).} Let \(\dot A=-\varphi(A)\) with \(\varphi>0\) on
\((A_{\min},A_0]\), and write the true horizon and the frozen-rate ratio as
\[ T=\int_{A_{\min}}^{A_0}\frac{dA}{\varphi(A)}, \qquad H^{\mathrm{loc}}=\frac{A_0-A_{\min}}{\varphi(A_0)} . \]
If \(\varphi\) is nondecreasing in \(A\) then \(T\ge H^{\mathrm{loc}}\), and the frozen-rate number
is conservative. If \(\varphi\) is nonincreasing in \(A\) then \(T\le H^{\mathrm{loc}}\), and it is
optimistic.

\emph{Proof.} Nondecreasing \(\varphi\) gives \(\varphi(A)\ge\varphi(A_0)\) for \(A\le A_0\), hence
\(T\ge(A_0-A_{\min})/\varphi(A_0)\); the second case reverses the inequality. The comparison is the
standard one for one-dimensional monotone dynamics (Smith, 1995). \ensuremath{\square}

Proportional extraction is the first case, with \(T=q^{-1}\ln(A_0/A_{\min})\) against
\(H^{\mathrm{loc}}=(A_0-A_{\min})/(qA_0)\); at \(A_{\min}=0\) the two disagree by being infinite
against finite. The sign is estimable from the data the classification already requires, by
regressing the decline rate on the stock; it is declared, not computed, in Section 6.5.

\textbf{Proposition 27 (Curvature correction).} On the barrier distance \(u=A-A_{\min}\), let
\(\dot u=-\varphi(u)\) and define the dimensionless curvature number
\[ \kappa=\frac{u\,\ddot u}{\dot u^{2}}=\frac{u\,\varphi'(u)}{\varphi(u)} . \]
For the power-law family \(\varphi(u)=cu^{p}\) one has \(\kappa=p\) identically, and for \(\kappa<1\)
\[ T=\frac{H^{\mathrm{loc}}}{1-\kappa}, \qquad T<\infty\iff\kappa<1, \]
while for \(\kappa\ge1\) the integral diverges at the barrier and the family reaches it only
asymptotically, so the frozen-rate ratio understates by an unbounded factor. Constant extraction is
\(\kappa=0\), where the frozen-rate ratio is exact; stock-proportional decline is \(\kappa=1\),
where the true horizon is infinite against a finite ratio; and \(\kappa=\tfrac12\) doubles the
horizon exactly, which is the smallest correction with a real effect. The correction is a one-number
statement about the bias, and it inherits its input requirement: \(\kappa\) needs \(\ddot u\), which
is a second difference of the same noisy series whose first difference the classification already
distrusts, so it is reported as declared or not at all. Where \(\varphi\) is not of the family, only
the sign statement of Proposition 26 is available.

\textbf{Proposition 28 (Reserve-life crossover).} Let reserves obey \(\dot R=-(1-\eta)P\), with
production \(P=P_0e^{gt}\), \(g>0\), and let \(\tau=R_0/P_0\) be the reserve-life ratio. Exhaustion of
the reserve class occurs at
\[ T=\frac1g\ln\Bigl(1+\frac{g\tau}{1-\eta}\Bigr), \]
and \(T\ge\tau\) holds to first order in \(g\tau\) exactly when \(\eta\ge\tfrac12g\tau\). A
reserve-life ratio therefore bounds the reserve class from above only where reclassification keeps
pace with growth. On the article's own pinned record, \(\tau\approx309\) yr at
\(g=0.03\,\mathrm{yr}^{-1}\), the condition requires \(\eta\ge\tfrac12g\tau\approx4.6\), which
exceeds unity: reclassification would have to outpace extraction itself, and no declared reserve
convention admits that. At \(\eta=0\) the same record gives \(T=g^{-1}\ln(1+g\tau)\approx77.6\) yr
against the tabulated 309. The direction of the error is therefore fixed by the model class rather
than by the data, and the classification of Section 6.5 is strengthened, not changed: the ratio
remains an arithmetic relation whose promotion to a forecast is unavailable at every admissible
elasticity.

\emph{Proof.} Integrate \(\dot R=-(1-\eta)P_0e^{gt}\) to \(R(T)=0\) and solve; the comparison with \(\tau\)
is the first-order expansion \(\ln(1+x)=x-x^2/2+O(x^3)\) at \(x=g\tau/(1-\eta)\). \ensuremath{\square}


\subsubsection{6.3 Upper barriers, exit times, and
maintainability}\label{upper-barriers-exit-times-and-maintainability}

The lower-barrier setting of Section 6.1 is one half of the story. For
each moiety \(m\), define the lower and upper exit times
\[\tau_m^- = \inf\{ t \ge 0 : S_m(t) \le \underline{B}_m(t) \}, \qquad \tau_m^+ = \inf\{ t \ge 0 : S_m(t) \ge \overline{B}_m(t) \}, \qquad \inf\varnothing = \infty,\]
and the overall admissibility exit time
\[\tau_{\mathrm{exit}} = \min_m \{ \tau_m^-, \tau_m^+ \};\] horizon
safety on \([0,T]\) is \(\tau_{\mathrm{exit}} > T\). Two disciplines
attach. First, equality at the hitting time ---
\(S_m(\tau_m^-) = \underline{B}_m(\tau_m^-)\) --- requires continuity of
both \(S_m\) and \(\underline{B}_m\) and appropriate initial separation;
if fluxes or barriers can jump, the stock can cross the barrier without
satisfying equality. Second, lower barriers need not be exhaustion
thresholds: the diagnostic distinguishes physical exhaustion
(\(S_m = 0\)), functional failure (\(S_m = B_m^{\mathrm{func}}\)), a
resilience or regime-shift threshold, an economically recoverable
reserve, and a minimum service-supporting stock --- the term
``exhaustion'' is reserved for \(S_m = 0\), and all other thresholds are
barrier violations.

Upper-barrier violations matter symmetrically: a system may satisfy
every lower-barrier condition while violating an upper barrier,
\(S_m(t) > \overline{B}_m(t)\). Atmospheric CO\textsubscript{2} accumulation violates an
upper concentration barrier while fossil carbon stocks decline; the
depletion diagnostic must check both barriers. And a stock can remain
above a barrier until an assessment horizon \(T\) and still be
unsustainable thereafter: if the assessment claims sustainability rather
than finite-horizon admissibility, it requires a terminal condition
\[x(T) \in K_{\mathrm{maint}},\] where
\(K_{\mathrm{maint}} = \{ x : \exists \text{ an admissible continuation satisfying all barriers for } t \ge T \}\)
--- the set from which barrier safety is indefinitely maintainable, the
viability kernel of the barrier set under the admissible controls
(Aubin, 1991). The finite-horizon diagnostic is a necessary but not
sufficient condition for sustainability; the full certificate requires
the terminal state to lie in the maintainability set.

One consequence for reporting. The exit time is itself a minimum, over moieties and over both barrier
signs, so a statement about its joint distribution is already a componentwise statement: no product
of marginal probabilities is required, and none is admissible in its place.


\subsubsection{6.4 Robust semantics}\label{robust-semantics}

For uncertain parameters \(\theta \in \Theta\) and admissible
disturbances \(d \in \mathcal{D}\), robust barrier safety is
\[\mathrm{RobustBarrierSafe}(x(\cdot)) \iff \underline{B}_m(t) \le S_m(t; \theta, d) \le \overline{B}_m(t) \quad \forall m, \forall t, \forall \theta \in \Theta, \forall d \in \mathcal{D},\]
and the depletion-horizon classification is fourfold: nominal
(\(\theta = \theta_0\), \(d = 0\)); worst-case
(\(\inf_{\theta,d} \tau_{\mathrm{exit}}(\theta,d)\)); probabilistic
(\(\Pr[\tau_{\mathrm{exit}} > T] \ge 1 - \varepsilon\)); and
scenario-conditioned (\(\tau_{\mathrm{exit}} \mid \theta = \theta_s\)).
The componentwise diagnostic applies to each scenario, and the
conjunctive criterion applies within each scenario and across scenarios.
No single number is promoted across the four classes without a declared
map.

\subsubsection{6.5 Application classifications at their exact
status}\label{application-classifications-at-their-exact-status}

The classification matrix records, quantity by quantity, what each
application computes and what each object is; the numerical exhibits
that follow are worked instances of the constructions, and the
classification of each row does not depend on the magnitudes.

\begin{longtable}[]{@{}
  >{\raggedright\arraybackslash}p{(\columnwidth - 8\tabcolsep) * \real{0.2000}}
  >{\raggedright\arraybackslash}p{(\columnwidth - 8\tabcolsep) * \real{0.2000}}
  >{\raggedright\arraybackslash}p{(\columnwidth - 8\tabcolsep) * \real{0.2000}}
  >{\raggedright\arraybackslash}p{(\columnwidth - 8\tabcolsep) * \real{0.2000}}
  >{\raggedright\arraybackslash}p{(\columnwidth - 8\tabcolsep) * \real{0.2000}}@{}}
\toprule\noalign{}
\begin{minipage}[b]{\linewidth}\raggedright
Time-like quantity
\end{minipage} & \begin{minipage}[b]{\linewidth}\raggedright
Question answered
\end{minipage} & \begin{minipage}[b]{\linewidth}\raggedright
G3P \(L_{\mathrm{hist}}^{\mathrm{anom}}\)
\end{minipage} & \begin{minipage}[b]{\linewidth}\raggedright
Phosphate \(T_{\mathrm{reserve}}\)
\end{minipage} & \begin{minipage}[b]{\linewidth}\raggedright
Fisheries \(\Theta_F\)
\end{minipage} \\
\midrule\noalign{}
\endhead
\bottomrule\noalign{}
\endlastfoot
\(J_A^{\mathrm{gross}}\), \(H_A^{\mathrm{gross}}\) & turnover /
dependency & no & no & gross-loss analogue only --- not a member
(\S{}6.5.4) \\
\(H_A^{\mathrm{loc}}\) & frozen net-rate ratio & no (anomaly, not stock)
& no (classification, not stock) & no (no net \(\dot B\)) \\
\(T_A\) & scenario hitting time & no & no & no \\
What it is & --- & record-relative statistical index & arithmetic ratio
of an economic class & removals-only pressure scale \\
\end{longtable}

\paragraph{6.5.1 Groundwater anomaly-persistence
indices}\label{groundwater-anomaly-persistence-indices}

The G3P column of the classification matrix is the groundwater case. The
Global Gravity-based Groundwater Product (G3P v1.12; G\"untner et al.,
2024; the GRACE line it descends from is Tapley et al., 2004) provides
monthly groundwater-storage anomalies relative to a reference period
rather than absolute aquifer volumes. For a basin-mean anomaly series
over the reported April 2002--September 2023 window, the linear-trend
anomaly persistence index is
\[L_{\mathrm{hist}}^{\mathrm{anom}} = \frac{a_{\mathrm{latest}} - a_{\mathrm{hist,min}}}{\bigl[ -\widehat{\dot a} \bigr]_+},\]
the fitted distance to the series' own historical minimum divided by the
fitted decline rate. The four-basin record: Indo-Gangetic \(-49.7\)
cm/yr with index \(\approx 2.7\) yr; North China Plain \(-18.6\) with
\(\approx 7.9\); Central Valley \(-16.1\) with \(\approx 9.5\); La
Mancha \(-3.2\) with \(\approx 21.4\).

Classification, stated at the product's own status: a statistical
anomaly index with units of time --- not the physical stock ratio
\(H_A^{\mathrm{loc}}\) and not a forecast of aquifer exhaustion. Its value depends on the product window, basin mask, anomaly reference, and linear-trend convention
--- the reference period is the product's declared long-term mean over April 2002 to December 2020,
the coverage ends September 2023, and the producer documents a faulty snow-water-equivalent entry for
June 2005 that propagates into groundwater storage, recommending that the month be excluded; the
window above spans it, so a re-derivation must declare whether June 2005 is retained, since it can
move both the fitted trend and the series' own historical minimum; a physical \(H_A^{\mathrm{loc}}\) requires an
absolute stock estimate and a net stock derivative (aquifer geometry or
saturated thickness together with storage parameters), not an anomaly
series alone.

A structural point sharpens the boundary. The access structure --- a
well or an index well --- is infrastructure rather than the resource: it
draws on stored water, the stock, which is replenished by recharge, the
flow. The anomaly series is a measure of the stored stock as observed at
the access point; it measures neither the access infrastructure nor the
recharge. An index built on it therefore cannot distinguish a drawdown
of stored water that is recoverable on the recharge time scale from one
that is not --- although heavy over-extraction can make the loss
permanent through compaction, subsidence, or saline intrusion, in which
case it is not recoverable at all. The index is exactly the
record-relative object analysed in Section 7.3, and its interpretive
boundary is that record-relativity (Section 7.7).

\paragraph{6.5.2 The applied depletion-horizon
tables}\label{the-applied-depletion-horizon-tables}

Component-resolved depletion horizons on public data products are
tabulated below, computed without fitting any dynamical parameter of the
reduced systems. Rows marked with a dagger (\dag{}) are quarantined and must
not be taken at face value: the Australia phosphate row dates to a
pre-2026 reserve vintage (\S{}6.5.3), and the Indo-Gangetic groundwater
magnitude sits an order of magnitude beyond published basin-mean trends
(the quarantine note below); both are retained only as worked instances
of the constructions, and neither enters any classification below.

\begin{longtable}[]{@{}
  >{\raggedright\arraybackslash}p{(\columnwidth - 8\tabcolsep) * \real{0.2000}}
  >{\raggedright\arraybackslash}p{(\columnwidth - 8\tabcolsep) * \real{0.2000}}
  >{\raggedright\arraybackslash}p{(\columnwidth - 8\tabcolsep) * \real{0.2000}}
  >{\raggedright\arraybackslash}p{(\columnwidth - 8\tabcolsep) * \real{0.2000}}
  >{\raggedright\arraybackslash}p{(\columnwidth - 8\tabcolsep) * \real{0.2000}}@{}}
\toprule\noalign{}
\begin{minipage}[b]{\linewidth}\raggedright
Basin
\end{minipage} & \begin{minipage}[b]{\linewidth}\raggedright
Trend (cm/yr)
\end{minipage} & \begin{minipage}[b]{\linewidth}\raggedright
2023 anomaly (cm)
\end{minipage} & \begin{minipage}[b]{\linewidth}\raggedright
Window minimum (cm, implied)
\end{minipage} & \begin{minipage}[b]{\linewidth}\raggedright
Horizon to window minimum (yr)
\end{minipage} \\
\midrule\noalign{}
\endhead
\bottomrule\noalign{}
\endlastfoot
Indo-Gangetic (N. India)\dag{} & \(-49.7\) & \(-414\) & \(-548\) &
\(\approx 2.7\) \\
North China Plain & \(-18.6\) & \(-145\) & \(-292\) & \(\approx 7.9\) \\
Central Valley (US) & \(-16.1\) & \(-84\) & \(-237\) &
\(\approx 9.5\) \\
La Mancha (Spain) & \(-3.2\) & \(-20\) & \(-88.5\) & \(\approx 21.4\) \\
High Plains (US) & \(-7.9\) & \(-160\) & \(-160\) & already at
minimum \\
global mean & \(-0.4\) & \(-14\) & \(-33.0\) & \(\approx 47.5\) \\
\end{longtable}

\dag{} \emph{Quarantine note, adjacent to the row it marks (recorded
data-vintage decision: the row stays first, daggered).} The
Indo-Gangetic magnitude must not be reused numerically: it sits an order
of magnitude beyond published basin-mean trends and awaits re-derivation
from the product's basin masks --- its full provenance is the supplementary's S5.4.

The basin rows are reported extractions from the G3P v1.12 basin series, used here only to exhibit the
index construction of Section 6.5.1 and never as product-endorsed values: the window-minimum column is implied
arithmetically through the index formula, every row must be re-derived from the product's basin masks before any
numerical reuse, and the extraction provenance --- including why the Indo-Gangetic row is the extreme case and
how its fitted segment convention enters --- is recorded in the supplementary's S5.4 with the vintage record it
belongs to. The classification status assigned below does not depend on the magnitudes.

\begin{longtable}[]{@{}
  >{\raggedright\arraybackslash}p{(\columnwidth - 6\tabcolsep) * \real{0.2500}}
  >{\raggedright\arraybackslash}p{(\columnwidth - 6\tabcolsep) * \real{0.2500}}
  >{\raggedright\arraybackslash}p{(\columnwidth - 6\tabcolsep) * \real{0.2500}}
  >{\raggedright\arraybackslash}p{(\columnwidth - 6\tabcolsep) * \real{0.2500}}@{}}
\toprule\noalign{}
\begin{minipage}[b]{\linewidth}\raggedright
Country
\end{minipage} & \begin{minipage}[b]{\linewidth}\raggedright
Reserves (kt)
\end{minipage} & \begin{minipage}[b]{\linewidth}\raggedright
Reserve-life horizon (yr)
\end{minipage} & \begin{minipage}[b]{\linewidth}\raggedright
Implied production (kt/yr)
\end{minipage} \\
\midrule\noalign{}
\endhead
\bottomrule\noalign{}
\endlastfoot
China & \(3{,}400{,}000\) & \(\approx 28\) & \(121{,}429\) \\
United States & \(1{,}000{,}000\) & \(\approx 45\) & \(22{,}222\) \\
Jordan & \(820{,}000\) & \(\approx 62\) & \(13{,}226\) \\
Morocco & \(50{,}000{,}000\) & \(\approx 1{,}250\) & \(40{,}000\) \\
Australia\dag{} & \(5{,}800{,}000\) & \(\approx 2{,}088\) & \(2{,}778\) \\
World (reserves) & \(74{,}000{,}000\) & \(\approx 309\) &
\(239{,}482\) \\
World (resources, \(\varepsilon = 0.10\)) & \(>300{,}000{,}000\) &
\(>1{,}125\) & \(240{,}000\) \\
\end{longtable}

The fisheries column reports the archived-depletion-horizon (ADH)
pure-decay proxy
\(\mathrm{ADH} = F^{-1}\log(\mathrm{SSB}_{\mathrm{now}}/(0.2\max\mathrm{SSB}))\)
under current fishing mortality \(F\), with median \(\approx 1.8\) yr
across the 43 assessed stocks with finite spawning-stock-biomass (SSB)
and \(F\) series --- the archived pull, kept in place as the headline
cohort by the recorded data-vintage decision --- computed with
\(\mathrm{ADH} = 0\) entered for the eight stocks already at or below
the reference, the zero convention of the source table's caption, which
the median includes. Two disclosures accompany the headline value (both
recorded): the archived 43-stock cohort is reproduced by neither public
RAM Legacy release --- S5's version-sensitivity record, cross-referenced
here, is the retraction and the row-by-row re-verification --- and the
v4.66 public-release broad-cohort reading (454 stocks, median \(3.39\)
yr) is the primary public-release comparison, executed in S5 and
detailed below. The qualifying positive sub-cohort (\(F > 0\) and
\(\mathrm{SSB}_{\mathrm{now}} > B_{\lim}\); 35 stocks) has median
\(2.9\) yr; both medians come from the archived pull alone. The cohort
is a selected class, not a random sample of assessed stocks: all 43
stocks are small pelagics (18 anchovy, 20 herring, 4 sprat, 1 sardine),
the fast-maturing class the companion review screen (Abaee, 2026, doi:10.5281/zenodo.22554297) selects by its annual-review eligibility criterion (42 of the
43 are that screen's annual-managed spectral-null stocks, per the source
caption). The \(\approx 1.8\) yr median is therefore a class-specific
diagnostic for fast-maturing, annually managed pelagics, not a statistic
of assessed fisheries in general; the executed broad-cohort comparison
(S5) runs the same protocol on the full public release --- 454 stocks,
median \(3.39\) yr, the upper end carried by the long-lived groups
(elasmobranchs \(11.5\), sebastids \(9.0\), pleuronectids \(6.0\) yr)
--- and only \(2\%\) of random 43-stock draws from that broad cohort
have medians at or below the class cohort's \(1.79\) yr.

The extract is the RAM Legacy cohort of Ricard et al.~(2012) at release v4.66 (Zenodo
14043031, 6 November 2024), recorded as the cohort file of the supplementary's S5 record; its vintage is
verified from the extract's own contents, four of six published \(F\) values reproducing that release
exactly, and the pull has been re-verified row by row against the formula --- all 43 rows
reproduce
\(\mathrm{ADH} = \max(0,\, F^{-1}\log(\mathrm{SSB}/B_{\lim}))\) with
\(B_{\lim} = 0.2 \max \mathrm{SSB}\). The value is reported with its
cohort conditions and is not promoted to a forecast. Because cohort
composition is database-version-dependent, every cohort statistic is
pinned to the archived pull: the quartile summary of \(F\) and
\(\log(\mathrm{SSB}/B_{\lim})\) over the cohort belongs to that pull
alone, and no cohort statistic is quoted from a different database
version. The cohort protocol is fully specified in the accompanying
supplementary material (S5) --- including the zero entries for stocks at
or below the reference, which enter the median --- together with the
version-sensitivity record: on the public RAM Legacy releases the same
protocol qualifies 415 stocks (v4.44, median \(2.57\) yr) and 454
(v4.66, median \(3.39\) yr) --- both reproduced to printed precision
under the recovered micro-specification (S5) --- neither reproducing the
archived 43-stock cohort --- the archived pull's 43-stock list and
extract-time series state --- supplied and re-verified --- differ from
both public releases.

The scope discipline is the tables' load-bearing content: none of the
reported numbers is a computed instance of any model's first-hitting
time --- the groundwater column is a trend-to-window-minimum
extrapolation, the phosphate column a reserve-life ratio, and the
fisheries column a pure-decay proxy with recruitment omitted. They are
descriptive, component-resolved diagnostics in the two-pool logic of the
taxonomy, not dynamical predictions. The equal-weight inverse-horizon
score of the four basins still above their window minimum and world
phosphate reserves, displayed as \textbf{Non-example 1} --- a deliberate
boundary of aggregation, not a score of the framework,
\[\Sigma_{\mathrm{reserves}} \approx \frac{1}{5}\left( \frac{1}{2.7} + \frac{1}{7.9} + \frac{1}{9.5} + \frac{1}{21.4} + \frac{1}{309} \right) \approx 0.130\ \mathrm{yr}^{-1},\]
is a ranking device, not a componentwise certificate: it mixes basins
and reserves, incommensurable objects under the type structure of Definition 47,
and is retained only to mark the boundary of legitimate aggregation. The
score is exhibited here as the worked instance of the noncompensation
boundary of Section 10.1: a positive aggregate coexisting with
componentwise deficits by construction --- admissible as communication,
inadmissible as certification.

\paragraph{6.5.3 The phosphate reserve-life
ratio}\label{the-phosphate-reserve-life-ratio}

The phosphate column of the classification matrix is the reserve-life
ratio. At constant current production \(C_G\), the reserve-life ratio is
\(T_{\mathrm{reserve}} = G_{\mathrm{reserve}}/C_G\); at approximately
\(74{,}000{,}000\) kt (\(74{,}000\) Mt) of world reserves and
\(240{,}000\) kt/yr of production (U.S. Geological Survey, 2026) this is
approximately \(309\) years. The arithmetic is internally consistent as
a reserve-life ratio to zero; it is not a physical exhaustion forecast,
because reserve classification changes with prices, technology,
exploration, and regulation --- the point made independently, and
forcefully, by Illakwahhi, Vegi, and Srivastava (2024) for the
single-source USGS data behind the influential phosphate depletion
estimates, and standard in mineral economics, where reserves have grown
through a century of rising production for copper (Tilton, 2003; Tilton
and Lagos, 2007). The reserves/resources split discipline is part of the
classification: a resource-threshold calculation
\(T_{\mathrm{resource},10\%} = 0.9\, G_{\mathrm{resource}}/C_G\) answers
a different question and must not share a column with the reserve-life
ratio without an explicit convention label; the reserve classification
is economic --- US reserves have remained near \(1{,}000{,}000\) kt
while cumulative production since 1996 is of order \(600{,}000\) kt ---
and the resource-based world horizon (\(\approx 1{,}125\) yr at
\(\varepsilon = 0.10\)) is more than three times the reserve-based
figure (\(\approx 309\) yr); the two-compartment split is what prevents
these from being collapsed into one number. The implied-production
column of the Section 6.5.2 table reproduces the production figure each
horizon assumes (production = reserves/horizon) and thereby exposes the
source arithmetic; the country horizons reproduce the recorded
MCS-vintage ratios. The vintage is pinned once: the single pinned source
of record is the Mineral Commodity Summaries (MCS) 2026 (U.S. Geological
Survey, 2026), and every figure this article quotes at pin status is
that vintage's --- the pinned source's 2025 world-production column
\(\approx 250{,}000\) kt and Australia's reserves \(120{,}000\) kt
(JORC-compliant). The displayed country rows remain at their recorded
pre-2026 vintage under the quarantine dagger (the Australian
\(5{,}800{,}000\) kt among them, retained rather than blanked), kept in
place as worked instances of the reserve-life construction; completing
the re-pin --- replacing the displayed rows row by row with the pinned
vintage's per-country reserve figures --- is the registered open data
action, it requires the per-country MCS 2026 reserve table, and no
displayed classification depends on it.

\paragraph{6.5.4 The fisheries removals-only pressure
time}\label{the-fisheries-removals-only-pressure-time}

The fisheries column of the classification matrix is the removals-only
pressure time. When \(\mathrm{SSB}_{\mathrm{now}} > B_{\lim} > 0\) and
\(F_{\mathrm{now}} > 0\), define
\(R_B = \log(\mathrm{SSB}_{\mathrm{now}}/B_{\lim})\) and
\[\Theta_F = \frac{R_B}{F_{\mathrm{now}}},\] the fishing-only
time-to-reference: the crossing time of the deliberately incomplete
comparison process \(\dot B = -F_{\mathrm{now}}B\). With
\(B_{\lim} = 0.2 \max \mathrm{SSB}\) this is the construction tabled as
ADH in Section 6.5.2; the two notations are kept because the boundary
hypotheses stated here (\(F_{\mathrm{now}} > 0\),
\(\mathrm{SSB}_{\mathrm{now}} > B_{\lim}\)) are exactly the conditions
of the positive sub-cohort of Section 6.5.2 (35 stocks, median \(2.9\)
yr); the reported Section 6.5.2 median (\(\approx 1.8\) yr) additionally
carries the eight zero entries for stocks at or below the reference, per
the zero convention of the source caption. It is a removals-only
pressure time scale --- the time unit comes from rescaling a
stock-reference margin by one isolated gross-loss rate. Because
recruitment, somatic growth, maturation, natural mortality, density
dependence, environmental forcing, and future policy are omitted,
\(\Theta_F\) is not a net biomass depletion diagnostic, not a
demographic hitting-time estimate, and not a member of the
\(J^{\mathrm{gross}}\)--\(H^{\mathrm{loc}}\)--\(T_A\) hierarchy. A
genuinely local biomass-decline ratio
\(H_B^{\mathrm{loc}} = (B - B_{\lim})/[-\dot B]_+\) would require a
compatible net \(\dot B\) estimate, and a demographic hitting time a
fully specified population model; RAM Legacy SSB and \(F\) data (Ricard
et al., 2012) do not by themselves supply these quantities or models.
Spawning biomass is not an abiotic support pool. The construction is
retained specifically to show why an isolated gross-removal time scale
must not be promoted to a net depletion diagnostic.

The collective implication for material-flows measurement is that the
three published ``depletion time'' numbers --- G3P index, phosphate
reserve-life, fisheries removals-only time --- answer three distinct
questions at three distinct evidentiary levels.

\begin{center}\rule{0.5\linewidth}{0.5pt}\end{center}

\textbf{Proposition 32 (Boundedness of the persistence index on the declared trend class).} Let
\(a_k=-\beta k+\varepsilon_k\) for \(k=1,\dots,n\), with \(\varepsilon_k\) independent of scale
\(\sigma\) and \(\beta>0\). Within the declared class, the record minimum lies at the end of the
record once the trend dominates the noise, so the fitted distance from the current level up to that
minimum is a noise-scale gap and the index
\[ \frac{a_n-\min_ka_k}{\text{fitted decline rate}} \]
is \(O_{\mathbb P}(\sigma/\beta)\): bounded in probability as the record length grows, and
independent of the stock. A simulation of 400 replicates at each of \(n=10^2,10^3,10^4,10^5\) with
\(\beta=\sigma=1\) returns a mean index of 0.21, 0.25, 0.21 and 0.25 yr, and a 90th percentile of
0.81, 0.84, 0.95 and 1.00 yr, while the stock implied by the same series falls by five orders of
magnitude; the median is zero at every length, because under a downward trend the current level is
usually the record minimum. The index is therefore flat in the record length at the noise-to-trend
scale \(\sigma/\beta\), which is the sense in which it cannot be extended into a horizon; the code
and the seed are archived with the supplementary material. Lengthening or densifying an anomaly
record cannot produce a horizon, because the statistic converges to a trend-detection quantity
expressed in years. This is a property of the declared linear-trend class and of the reported
simulation, not a claim about any basin.



\subsubsection{6.6 What an aggregate record fixes}\label{what-an-aggregate-record-fixes}

\textbf{Definition 38 (Identifiability set of an event time).} Fix the readout, the declared boxes and
bounds, and an observed aggregate record \(z\). The \emph{fibre} \(\mathcal F(z)\) is the set of admissible
component initial states whose induced aggregate trajectory equals \(z\). The \emph{identifiability set} of
the first component exit time is
\[
\mathrm T(z) \;=\; \bigl\{\tau(x_0) : x_0 \in \mathcal F(z)\bigr\},
\]
and the event time is identifiable from the aggregate record if and only if \(\mathrm T(z)\) is a
singleton.

\textbf{Proposition 39 (The set is a polytope image, and it is usually an interval).} On the decay pair of
Proposition 31, with \(\dot x_i = -x_i\), \(Z(t) = 100 e^{-t}\) and barrier \(x_i \ge 1\), the fibre is
\(\{x_1 + x_2 = 100,\ 1 \le x_1, x_2 \le 99\}\) and
\[
\mathrm T(z) \;=\; \bigl[\,0,\ \log 50\,\bigr] \;=\; [\,0,\ 3.9120\,] \ \text{yr},
\]
its supremum attained at the balanced start \((50, 50)\) and its infimum approached as either component
nears its barrier. Both recorded instances lie in the set, \(0.6931\) yr from \((2, 98)\) and \(3.9120\) yr
from \((50, 50)\), and the aggregate record is silent on everything between them. More generally, where the declared constraints are linear the fibre is a polytope, and a continuous
event time has an interval image on a connected fibre: \(\mathrm T(z)\) is that interval, and its endpoints
are its maximum and minimum over the polytope. Those two values are programmes rather than aspirations
exactly where the event time is declared in polyhedral form, \(\tau(x) = g(\min_{j \le k}(a_j^{\top} x + b_j))\)
with \(g\) continuous and increasing: the upper endpoint is the single linear programme over \((x, r)\) with
\(r \le a_j^{\top} x + b_j\) for every \(j\), and the lower endpoint is the minimum over the polytope's vertices,
finite to enumerate, because a minimum of affine functions is concave and a concave function attains its
minimum at an extreme point. No such reduction holds for an arbitrary continuous event time: its extremum sits
where it sits, and in the instance above it is attained at the balanced start, in the interior of the fibre. Monotonicity of the event time along the fibre
is neither needed nor true in general --- an exit time is typically the minimum of several monotone branch
functions, which is not itself monotone --- and where the declared set is non-convex the fibre can split,
in which case the same two programmes bound the value set from outside instead of recovering it, and identifiability from an aggregate record holds exactly when the declared
component data reduce the fibre to a point --- which is a statement about the input, never about the
aggregate.

\emph{Proof.} \(\tau(x_0) = \min\{\log x_1(0), \log x_2(0)\}\) with \(x_2(0) = 100 - x_1(0)\): the function is
increasing then decreasing in \(x_1\) on \([1, 99]\), symmetric about \(50\), with maximum \(\log 50\) and
limit \(0\) at the ends, so its image is \([0, \log 50]\). The fibre is an intersection of the declared
boxes with the invariant hyperplane, hence a convex polytope; \(\tau\) is continuous, so its image is a
connected interval. Its maximum is not a vertex value --- \(\min_j\) of affine functions is concave, and a
concave function attains its maximum where it likes, here at \(x_1 = x_2 = 50\) --- while its minimum over the
segment is attained at an end, where \(\tau\) tends to \(0\). The two programmes of the statement are read off
those two facts, and the fact that one of them is an interior value is why the polyhedral form is quoted with
its hypothesis instead of as a general method. \ensuremath{\square}
\subsection{7. First-Passage Semantics on Declared
Surrogates}\label{first-passage-semantics-on-declared-surrogates}

\subsubsection{7.1 Two objects, not one}\label{two-objects-not-one}

The ledger's own first-passage object is the model hitting time of
Definition 5 --- a quantity on trajectories of the mass-conserved ledger
or of a named reduced system. The public-data quantities of Section 6.5
are constructed proxies on observed series. The distinction is the entry
discipline of this section: the surrogates below do not compute the
ledger's hitting time, do not complete the ledger stochastically, and do
not identify physical failure thresholds.

The surrogates used here are standard objects, and the section's content is the discipline attached
to them rather than the passage-time formulae. The inverse-Gaussian law and its parameterisations are
those of the first-passage literature on Brownian motion with drift (Chhikara and Folks, 1989;
Redner, 2001); the stochastic machinery is It\^o calculus in its textbook form (\O{}ksendal, 2003). The
comparison between a surrogate mean and the deterministic margin is a monotone-dynamics statement
(Smith, 1995), and the barrier construction is the deterministic safety argument used in hybrid-systems
verification, whose stochastic extension is the nearest formal relative of the record-relative
discipline of Section 7.4 (Prajna and Jadbabaie, 2004; Prajna et al., 2007). None of these sources
supplies a calibrated input: every drift, barrier and noise scale entering Sections 7.2 to 7.6 is
declared by the analysis, and the results are statements about the declared class. Transfer noise obeys the same typing as
transfer fluxes: process noise on a conserved moiety must be zero-sum across the pools it moves between,
or must carry an explicit boundary term; where it does neither, the violation is a residual and belongs in
the budget of Definition 23, not in the drift --- a diffusion that creates mass is not a conservative ledger
with noisy data but a different object. The converse discipline holds as well: a bracket on the drift of a
selected set does not transfer to a bracket on its support or on its exit time, so no drift bound in this
section is a deterministic hitting time, and the surrogate means are means.


\subsubsection{7.2 The observed-drift Brownian
surrogate}\label{the-observed-drift-brownian-surrogate}

\textbf{Definition 6 (Observed-drift Brownian surrogate).} \emph{Let
\(A_0\) (local to Sections 7.2--7.4; not the half-saturation constant of
Section 2.2) be the latest observed anomaly and
\(\mu = \widehat\mu < 0\) the fitted drawdown rate. On the scale of the
tabulated series define}
\[A(t) = A_0 + \mu t + \varsigma W_t, \qquad A(0) = A_0 > A_{\min}^{\mathrm{win}},\]
\emph{where \(W\) is a standard Wiener process, \(\varsigma > 0\) a
chosen noise scale, and the process is stopped at first reaching the
record-relative barrier \(A_{\min}^{\mathrm{win}}\).}

This is a statistical surrogate for the empirical trend extrapolation.
It is not a hydrological constitutive law, is not mass-conserving, and
is not a perturbation or stochastic completion of the ledger's
active-pool equation or of the finite-donor primitive system of Section
2.2. The non-completion non-claim is part of the definition.

\subsubsection{7.3 The inverse-Gaussian groundwater first
passage}\label{the-inverse-gaussian-groundwater-first-passage}

\textbf{Proposition 18 (Inverse-Gaussian first passage --- a standard
fact, stated for notation).} \emph{Let
\(T_{\mathrm{GW}} = \inf\{ t > 0 : A(t) \le A_{\min}^{\mathrm{win}} \}\)
for the process of Definition 6 and
\(d = A_0 - A_{\min}^{\mathrm{win}} > 0\). Conditional on treating
\(\mu\) and the barrier as fixed,}
\[T_{\mathrm{GW}} \sim \mathrm{IG}(\nu, \lambda), \qquad \nu = \frac{d}{|\mu|}, \qquad \lambda = \frac{d^2}{\varsigma^2},\]
\emph{in the mean--shape parameterization; in particular
\(\mathbb{E}[T_{\mathrm{GW}}] = \nu = H^{\mathrm{win}}_{\mathrm{GW}}\)
--- the deterministic horizon to the window minimum, \(= d/|\mu|\) ---
and
\(\operatorname{Var}(T_{\mathrm{GW}}) = \nu^3/\lambda = d\,\varsigma^2/|\mu|^3\).}

\emph{Proof.} The first-passage time of a Brownian motion with constant
negative drift to a lower barrier is inverse Gaussian --- the classical
first-passage result (Chhikara and Folks, 1989; Redner, 2001) --- with
the stated mean and shape parameters; the standard inverse-Gaussian
moments give the displayed mean and variance. \ensuremath{\square}

The mean of the stochastic surrogate equals the deterministic
trend-to-window-minimum ratio of Section 6.5.1; that equality is the
precise sense in which the tabled groundwater numbers are first-passage
means of a declared surrogate.

\textbf{Corollary 19 (Zero-noise limit and median).} \emph{As
\(\varsigma \to 0^+\),
\(T_{\mathrm{GW}} \to H^{\mathrm{win}}_{\mathrm{GW}}\) in probability,
and at \(\varsigma = 0\) the deterministic trajectory reaches the
barrier exactly there. For every finite \(\varsigma > 0\) the
inverse-Gaussian median \(m\) satisfies}
\[m < \nu = H^{\mathrm{win}}_{\mathrm{GW}}, \qquad F_T(\nu) = \frac{1}{2} + e^{2\lambda/\nu}\,\Phi\!\left( -2\sqrt{\lambda/\nu} \right) > \frac{1}{2}.\]
\emph{The variance scales as \(\varsigma^2\) and the standard deviation
and small-noise quantile widths as \(\varsigma\). The median below the
mean is the inverse Gaussian's right skew toward short passage times;
the inequality must not be inverted.}

\emph{Proof.} Evaluate the inverse-Gaussian CDF
\(F_T(t) = \Phi(\sqrt{\lambda/t}(t/\nu - 1)) + e^{2\lambda/\nu}\Phi(-\sqrt{\lambda/t}(t/\nu + 1))\)
at \(t = \nu\): the first term is \(\Phi(0) = 1/2\) and the second is
strictly positive for finite \(\lambda\), so the median lies strictly
below the mean; the concentration statement follows from the variance. \ensuremath{\square}

These are conditional distributional statements about the surrogate.
They are not corrections to the tabled years, and they do not show that
physical water mass is depleted faster.

\subsubsection{7.4 The record-relative barrier
discipline}\label{the-record-relative-barrier-discipline}

The barrier \(A_{\min}^{\mathrm{win}}\) is selected from the same finite
observation window used to estimate \(\widehat\mu\). It is therefore a
path-dependent, record-relative threshold, not an independently
identified hydrological failure floor; future passage below it
represents a record-breaking stress event under the surrogate, not
physical exhaustion. Three boundary facts complete the discipline.

\begin{enumerate}
\def\labelenumi{\arabic{enumi}.}
\tightlist
\item
  \textbf{Already-at-minimum.} If \(A_0 = A_{\min}^{\mathrm{win}}\), the
  stopping-time convention gives \(T_{\mathrm{GW}} = 0\)
  deterministically for every \(\varsigma\); the inverse-Gaussian family
  has a degenerate boundary limit concentrated at zero, and
  \(\mathrm{IG}(0,0)\) is not an ordinary inverse-Gaussian distribution.
  Zero cells report zero relative to the selected observational barrier
  --- not zero physical uncertainty and no confirmation of collapse.
\item
  \textbf{Independent physical thresholds.} If an independent physical
  threshold \(A^\sharp < A_{\min}^{\mathrm{win}}\) is specified, the
  same constant-drift surrogate gives the conditional mean
  \(\mathbb{E}[T^\sharp] = (A_0 - A^\sharp)/|\mu|\), longer than the
  record-relative proxy because the barrier is lower. This is a
  statement within the surrogate, not a general lower-bound theorem for
  the physical ledger, whose drift and state coupling may differ.
\item
  \textbf{Classification.} The load-bearing content is the
  interpretation boundary itself: a record-relative barrier makes the
  passage time a property of the observation window, and no reading of
  the tabled numbers escapes that qualification.
\end{enumerate}

\subsubsection{7.5 The geometric-Brownian fisheries first
passage}\label{the-geometric-brownian-fisheries-first-passage}

\textbf{Proposition 20 (Geometric-Brownian correction --- a standard
fact, stated for notation).} \emph{Let
\(dB_t = -hB_t\, dt + \varsigma B_t\, dW_t\) under the It\^o convention
with \(h > 0\) and \(0 < B_{\min} < B_0\), and
\(T_{\mathrm{fish}} = \inf\{ t > 0 : B_t \le B_{\min} \}\). Then}
\[T_{\mathrm{fish}} \sim \mathrm{IG}(\nu_F, \lambda_F), \qquad \nu_F = \frac{\log(B_0/B_{\min})}{h + \varsigma^2/2}, \qquad \lambda_F = \frac{\log(B_0/B_{\min})^2}{\varsigma^2},\]
\emph{so
\(\mathbb{E}[T_{\mathrm{fish}}] = \log(B_0/B_{\min})/(h + \varsigma^2/2)\);
as \(\varsigma \to 0^+\) this converges to the deterministic pure-decay
horizon when \(h = F\) and \(B_{\min} = B_{\lim}\).}

\emph{Proof.} It\^o's lemma (\O{}ksendal, 2003) gives
\(d\log B_t = -(h + \varsigma^2/2)\, dt + \varsigma\, dW_t\), so the
logarithmic threshold is a Brownian first-passage problem with initial
distance \(\log(B_0/B_{\min})\) and downward drift
\(h + \varsigma^2/2\); Proposition 18 applies. Under the Stratonovich
convention the log-drift would be \(-h\) and the deterministic limit
would match the pure-decay horizon \(h = F\) exactly: the
\(\varsigma^2/2\) shortening is the It\^o choice, not a property of the
physical process. \ensuremath{\square}

For fixed arithmetic drift and the It\^o parameterization, the
finite-noise mean is strictly shorter than the deterministic horizon.
This is a property of the chosen surrogate parameterization; it is not a
universal claim that environmental variability accelerates physical
biomass loss. The construction joins the removals-only classification of
Section 6.5.4 --- the same pure-decay process, now under a declared
stochastic surrogate.

\subsubsection{7.6 The constant-production phosphate passage
time}\label{the-constant-production-phosphate-passage-time}

Under the deterministic surrogate \(\dot G = -P\) with constant
production \(P > 0\), the first-passage time to a fixed threshold
\(G_{\min} \in [0, G_0)\) is
\[T_{\mathrm{phos}} = \frac{G_0 - G_{\min}}{P},\] the reserve-life ratio
being the \(G_{\min} = 0\) special case and a threshold fraction
\(\varepsilon G_0\) giving \((1-\varepsilon)G_0/P\). This is a
conditional reserve-classification proxy under constant production;
because reserves are an economic classification rather than a fixed
physical stock, it is not a forecast of geological exhaustion without an
explicit resource and production model. No stochastic phosphate
extension is required for the interpretation.

The classification addresses the status of the ratio, not the adequacy of the resource base. The
reserve-life convention and its long defence against depletion-pessimistic readings are the subject of
a separate literature (Tilton, 2003; Tilton and Lagos, 2007), whose arguments about substitution,
price-induced discovery and economic recovery are exactly the premises this article registers as
carried rather than discharged; nothing here contradicts them.


\subsubsection{7.7 The explicit
non-claims}\label{the-explicit-non-claims}

The first-passage semantics close with seven explicit non-claims, all of
which hold in this article:

\begin{enumerate}
\def\labelenumi{\arabic{enumi}.}
\tightlist
\item
  The Brownian and geometric-Brownian processes are not stochastic
  completions of the ledger and do not conserve its mass compartments.
\item
  No theorem relates \(\widehat\mu\) to \(-\dot A\) of the reduced
  systems, to the finite-donor primitive system, or to the institutional
  delay equations.
\item
  The model hitting time \(T_A\) of Definition 5 is not shown to be
  inverse Gaussian: it would be inverse Gaussian only if the active-pool
  residual were Brownian with constant drift, which the coupled balance
  (2) does not supply --- the tabled groundwater numbers inherit
  inverse-Gaussian means from Definition 6's surrogate and from nothing
  else.
\item
  The historical groundwater minimum is not an independently identified
  physical failure barrier.
\item
  A shorter surrogate median or It\^o mean is not evidence of faster
  physical depletion.
\item
  The gross active-pool horizon \(H_A^{\mathrm{gross}}\) of Definition 3
  and its productivity-illusion interpretation --- the misreading of a
  large gross-turnover horizon as evidence of slow net depletion, the
  false implication recorded in Section 6.1 --- are not first-passage
  results treated here.
\item
  The fisheries calculation is not a stage-structured fisheries model,
  and the phosphate calculation is not a geological-reserve model.
\end{enumerate}

\subsubsection{7.8 Parameter and observation
uncertainty}\label{parameter-and-observation-uncertainty}

The inverse-Gaussian results condition on the drift, barrier, and noise
scale. In the groundwater application \(\widehat\mu\) is estimated from
a finite, potentially autocorrelated record and the barrier is selected
from that same record; measurement error, serial dependence, seasonal
forcing, spatial aggregation, trend breaks, and common climatic drivers
are separate uncertainties, and integrating any of them out yields a
predictive mixture rather than a single inverse-Gaussian law. A residual
scale estimated from the same window does not by itself identify process
noise. No calibrated predictive distribution is claimed; the full
uncertainty treatment belongs to an empirical identification study, not
to this article. For industrial-ecology measurement this discipline is
the practical message: first-passage distributions on declared
surrogates are usable as descriptive statistics, but their drift,
barrier, and noise inputs each carry their own identification story, and
a calibrated forecast requires that story to be discharged.

\begin{center}\rule{0.5\linewidth}{0.5pt}\end{center}

\subsection{8. Domain Templates at Registered
Status}\label{domain-templates-at-registered-status}

\subsubsection{8.1 The phosphorus
template}\label{the-phosphorus-template}

The phosphorus domain enters at registered template status: an identification ladder for the
resource-product-waste-detritus structure of Section 2.3, whose constitutive content --- yield and loss
functions, recovery fractions, the price response of the reserve classification --- is declared and not
established, and whose reserve and production quantities carry their source vintage (U.S. Geological Survey,
2026). The ladder's competing-model structure and its falsification protocols are recorded obligations rather
than results. The ladders themselves are the supplementary's S2, next to the groundwater ladder this one pairs
with, and the registered detail these paragraphs used to carry --- the competing-model structure and the
falsification protocols in full --- is its S14; nothing in Sections 6 to 10 uses either.

\subsubsection{8.2 The groundwater template and the two-pool
gap}\label{the-groundwater-template-and-the-two-pool-gap}

The groundwater template enters at registered status with one admitted object and one declared gap: the
admitted object is the one-pool affine approximation behind the anomaly-persistence index of Section 6.5.1, and
the two-pool model --- active storage with a slow donor pool, the two-compartment structure of Section 2.2 ---
is not established. The registered identification requirements for closing that gap, and the discipline that
leakage terms may not absorb unexplained residuals, are the subject of the supplementary's S2.1, and none of
them is met by the record used in Section 6.5; the registration as this article previously stated it is the
supplementary's S14.

\subsubsection{8.3 Extractor-side harvest
economics}\label{extractor-side-harvest-economics}

On the extractor side, the same discipline applies to economic steady
states. In the open-access equilibrium of a single-species fishery
(Clark, 1990), the bioeconomic equilibrium stock
\(S_{\mathrm{OA}} = c/(pq)\) is set by cost, price, and catchability ---
and is infeasible as a management target under a conservation floor
\(S_{\min} > S_{\mathrm{OA}}\): the unregulated equilibrium lies below
the floor, and no open-access trajectory is viable against it. The
modified golden rule in its constant-unit-cost form,
\(g'(S_\delta) = \delta\), sets the optimal steady stock for the
discount rate \(\delta\) (Clark's general form carries an additional
marginal-stock-effect term); a harvest tax shifts the open-access
equilibrium to \(S_{\mathrm{OA}} = c/((p - \tau)q)\) --- the tax moves
the economic equilibrium, but it does not move the physical floor. The
distinction is the extractor-side counterpart of the accounting
discipline of Section 6: instrument parameters and constraint thresholds
are different objects, and no tax schedule substitutes for a constraint
the ledger must satisfy. The growth function \(g\) of this paragraph is
a declared constitutive readout on the stock for this extractor-side
remark only; it is not a primitive of the closed natural block of
Section 2, and nothing in this section is promoted into the typed
ledger.

\begin{center}\rule{0.5\linewidth}{0.5pt}\end{center}

\subsection{9. The Interface with Institutional Delay
Dynamics}\label{the-interface-with-institutional-delay-dynamics}

The partition between this article and the companion delay-dynamics analysis (Abaee, 2026, doi:10.5281/zenodo.22554217) is fixed by an interface contract. This article owns the closed material accounting: the
primitive ledger equations and full routing, the conservation and
positivity theorems of Section 4, the componentwise deficit and
depletion diagnostics of Sections 5--6, and the closed-donor no-rest and
extraction-integrability limitations. The companion owns the open
frozen-donor retarded systems and their bifurcation results. The
interface is viable, but not because the closed primitive ledger
dynamically reduces to the open working system: the two are different
completions, and the contract records both the exact shared object and
the rejected mapping.

\textbf{The exact shared object.} Under the single-resource
specialization of Section 5.4 (\(\mu = \nu = \rho = 0\) --- the product,
waste, and price parameters of the unreduced ledger: its
macroeconomic-feedback, recycling, and price-response channels, per the
Section 2.2 gloss --- and the mining intensity \(C^A = 0\)) --- with the
local stock equation \(\dot N = R - qEN\), the deficit identity
\[D(t) := qE(t)N(t) - R(N(t), A(t)) = -\dot N(t), \qquad \Lambda(t) := [D(t)]_+ = [-\dot N(t)]_+\]
holds for every trajectory of either the specialized ledger or the
reduced core (Remark 16). The identity is the one object both analyses
may use without substantive duplication; the reduced core's constitutive
replacement \(R(N,A) \to rN(1 - N/K)\) is separately an approximation
and carries its own finite-time scope (Theorem 1 and Remark 2: the
replacement is pointwise on the interior support region and non-uniform
through the depleted-pool boundary). The contract fixes more than the
deficit identity: the shared object includes the nonnegative orthant and
the sign pattern of harvest as an outflow from the living stock --- a
companion model that routes the ``unsustainable portion'' of a flow into
a different compartment changes the incidence and thereby leaves the
interface (Section 2.5).

\textbf{Proposition 43 (The institutional-failure subsystem is exactly closed).} \emph{Under the institutional-failure
specialization of Section 5.4, the ecological-institutional subsystem on \((N, A^{\mathrm{act}}, A^{\mathrm{geo}}, U, Z, E)\) is an exact closed projection of the ledger for every parameter value: no singular limit, no small
parameter and no timescale separation is required, and no macroeconomic variable enters its vector field.}
\emph{Proof.} The six right-hand sides displayed above depend only on the block's own states and the delayed memory; a
subsystem whose vector field involves no excluded variable is invariant, and invariance is the projection. The
converse fails by construction --- the excluded block is driven by the subsystem --- so the projection is
one-way, and Section 5's macroeconomic readouts remain available without being determined. \(\square\)

What the proposition does not claim is the next step, and the distinction is the point of stating it: the
memory-effort pair \((Z, E)\) is the gated three-state core and working four-state core of the companion
delay-dynamics analysis (under review; its eq. (1), the gated three-state core, and its Section 2.4, which
relates that core to the working four-state model), not an object of this article, and the semiconjugacy
condition \(D\pi(\xi) f(\xi) = F(\pi(\xi))\) on the history phase space --- which would carry closed-block
results across to that system --- is made under the citation and is not re-proved here.

\textbf{The non-reduction boundary.} There is no exact dynamic reduction
from the closed primitive finite-donor ledger to the open working system
--- not as a projectable reduction and not as a regular perturbation.
The reasons are mathematical:

\begin{enumerate}
\def\labelenumi{\arabic{enumi}.}
\tightlist
\item
  The primitive ledger uses the intrinsic donor-limited target
  \(A^{\mathrm{eq,intrinsic}}\); the working system uses the derived
  target
  \(A^{\mathrm{eq,W}} = A^{\mathrm{eq,intrinsic}} + \kappa_A K/\omega_A\).
  The three registered numbers display the separation:
  \(A^{\mathrm{eq,intrinsic}} = 50\), the working active pool
  \(A^{\mathrm{act,*}} = 397.87\), and
  \(A^{\mathrm{eq,W}} = 50 + \kappa_A K/\omega_A = 5050\) --- the two
  equilibria differ by a factor of eight and the two targets by two
  orders of magnitude.
\item
  At the working equilibrium the two \(A^{\mathrm{act}}\) vector fields,
  written at the same state \((N, A^{\mathrm{act}}, U)\), differ by
  \(\omega_A\bigl(A^{\mathrm{eq,W}} - A^{\mathrm{eq,intrinsic}}\sigma\bigr) - \gamma_U U = \kappa_A K - \gamma_U U\)
  under the registered scale separation (\(\sigma \approx 1\)):
  approximately \(0.535\) stock units per year at the working point's
  quasi-rest detritus level (where \(\gamma_U U = T^* \approx 4.47\),
  the gross uptake at the working point) and \(\kappa_A K = 5.000\)
  stock units per year at \(U = 0\) --- an \(O(1)\) to \(O(\kappa_A K)\)
  discrepancy, not a small residual. The two same-state flux readings
  behind it are the working recharge
  \(\omega_A(A^{\mathrm{eq,W}} - A^{\mathrm{act,*}}) = 4.652\) and the
  closed donor flow \(e_{GA} - e_{AG} \approx -0.348\), whose signed
  difference is \(4.652 + 0.348 = 5.000 = \kappa_A K\); the difference
  is \(U\)-dependent because the working field omits the detritus return
  \(\gamma_U U\) that the closed field carries --- the U-handling split
  is part of this obstruction, and \(B^* - R^* = T^*\) is the working
  system's turnover balance, not the field difference.
\item
  The working point requires continuing geological support --- the flux
  \(\omega_A(A^{\mathrm{eq,W}} - A^{\mathrm{act,*}}) = 4.652133\ldots\)
  stock units per year, supplied every year by a donor the working
  system treats as a parameter --- and is not a rest point of the closed
  finite-donor system (Theorem 12). At the same state the closed
  primitive donor flow is
  \(e_{GA} - e_{AG} = \omega_A(A^{\mathrm{eq,intrinsic}} - A^{\mathrm{act,*}}) \approx -0.348\):
  the donor gains in the closed ledger where the working completion has
  it losing \(4.652\) --- the two fields have opposite signs on the
  donor coordinate, not merely different magnitudes. The working-point
  figures \(E^* \approx 2.090\), \(N^* = 89.526\),
  \(A^{\mathrm{act,*}} = 397.87\), and the recharge \(4.652\) are
  imported at the companion's registered precision; the reverse check
  \(qE^*N^* = 0.001 \times 2.090 \times 89.526 \approx 0.187\) is
  consistent to the quoted digits.
\item
  The cumulative donor-draw quantity
  \(\varepsilon_G(T) = G_0^{-1}\int_0^T |e_{GA} - e_{AG}|\, dt\) is a
  diagnostic of the derived-target completion, not a trajectory-tracking
  error between the two fields; no finite-time tracking theorem between
  the completions holds.
\item
  The closed primitive system makes sustained extraction integrable
  (Theorem 14) and therefore cannot possess the working positive-flux
  rest indefinitely.
\end{enumerate}

The five reasons form a trichotomy: (1)--(3) are short-time obstructions
--- the two \(A^{\mathrm{act}}\) fields differ by
\(\kappa_A K - \gamma_U U\) at the same state, which is \(O(1)\) at the
working point's quasi-rest detritus level (\(\approx 0.535\)) and at
most \(\kappa_A K = O(5)\) (at \(U = 0\)), and trajectories of the two
systems diverge on \(O(1)\) time scales; (5) is the long-time
obstruction --- extraction on the closed ledger is \(L^1\) in time
(Theorem 14); (4) is neither --- \(\varepsilon_G\) is not a tracking
error between the two fields at any time scale. Collectively:

\textbf{Theorem (Non-reduction of the open working completion).}
\emph{There is no exact dynamic reduction, no regular perturbation, and
no finite-time tracking correspondence from the closed primitive ledger
(2) to the open working system, because (i) the targets differ by
\(\kappa_A K/\omega_A = 5{,}000\) stock units (structural); (ii) the
\(A^{\mathrm{act}}\) fields differ by \(O(1)\) at the working point
(short-time); (iii) the working point is not a rest point of (2)
(Theorem 12; equilibrium); (iv) \(\varepsilon_G\) is not a tracking
metric (diagnostic misuse); and (v) extraction on (2) is \(L^1\) in time
(Theorem 14; long-time).}

The mapping type for exact dynamic reduction is rejected; the permitted
relation is analogy for shared mechanism language plus diagnostic
reconstruction of omitted mass flows. The companion's global periodic
results are properties of its reduced systems and do not transfer to the
closed primitive ledger; in particular, Hopf or periodic orbits of the
frozen-donor working system are not properties of (2). In the other
direction, the working system is an open projection: omitted turnover is
routed to a diagnostic detritus or inert sink, imposed recharge
corresponds to geological draw, and the reduced trajectory's mass
discrepancy is reconstructible from the omitted flows.

\textbf{The frozen-donor limit is a corollary of the structural clause
(i).} Rescaling the donor as \(G = G_0 g\) with \(g(0) = 1\) gives
\(\dot g = -G_0^{-1}(e_{GA} - e_{AG})\). The limit \(G_0 \to \infty\)
freezes \(g\) but does not restore the working completion's derived
target: the limiting recharge field still uses
\(A^{\mathrm{eq,intrinsic}}\), not \(A^{\mathrm{eq,W}}\), so the scaling
is not a regular perturbation of the working vector field. Local Hopf
persistence of the working system under this primitive scaling is not
claimed; a different derived-target completion would be required before
a regular-perturbation theorem could be formulated.

\textbf{The long-time finite-budget interpretation.} With the donor
\(G(t)\) included as a state, the closed system is an autonomous
retarded equation with a slow donor coordinate. The companion's
\(\tau_+ \approx 150\) yr upper cycle is a frozen-donor object; on the
closed system it can persist only as a transient on the finite donor
budget. The transient-duration statement is an order/budget bound, not
an asymptotic estimate: under a sustained lower extraction flux
\(c > 0\) the duration is bounded above by \(G_0/c\). The scale must
name its flux: at the closed-block extraction rate
\(c = qE^*N^* \approx 0.187\) stock units per year the budget bound is
\(G_0/c \approx 2 \times 10^6\) years, while at the working completion's
recharge flux \(B^* \approx 4.652\) stock units per year the draw scale
is \(G_0/B^* \approx 8.6 \times 10^4\) years --- the ``tens of thousands
of years'' heuristic uses the working flux (at
\(G_0/A^{\mathrm{act}*} = 10^3\)), and both scales sit far above the
institutional delays of the companion family. Whether the frozen-donor
local Hopf structure persists as a slowly drifting transient in the
closed donor system is an open slow-passage problem; the mass budget
alone does not establish it. The implication for industrial-ecology
coupling is concrete: a closed physical ledger and an open working
system can share one diagnostic identity without sharing a dynamics, and
the temptation to import one system's theorems into the other must be
resisted.

\begin{center}\rule{0.5\linewidth}{0.5pt}\end{center}

The companion analysis supplies the one numerical fact about governance timing that this article may
use without re-deriving it: at its calibrated point, annual review under the mobilising law is
unstable, and stability returns above a review interval of about 6.5 years, with the two subcritical
Hopf crossings of the continuous-delay problem certified near 3.7 years and 150 years. The readout of
Section 6.4 inherits the consequence. A margin computed at an interval longer than the admissible
one is not conservative, because the declared object it is computed on has already changed between
reviews; the admissible interval is a property of the institutional block, is reported there, and is
not estimated here.


\subsection{10. What the Ledger Does Not
Support}\label{what-the-ledger-does-not-support}

\subsubsection{10.1 Compensatory aggregation is rejected, not merely
discouraged}\label{compensatory-aggregation-is-rejected-not-merely-discouraged}

On the unreduced ledger, no scalar weighting of components is authorized
by the accounting itself. The precise statement is on the feasible
balance domain (Definition 1): a weighted sum \(\sum_m w_m s_m\)
certifies componentwise adequacy \(s_m \ge d_m\) for all \(m\) on
\(\mathcal{B}(x,t)\) only through an implication proved from the
physical restrictions defining the domain --- and, in general, no such
implication holds. The feasible domain can admit vectors with a positive
aggregate and a negative component, exactly the geometry documented in
the composite-indicator literature (Munda and Nardo, 2009). In short, a
positive weighted sum never certifies componentwise nonnegativity;
scalar summaries may rank and communicate, but certification requires
the vector.

\textbf{Proposition (No weighted certification).} \emph{Let
\(\mathcal{B}(x,t)\) be the feasible balance domain (Definition 1). If
\(\mathcal{B}(x,t)\) contains a vector with \(b_i < 0\) and
\(w^\top b > 0\) for a fixed \(w \ge 0\), then the certificate
\(\{w^\top b \ge 0\}\) does not imply \(b \ge 0\); if
\(\mathcal{B}(x,t)\) is not known to exclude that pattern, no
nonnegative weighting is authorized as a componentwise certificate.}

\emph{Proof.} On the exhibited vector the componentwise predicate fails
(\(b_i < 0\)) while the weighted predicate holds (\(w^\top b > 0\)); the
two predicates differ on \(\mathcal{B}(x,t)\). The inverse-horizon score
of Section 6.5.2 (Non-example 1) is the worked witness: a positive
aggregate coexisting with componentwise deficits by construction. \ensuremath{\square}

The conditional form is not an accident of a particular domain: the
compensating pattern is constructible against \emph{any} weight, so the
failure is universal to the method of weighted certification.

\textbf{Theorem (Universal failure of weighted certification).}
\emph{Let the ledger have \(m \ge 2\) components. For every weight
vector \(w \in \mathbb{R}_+^m\), \(w \neq 0\), there exist a typed
ledger as in Section 2, a demand vector \(d\), and an admissible
state--operation pair \((x, u)\) with \(u \in \mathcal{U}(x,t)\) whose
attainable balance \(b = \mathcal{O}(x,u,\theta) - d\) satisfies
\(w^\top b \ge 0\) while \(b_m < 0\) for some component \(m\). No
nonnegative weighting is therefore a valid componentwise certificate in
general.}

\emph{Proof.} Take the two-compartment ledger with stock
\(x = (x_1, x_2) \in \mathbb{R}^2_+\), one donor-limited transfer flux
\(f = kx_1\) from compartment 1 to compartment 2 (\(k > 0\)), identity
readout \(\mathcal{O}(x, u, \theta) = x\), and demands
\(d = (d_1, d_2)\). Every nonnegative state is admissible with the
declared flux, because donor limitation holds (\(f = 0\) at
\(x_1 = 0\)), so the balance \(b = x - d\) is attainable for every
\(x \in \mathbb{R}^2_+\); the admissible operating set is the flux
declaration itself. Fix \(w\) and let \(j\) be an index with
\(w_j > 0\). If some \(m \neq j\) has \(w_m = 0\), place the deficit
there: choose \(x_m < d_m\) and \(x_j \ge d_j\), giving
\(w^\top b = w_j (x_j - d_j) \ge 0\) while \(b_m < 0\). Otherwise
\(w_m > 0\) for every \(m \neq j\): choose any \(m \neq j\) with
\(x_m < d_m\), and choose
\(x_j \ge d_j + \bigl( w_m (d_m - x_m) + \varepsilon \bigr)/w_j\) for
any \(\varepsilon > 0\). Then
\[w^\top b = w_m (x_m - d_m) + w_j (x_j - d_j) \ge w_m (x_m - d_m) + w_m (d_m - x_m) + \varepsilon = \varepsilon \ge 0,\]
while \(b_m = x_m - d_m < 0\). The pair \((x, d)\) is the witness for
\(w\). \ensuremath{\square}

The construction never uses the dynamics: the failure is a property of
nonnegative weightings over mixed-sign balances, not of the
donor-limited positivity mechanism. Consequently, on any feasible
balance domain containing a compensating pair --- and Definition 1's
domain is exactly that whenever the operating set admits states with
mixed balances --- no nonnegative weighting can stand in for the
conjunctive criterion of Section 3.2. The reading is the algebraic form
of the weak-comparability thesis stated in Section 1.1.

The companion assessment analysis (Abaee, 2026, doi:10.5281/zenodo.22545740) proves
the dynamic form of the same separation for transition operators; the
ledger side contributes the static prerequisite: the aggregation
question is only well posed after the balance domain is declared, and
the burden of proof sits on the aggregation, not on the componentwise
report. The same separation holds at the service layer: a weighted sum
on \(\mathcal{B}(x,t)\) cannot see the directional support gap
\((1 - \alpha_{\mathrm{reg}})\bar s\) of Definition 2 --- an aggregate
that mixes regenerative and non-regenerative provenance never reports
which component carries the gap. In the terms of the
ecological-economics literature of Section 1.1, componentwise adequacy
on the typed ledger is strong sustainability as a conjunctive predicate;
a positive weighted sum is weak sustainability as a ranking device. The
ledger authorizes the first and does not authorize the second.

\textbf{Proposition 29 (The certifying aggregator is unique).} Let \(r\in\mathbb R^{n}\) be component
adequacies expressed in common units, and let \(\mathcal A:\mathbb R^{n}\to\mathbb R\) be monotone,
\(r\le r'\Rightarrow\mathcal A(r)\le\mathcal A(r')\), calibrated, \(\mathcal A(c\mathbf 1)=c\) for
every \(c\), and certifying, \(\mathcal A(r)\ge c\Rightarrow r\ge c\mathbf 1\) for every \(c\). Then
\(\mathcal A(r)=\min_ir_i\).

\emph{Proof.} Certification at \(c=\mathcal A(r)\) gives \(\min_ir_i\ge\mathcal A(r)\). Since
\(r\ge(\min_ir_i)\mathbf 1\), monotonicity and calibration give \(\mathcal A(r)\ge\min_ir_i\).
\ensuremath{\square}

For any admissible aggregator the compensation premium
\[ \Pi(r)=\mathcal A(r)-\min_ir_i\ \ge\ 0 \]
measures the part of a published aggregate that was produced by cross-component trades. Calibration
cannot be dropped: the additively separable functional \(\sum_i\min(0,r_i)\) is continuous, monotone,
faithful in the sense that a nonnegative value admits no deficit, and complete in the sense that no
deficit forces a negative value; it is not the minimum and it is not calibrated. It reports the depth
of the deficits where the minimum reports their existence, and either may be published provided
neither is called the other. Nor can monotonicity be strengthened to strict monotonicity while
keeping both directions: if a continuous \(\mathcal A\) satisfied \(\mathcal A(r)\ge0\Leftrightarrow r\ge0\), then \(\mathcal A\) would vanish on the faces of the nonnegative orthant, since it is
negative wherever some component is and nonnegative on the orthant, and continuity forces equality
at the boundary; a strictly increasing functional then cannot take one value at \((0,1)\) and another
at \((0,2)\). Certification and strict monotonicity are incompatible, which is why the minimum is not
an arbitrary convention.

\textbf{Proposition 30 (Worst concealed deficit).} Let the component balances lie in declared bounds
\(\ell\le b\le u\), and let \(z\) be a published aggregate value \(z=w^{\top}b\) with \(w\ge0\),
\(w\ne0\). The most severe deficit an aggregate of that value can conceal in component \(j\) is
\[ \delta^{*}_j(z)=\max\ \{-b_j:\ w^{\top}b=z,\ \ell\le b\le u\}, \]
a linear programme, finite whenever the bounds are, and positive whenever any admissible \(b\) has
\(b_j<0\). Together with the premium it turns the noncompensation result into a two-sided audit: the
premium measures what the aggregate bought with trades, and \(\delta^{*}_j\) measures what it can
hide. Reported on declared bounds only, it is arithmetic on a published construction; where the
bounds are not declared, the value is reported as not established.

\textbf{Proposition 31 (Aggregates do not transport event times).} Let two stocks and their sinks obey
\(\dot x_i=-kx_i\), \(\dot w_i=kx_i\) with \(k>0\), so that each pair conserves material and remains
nonnegative, and let the aggregate \(Z=x_1+x_2\) obey the exact closed dynamics \(\dot Z=-kZ\). The
initial states \((x_1,x_2)=(2,98)\) and \((50,50)\) generate the identical aggregate trajectory
\(Z(t)=100e^{-kt}\) and, at the common lower barrier \(x_i\ge1\), first hitting times \((\log2)/k\)
and \((\log50)/k\). No function of the aggregate trajectory determines the component event time.
Dynamical closure and barrier observability are therefore separate obligations: an exact reduced
model is not thereby competent about the events its components are declared to suffer.

\emph{Proof.} Direct evaluation of each exponential at its own barrier; the two aggregate trajectories
coincide identically. \ensuremath{\square}

The asymmetry is exact and one-directional. A coarse aggregate refutes, because a violated aggregate
barrier is a violated component barrier, and it alarms, because a positive premium is a measured
quantity; it does not certify, because the concealed deficit of Proposition 30 can be made positive
at any published aggregate value the domain admits.


\subsubsection{10.2 The double-counting
discipline}\label{the-double-counting-discipline}

Five rules, each carried by a proved or defined statement of this
article, jointly prevent double counting and phantom mass:

\begin{enumerate}
\def\labelenumi{\arabic{enumi}.}
\tightlist
\item
  \textbf{One balance per moiety.} Conservation laws attach to declared
  moieties (Theorems 7--8); adding unlike units --- biomass, money,
  biodiversity indices, exergy --- into one conserved scalar is not
  authorized by any conservation theorem.
\item
  \textbf{Explicit stoichiometry.} Entries are added within an incidence
  row only when their types and units agree; every conversion is an
  explicit coefficient in \(S_{\mathcal{T}}\), never an implicit sum
  (Section 2.1).
\item
  \textbf{Yield routing.} A transformation represented with yield below
  one must route the omitted fraction to a represented compartment or a
  declared boundary flow; otherwise the moiety balance holds only after
  silently dropping the moiety (Theorem 15).
\item
  \textbf{No ghost sinks.} Every primitive with an outflow from some
  compartment must have its routed inflow represented, and every inflow
  its source --- opposite-sign incidence of the same primitive is
  checkable column by column --- and the six-state cancellation of
  Section 4.9 shows the check passing, while the same template shows
  that cancellation without donor-limited admissibility establishes
  nothing.
\item
  \textbf{Classification labels stay out of the columns.} Reserve-life
  and resource-threshold quantities answer different questions and share
  a column only under an explicit convention label (Section 6.5.3);
  diagnostic labels never determine material routing (Section 2.5).
\end{enumerate}

Every unit of mass is routed once: the unit-sum split routing of Theorem
8 and the column-sum-zero incidence of Theorem 9 are the mechanism, and
the yield-routing obligation of Theorem 15 is its enforcement. A flux
omitted from the ledger is not thereby conserved; a recovery claimed in
a quality-neutral loop is not thereby real; and a diagnostic index that
mixes basins and reserves (the ranking device of Section 6.5.2) is
identified as such in the text that reports it. Double counting is a
representation error, and the representation that prevents it is the
contribution.

\textbf{Remark 33 (One-signed bias of an aggregate overshoot date).} Let components carry biocapacity
\(b_i\) and demand \(d_i\), and form an aggregate overshoot date from the aggregates:
\[ \tau_{\mathrm{agg}}=365\,\frac{\sum_ib_i}{\sum_id_i}=365\sum_iw_ir_i, \qquad w_i=\frac{d_i}{\sum_jd_j}, \qquad r_i=\frac{b_i}{d_i}. \]
Then \(\tau_{\mathrm{agg}}\) is a demand-weighted mean of the component ratios, so
\(\tau_{\mathrm{agg}}\ge\tau_{\min}:=365\min_ir_i\), with equality exactly where every component of
positive weight coincides. The difference
\[ \Pi_\tau=\tau_{\mathrm{agg}}-\tau_{\min}=365\sum_iw_i\bigl(r_i-r_{\min}\bigr) \]
is the compensation premium of Section 10.1 in unit disguise: one-signed, exactly decomposable,
unbounded in the dispersion of the components, and computable from published component tables. An
aggregate overshoot date is therefore optimistic in a known direction, which is a statement about
the construction rather than about any territory, and it is the one applied claim here that requires
no recomputation of a public dataset to state. A figure can be read straight off the published component table: on the world totals of the 2025 edition of the National Footprint and Biocapacity Accounts --- the edition whose series contains that year, since the release Lin et al. (2018) document runs from 1961 to 2014 --- the 2022 aggregate ratio is \(0.584\) of biocapacity to demand, so \(\tau_{\mathrm{agg}} = 213\) d, and because the carbon component carries zero
biocapacity the minimum component ratio is \(0\) --- hence \(\Pi_\tau = 213\) d, the whole aggregate date.
Restricted to the five components of positive biocapacity with renormalised weights, the same arithmetic
gives \(\tau_{\mathrm{agg}} = 538\) d against \(\tau_{\min} = 365\) d, a premium of \(173\) d, and the
restricted premium runs \(547\) d (1961), \(346\) d (1980), \(251\) d (2000), \(173\) d (2022): it shrinks as the
components converge, exactly as the display says it must. Which rows enter the ratio set is fixed by the accounts, not by this
construction: the Guidebook to the National Footprint and Biocapacity Accounts states that no biocapacity
figure is computed for carbon uptake, because carbon demand is charged against forest land biocapacity and a
carbon biocapacity of its own would double count it (Global Footprint Network, 2021, Section 9.1.2). The zero
row is therefore a published convention, and it is what makes \(\tau_{\min} = 0\) identically, so the premium
above equals the whole aggregate date by construction rather than by dispersion. Excluding the carbon
\emph{demand} instead is a different object --- the non-carbon components only --- on which
\(\tau_{\mathrm{agg}} > 365\) d states that no included component overshoots within the year, not that no
overshoot occurs. Both figures are world totals of one data vintage: the accounts are recomputed for every
year of the series at each release, so the 2022 row of the 2025 edition (213 d) need not reproduce the date
announced for 2022 (28 July, the 209th day of that year), and that four-day difference is vintage, nowcasting and aggregation level rather than arithmetic. On the newest edition of the accounts that can be fetched without an account --- the 2018 edition, whose series ends in 2014 --- reading the publisher's world aggregate instead of summing the national rows moves that date by up to 9 d over the last decade of the release, while a whole edition step on a closed year moves it by at most 1.9 d. A gap of a few days is therefore evidence about neither the arithmetic nor the vintage alone: which level the accounts are aggregated at is a choice the reader makes, and on this release it is worth as much as an edition. The figures are world totals, not any territory's.


\textbf{Remark 37 (Two component ratios are identities at a world aggregate).} \emph{On the world aggregate of the published accounts the restricted minimum \(\tau_{\min}\) is 365 d in every year of the series, because the two components whose demand is the area itself --- cropland and built-up land --- have \(b_i = d_i\) identically; the restricted premium is then the distance of the aggregate date past the year boundary, and its decline measures the dispersion of the three yield-based components around a level the accounts fix, not a convergence the ecosystem exhibits.} \emph{Verification.} On the 2018 edition, \(r_{\mathrm{crop}} = r_{\mathrm{built}} = 1\) exactly in 54 of 54 years and \(\tau_{\min}^{\mathrm{NOC}} = 365\) d in 54 of 54; over the last decade of that series the ratio \(\tau_{\mathrm{agg}}^{\mathrm{NOC}}/\tau_{\mathrm{agg}}^{\mathrm{ALL}}\) equals \(1/(1-s_{\mathrm{carbon}})\) to within 4.4e-16, where \(s_{\mathrm{carbon}}\) is carbon's share of demand and runs between 58.6% and 61.6%. The computation is one weighted sum of the published component rows, and the record is the supplementary's S17. \emph{Consequence.} A premium computed at a world aggregate and a premium computed on a country's own component table are not two readings of one ecological fact, because only at the aggregate level are two of the ratios pinned by the convention that prices world demand at world-average yields; the level at which the accounts were read belongs to the number the way a data vintage does, and Section 10.2's sensitivity sentence says so.

\subsubsection{10.3 Negative and boundary content is
first-class}\label{negative-and-boundary-content-is-first-class}

The classification results of Section 6.5 are negative results stated as
such: the anomaly index is not a stock ratio; the reserve-life ratio is
not a forecast; the removals-only time is not a depletion diagnostic.
The non-reduction boundary of Section 9 is a rejected mapping with five
mathematical reasons. The sink obstructions of Section 2.4 are
empty-kernel mechanisms. None of these is a failure of the framework: a
quantity that answers exactly one question, stated with that question,
is the framework working --- and the framework's claim about itself is
limited to the accounting layer it establishes.

\subsubsection{10.4 Limitations}\label{limitations}

\begin{enumerate}
\def\labelenumi{(\roman{enumi})}
\tightlist
\item
  The two-pool exact specialization of the groundwater template remains
  open; the admitted object is the one-pool affine approximation, and no
  two-pool model is claimed as established (Section 8.2). (ii) The
  phosphorus and groundwater rows are registered template obligations:
  no constitutive content exists behind their identification ladders.
  (iii) The first-passage propositions of Section 7 concern declared
  stochastic surrogates, not the ledger: they do not compute the
  ledger's hitting time, do not conserve its mass, and carry the
  record-relative-barrier and non-claim disciplines. (iv) The applied
  records of Section 6.5 are classified diagnostics at their stated
  evidentiary levels --- statistical index, arithmetic ratio,
  removals-only pressure scale --- and none is a calibrated
  early-warning system or a forecast. (v) The non-reduction boundary of
  Section 9 is permanent mathematics, not a gap: the closed primitive
  ledger and the open working system are different completions. (vi) The
  conditional hybrid balance of Theorem 15 stays conditional, with its
  jump-interpretation and yield-routing obligations open per
  application. (vii) The support-saturation limits of Section 2.6 are
  local and finite-time; neither is a full-system reduction. (viii) The
  article asserts nothing empirical about any named resource system
  beyond the classifications of public data products stated at their
  source status.
\end{enumerate}

\begin{center}\rule{0.5\linewidth}{0.5pt}\end{center}

\subsection{11. Conclusion}\label{conclusion}

The two failure modes of the introduction --- compensatory aggregation
and classification drift --- are representation errors, and the
representation that prevents them is the article's content. Conservation
is proved from the incidence structure, not assumed; positivity is
proved from donor limitation, not asserted; services are readouts, not
mass; and depletion time is three quantities, not one. The closed
finite-donor ledger carries its complete theorem set --- the
natural-block mass identity, orthant invariance, no interior rest at
positive effort, the vanishing-extraction rest set (extinction, carrying
capacity, and the frozen-biomass face), extraction integrability --- and
the non-reduction boundary records, with reasons, why the closed ledger
and the open working systems of institutional dynamics are different
completions sharing one exact object. The depletion numbers of the
applied record answer, each, exactly the question their construction
poses: a statistical index of record-relative stress, an arithmetic
ratio of an economic classification, a pressure scale of one gross loss
rate. Stating them with those questions --- against the published
critiques they corroborate --- is what makes them usable; collapsing
them into one ``time to depletion'' is what makes them false.

The article likewise locates substitution within the ledger rather than
alongside it. The weak and strong regimes of Section 1.1 are two
readings of one typed stock--flow ledger --- distinguished by whether
the material cycle closes at the rate of use; the reading is developed
once, in the introduction, and is not re-argued here. They are not rival
doctrines: the vector reading is the one that carries the certificate,
the scalar \(b\cdot M\) check of Section 1.1 is the operational reading
on top of it, and a substitute is either a recycled flux returned to the
regenerating pool or a non-renewable drawdown on a second compartment
--- different ledger entries with different statuses. An
exhaustion-horizon estimate built on a reserve figure carries the
substitution and technology premises of the reserve classification, not
a physical forecast. For ecological-economics measurement, that is the
closing statement: not rival doctrines but two readings of one ledger,
and the vector reading is what carries the certificate.

\begin{center}\rule{0.5\linewidth}{0.5pt}\end{center}

\begin{center}\rule{0.5\linewidth}{0.5pt}\end{center}

\subsection{References}\label{references}

Abaee, A., 2026. Delay-induced regime change in harvested stocks: the
mobilising and protective channels of institutional feedback, and the
review interval as control. Zenodo.
https://doi.org/10.5281/zenodo.22554217. Companion delay-dynamics study.

Abaee, A., 2026. Periodic review as sampled governance: sample-and-hold
dynamics of assessment-driven effort control, a selected 42-stock
spectral screen, and the Northern Cod case. Zenodo.
https://doi.org/10.5281/zenodo.22554297. Companion review-screen study.

Abaee, A., 2026. The limits of compensatory aggregation: a formal
separation of weak and strong sustainability assessment. Zenodo.
https://doi.org/10.5281/zenodo.22545740. Companion assessment-separation
study.

Ahuja, R.K., Magnanti, T.L., Orlin, J.B., 1993. Network Flows: Theory, Algorithms, and Applications. Prentice-Hall, Englewood Cliffs.

Aubin, J.-P., 1991. Viability Theory. Birkh\"auser, Boston.

Baez, J., Li, X., Libkind, S., Osgood, N.D., Patterson, E., 2023. Compositional modeling with stock and flow diagrams. In: Applied Category Theory 2022. Electronic Proceedings in Theoretical Computer Science 380, 77--96. https://doi.org/10.4204/EPTCS.380.5

Blomqvist, L., Brook, B.W., Ellis, E.C., Kareiva, P.M., Nordhaus, T.,
Shellenberger, M., 2013. Does the shoe fit? Real versus imagined
ecological footprints. PLoS Biology 11, e1001700.
https://doi.org/10.1371/journal.pbio.1001700

Brunner, P.H., Rechberger, H., 2004. Practical Handbook of Material Flow
Analysis. Lewis Publishers, Boca Raton.

Chhikara, R.S., Folks, J.L., 1989. The Inverse Gaussian Distribution:
Theory, Methodology, and Applications. Marcel Dekker, New York.

Daly, H.E., 1990. Toward some operational principles of sustainable
development. Ecological Economics 2, 1--6. Clark, C.W., 1990.
Mathematical Bioeconomics: The Optimal Management of Renewable
Resources, 2nd ed.~Wiley, New York.

Ekins, P., Simon, S., Deutsch, L., Folke, C., De Groot, R., 2003. A
framework for the practical application of the concepts of critical
natural capital and strong sustainability. Ecological Economics 44,
165--185.

Eurostat, 2001. Economy-wide Material Flow Accounts and Derived
Indicators: A Methodological Guide. Eurostat, Luxembourg.

Feinberg, M., 2019. Foundations of Chemical Reaction Network Theory.
Springer, Cham.

Fischer-Kowalski, M., Krausmann, F., Giljum, S., Lutter, S., Mayer, A.,
Bringezu, S., Moriguchi, Y., Sch\"utz, H., Schandl, H., Weisz, H., 2011.
Methodology and indicators of economy-wide material flow accounting:
state of the art and reliability across sources. Journal of Industrial
Ecology 15, 855--876.

Gale, D., 1957. A theorem on flows in networks. Pacific Journal of Mathematics 7, 1073--1082.

Global Footprint Network, 2021. Working Guidebook to the National Footprint and Biocapacity Accounts, 2021 edition. Global Footprint Network, Oakland, CA.

G\"untner, A., Sharifi, E., Haas, J., et al., 2024. Global Gravity-based
Groundwater Product (G3P), V. 1.12. GFZ Data Services.
https://doi.org/10.5880/G3P.2024.001

Illakwahhi, D.T., Vegi, M.R., Srivastava, B.B.L., 2024. Phosphorus'
future insecurity, the horror of depletion, and sustainability measures.
International Journal of Environmental Science and Technology 21,
9265--9280. https://doi.org/10.1007/s13762-024-05664-y

Jacquez, J.A., Simon, C.P., 1993. Qualitative theory of compartmental
systems. SIAM Review 35, 43--79.

Lin, D., Hanscom, L., Murthy, A., Galli, A., Evans, M., Neill, E.,
Mancini, M.S., Martindill, J., Medouar, F.-Z., Huang, S., Wackernagel,
M., 2018. Ecological footprint accounting for countries: Updates and
results of the National Footprint Accounts, 2012--2018. Resources 7, 58.
https://doi.org/10.3390/resources7030058

Martinez-Alier, J., Munda, G., O'Neill, J., 1998. Weak comparability of
values as a foundation for ecological economics. Ecological Economics
26, 277--286.

Meadows, D.H., Meadows, D.L., Randers, J., Behrens III, W.W., 1972. The
Limits to Growth. Universe Books, New York.

Munda, G., Nardo, M., 2009. Noncompensatory/nonlinear composite
indicators for ranking countries: a defensible setting. Applied
Economics 41, 1513--1523.

Neumayer, E., 2013. Weak versus Strong Sustainability: Exploring the
Limits of Two Opposing Paradigms, 4th ed.~Edward Elgar, Cheltenham.

\O{}ksendal, B., 2003. Stochastic Differential Equations: An Introduction
with Applications, 6th ed.~Springer, Berlin.

Prajna, S., Jadbabaie, A., 2004. Safety verification of hybrid systems using barrier certificates. In: Hybrid Systems: Computation and Control VII. Lecture Notes in Computer Science 2993, 477--492.

Prajna, S., Jadbabaie, A., Pappas, G.J., 2007. A framework for worst-case and stochastic safety verification using barrier certificates. IEEE Transactions on Automatic Control 52, 1415--1428.

Redner, S., 2001. A Guide to First-Passage Processes. Cambridge
University Press, Cambridge.

Ricard, D., Minto, C., Jensen, O.P., Baum, J.K., 2012. Examining the
knowledge base and status of commercially exploited marine species with
the RAM Legacy Stock Assessment Database. Fish and Fisheries 13,
380--398.

Smith, H.L., 1995. Monotone Dynamical Systems: An Introduction to the Theory of Competitive and Cooperative Systems. Mathematical Surveys and Monographs 41. American Mathematical Society, Providence.

Tapley, B.D., Bettadpur, S., Ries, J.C., Thompson, P.F., Watkins, M.M.,
2004. GRACE measurements of mass variability in the Earth system.
Science 305, 503--505.

Tilton, J.E., 2003. On Borrowed Time? Assessing the Threat of Mineral
Depletion. Resources for the Future, Washington, DC.

Tilton, J.E., Lagos, G., 2007. Assessing the long-run availability of
copper. Resources Policy 32, 19--23.

U.S. Geological Survey, 2026. Mineral Commodity Summaries 2026:
Phosphate Rock. USGS, Reston, VA.
https://www.usgs.gov/centers/national-minerals-information-center/mineral-commodity-summaries

United Nations, 2025. System of National Accounts 2025. United Nations Statistics Division, New York. Endorsed by the United Nations Statistical Commission at its fifty-sixth session, March 2025. https://unstats.un.org/unsd/nationalaccount/sna2025.asp

United Nations, European Commission, International Monetary Fund, Organisation for Economic Co-operation and Development, World Bank, 2014. SEEA Central Framework: 2012 Technical Implementation. Statistical Papers, Series M No. 96. United Nations, New York.

Wackernagel, M., Beyers, B., 2019. Ecological Footprint: Managing our
Biocapacity Budget. New Society Publishers, Gabriola Island, BC.

\begin{center}\rule{0.5\linewidth}{0.5pt}\end{center}

\subsection{Supplementary material}\label{supplementary-material}

The accompanying file \texttt{paper3\_supplementary\_v11.md} carries: the
ten-state admissibility template and its three audited negative
witnesses (the non-donor-limited geological exchange, the
variance-closure failure, the undefined output functional); the
registered identification ladders of the phosphorus and groundwater
templates at full detail; the split-assignment mechanism table; the
statement inventory with the status of every statement in the main text
(theorem with displayed proof, conditional theorem, definition,
application record, or boundary statement) --- read with the S6
statement-status naming offset, which maps the supplementary's status
words to the main text's current labels; and the fisheries cohort record
with the archived-pull verification and the executed broad-cohort
comparison (S5). At this revision it additionally carries the proof obligations attached to each entry of the certification state of Section 3.1 with the certificate vectors of the classified indicators (S7), the linear programmes of Definitions 21--22 and Theorem 24 with their input requirements and the reading rule for an infeasible programme (S8), and the worked exhibits of Sections 6.2, 6.5 and 10.1 with their reproduction record and the statement
inventory extended to the new labels (S9). What v39 folded out of the article is here rather than gone: the
G3P basin-row provenance (S5.4), the identifiability and value records (S10, S11) and the registered
template detail that left Sections 8.1 and 8.2 (S14 to S16). And S17, added with this revision, is the
overshoot-date recomputation behind Remark 37, with the two edition tables it read named and hashed.

\section*{Declarations}

\subsection*{Data availability}

All computations underlying Section 6.5 are descriptive arithmetic on
the cited public data products: the G3P groundwater anomaly product
v1.12 (G\"untner et al., 2024; GFZ Data Services,
doi:10.5880/G3P.2024.001), the U.S. Geological Survey Mineral Commodity
Summaries 2026 (the January 2026 release), and the RAM Legacy Stock Assessment Database, release v4.66 (Ricard et al., 2012; Zenodo 14043031; the cohort is the
dated extract recorded in the supplementary's S5 record, and no cohort statistic is quoted from any other release). The overshoot-date arithmetic of Section 10.2 and Remark 33 is read off the component tables of the National Footprint and Biocapacity Accounts, whose current edition the publisher distributes free but behind registration. The two editions used for the sensitivity check recorded in the supplementary's S17 --- the 2018 edition, with series 1961 to 2014, and the 2017 edition, with series 1961 to 2013, both deposited by the publisher under CC BY-SA 4.0 --- are archived with the reproduction bundle, together with the checksums of the tables as read and the script that computes both conventions from them. The parameter tables of Section 2 are declared parameterizations. No other data were used.

\subsection*{Code availability}
The scripts that generate every computed figure in this article --- the persistence-index simulation of Section 6.5, the curvature and crossover arithmetic of Section 6.2, and the compensation-premium and worst-concealed-deficit linear programmes of Section 10.1 --- are archived with the supplementary material, with the random seed and the library versions recorded in the archive manifest. A fourth script is deliberately not claimed: the overshoot-date recomputation that supports Remark 37 and the supplementary's S17 is archived as a separate analysis record beside the reproduction bundle, because the bundle's manifest scopes it to arithmetic on declared figures and forbids it from ingesting data. That record reads only the two edition tables named above, needs nothing beyond the standard library, prints the aggregation diagnostic Section 10.2 reports, and re-runs on a newer edition of the accounts by changing two file names rather than any line of arithmetic.

\subsection*{Funding}

None declared.

\subsection*{Declaration of competing interest}

None.

\subsection*{AI declaration}
GLM (Z.ai), Qwen (Alibaba Cloud) and DeepSeek AI assisted with drafting and iterative review.

\end{document}
